% ============================================================
%  MÓDULO 10 — Catálogo Completo de Agentes e Subagentes
%  ARQUIVO GERADO — não editar à mão. Rode: python3 gen_mod10.py
% ============================================================
\chapter{Catálogo Completo — 248 Fichas de Agente (248 registradas no runtime)}

\roteirodidatico
{Este catálogo é um inventário de perfis descritos no projeto. Para encontrar um agente, comece pela categoria e leia a ficha como uma declaração de escopo e instruções.}
{\emph{Ficha} é documentação; \emph{slug} é o identificador usado na configuração; \emph{runtime} é a sessão do aplicativo; \emph{ativação} é uma chamada efetiva do perfil. Nas permissões, \cod{allow} autoriza uma ação e \cod{deny} a bloqueia, conforme a política carregada.}
{Escolha um perfil, compare seu nome humano ao identificador, leia as instruções e confira permissões e requisitos. Depois localize um registro que demonstre se houve chamada.}
{A reconciliação entre fichas e configuração mede consistência declarativa. Não mede desempenho, adequação da resposta nem taxa de sucesso.}
{O que a ficha permite concluir sobre o perfil e que observação adicional seria necessária para afirmar que ele executou bem uma tarefa?}

\begin{resultado}
O \pth{opencode.json} registra \textbf{248} agentes (1 orquestrador, 247 subagentes), e o acervo \pth{agents/catalog/} traz \textbf{248} fichas. Cada agente tem exatamente uma ficha e cada ficha um agente: a cobertura documental é \textbf{1:1} (0 fichas sem registro, 0 agentes sem ficha). As 248 fichas abaixo são geradas deterministicamente por \pth{gen\_mod10.py} e agrupam-se por categoria.
\end{resultado}

\begin{aviso}
\textbf{Convenção de nomes (verificada em 29/09/2026, SPEC-935-R618).} O \pth{frontmatter} de cada card usa o \emph{nome humano} (ex.: \emph{Cloud SQL PostgreSQL Specialist}), enquanto \pth{opencode.json} registra o \emph{slug} em kebab-case (ex.: \cod{cloud-sql-postgres-specialist}). Uma versão anterior deste capítulo comparava os dois textos literalmente e reportava trinta e nove fichas como não registradas — falso positivo de convenção, não lacuna de runtime: reconciliadas pelo slug e por alias declarado, todas as 248 fichas estão registradas. Ficha documentada ainda $\neq$ agente testado: registrar não prova que o agente funcione.
\end{aviso}

\begin{processo}
GATILHO. o leitor quer saber se o runtime entrega o que o acervo promete, ou se documentacao e disponibilidade sao coisas distintas.
ROTA. reconciliar ficha contra registro antes de contar inventario.
FUNCAO. transformar o catalogo enumerativo em um teste de consistencia executavel, e nao em uma lista de nomes.
\end{processo}

\section[Leitura guiada]{Leitura guiada — Módulo 10: o catálogo como teste de consistência}

\begin{descricao}
\textbf{O fenômeno.} Um capítulo que enumera 248 agentes produz a impressão de catálogo fiel. Mas a enumeração é a forma mais fácil de mascarar um erro de reconciliação: se o procedimento de contagem compara o nome humano do card com o slug do runtime, a lista parece cheia de agentes faltantes que não faltam. Foi exatamente o que aconteceu: o capítulo afirmava não estarem registradas fichas que estavam todas registradas, apenas sob outra convenção de nome.
\end{descricao}

\begin{definicao}
\textbf{Grandezas e definições.}\par
(1) \emph{Ficha}: registro documental em \pth{agents/catalog/}, identificado pelo \emph{nome humano} no frontmatter.\par
(2) \emph{Agente registrado}: entrada na chave \pth{agent} do \pth{opencode.json}, identificada pelo \emph{slug} em kebab-case.\par
(3) \emph{Reconciliação}: correspondência bijetiva entre ficha e agente, resolvida por nome exato, alias declarado ou slug.\par
(4) \emph{Cobertura documental}: a razão entre o número de agentes com ficha e o número de agentes registrados.\par
(5) \emph{Falso positivo de convenção}: contagem que reporta ausência onde há apenas diferença de grafia entre as duas colunas.
\end{definicao}

\begin{metodo}
\textbf{Dedução passo a passo.} Por que o erro numérico é inofensivo e o erro de reconciliação não?\par
(1) Uma lista de nomes não tem invariante: qualquer contagem produz um inteiro, e um inteiro plausível é indistinguível de um inteiro errado.\par
(2) Uma reconciliação tem invariante: o conjunto de fichas e o conjunto de agentes devem ter a mesma cardinalidade, e cada elemento de um deve ter exatamente um parceiro no outro.\par
(3) Sem a convenção de nomes explicitada, a comparação literal entre duas colunas de nomes mede \emph{grafia}, não presença.\par
(4) A soma bruta de dois conjuntos que se correspondem por alias conta cada agente duas vezes: o total fica inflado e a discrepância aparente, quando na verdade inexiste.\par
(5) Logo, a contagem correta exige reconciliar antes de somar, e o resíduo da reconciliação é o número que merece atenção.
\end{metodo}

\begin{resultado}
\textbf{Exemplo resolvido.} Estado da configuração consultada nesta geração, com $N = 248$ agentes no \pth{opencode.json} e $F = 248$ fichas em \pth{agents/catalog/}.\par
(1) Contagem ingênua por igualdade textual: $42$ fichas sem correspondência exata, o que sugeriria cobertura de $\frac{206}{248} \approx 83\%$ — número plausível, e por isso mesmo perigoso.\par
(2) Reconciliação por slug: os $42$ casos se resolvem em $24$ por equivalência direta de slug e em $18$ por alias declarado (ex.: \emph{Synthetic University — Orquestrador Acadêmico Transversal} $\to$ \cod{university_synthetic}).\par
(3) Resíduo: $42 - 24 - 18 = 0$ fichas sem agente; e $0$ agentes sem ficha.\par
(4) Cobertura documental correta: $\frac{248}{248} = 1 = 100\%$.\par
(5) Leitura correta do número verdadeiro: uma versão anterior somava os dois acervos como se fossem disjuntos e obtinha um total maior que o número real de agentes — supercontagem devida a aliases contados em duplicidade.
\end{resultado}

\begin{resultado}
\textbf{Ordem de grandeza e verificação.} A cobertura de $100\%$ é exata e, ainda assim, diz pouco: as duas maiores categorias ($128$ e $52$ fichas, de 11 no total) somam $\frac{180}{248} \approx 72,6\%$, de modo que a média por categoria ($248/11 \approx 22,5$) descreve mal a distribuição — o mesmo mecanismo de concentração que oculta o desequilíbrio no Módulo 8. \emph{Verifique você mesmo:} rode \cod{python3 gen_mod10.py} e confira quantos elementos são emitidos; alterar o \pth{opencode.json} e reexecutar deve mover esse número na mesma medida.
\end{resultado}

\begin{aviso}
\textbf{Limite de validade.} Reconciliação é consistência declarativa, não evidência funcional: um agente registrado e documentado pode ainda estar quebrado, jamais invocado ou sem permissão adequada. O resultado $100\%$ atesta que nome e documento batem, não que o agente funcione: atestar comportamento é o gate de teste, descrito no Módulo 4.
\end{aviso}

\begin{resultado}
\textbf{Exercícios.} (1) Acrescente um agente fictício ao \pth{opencode.json} sem card correspondente e calcule a cobertura resultante. (2) Remova o alias de \emph{Synthetic University} no gerador e observe o teste de reconciliação falhar. (3) Explique, em uma frase, por que \emph{inventário completo} e \emph{inventário correto} são propriedades diferentes.
\end{resultado}

\section{Panorama por categoria}

\begin{resultado}
Contagem por categoria (ordem decrescente de fichas):\par
\smallskip
{\footnotesize
\begin{tabularx}{\textwidth}{lrX}
\toprule
Categoria & Fichas & Exemplo de agentes \\
\midrule
Ferramentas e apoio transversal & 128 & Gestor de decisões arquiteturais, Moderador de fórum e debates multiagente, Pipeline de extração e análise de documentos \\
Produção acadêmica & 52 & 00 editor chefe phd, 01 agente diagnostico escopo, 02 agente busca curadoria \\
Engenharia & 16 & Abnt latex modular, Programador de, Docx abnt converter \\
Pesquisa e evidências & 14 & Auditoria reprodutibilidade, Bibtex crossref auditor, Inferencia causal did iv rdd \\
Auditoria e qualidade & 13 & Auditor, Preparação de livros para ajuste fino, Consulta de livros por MCP \\
Escrita e tradução & 13 & Author voice guardian, Back translation verifier, Cultural episteme agente \\
Apoio jurídico & 4 & Especialista em síntese de documentos jurídicos, Auxjuris email drafter, Auxjuris jurídico assistant \\
Dados e proveniência & 3 & Especialista em Haystack RAG, Mira chart, Mira chart race \\
Orquestração & 3 & Mira deck academico, Mira new, Mira planner \\
Inferência e modelagem & 1 & Colibri agente \\
Testes & 1 & Executor do ambiente de avaliação \\
\bottomrule
\end{tabularx}\par}
\end{resultado}

\clearpage
\section{Categoria: Orquestração (3 perfis)}

\begin{descricao}
Agrupa perfis ligados à coordenação, divisão e acompanhamento de tarefas. No Core, suas contribuições se relacionam à orquestração e ao Blackboard (Módulo 2).
\end{descricao}

\elemento{Mira deck academico (v1.0.0)}{\metainfo{Categoria}{Orquestração} \hfill \metainfo{Tipo}{mira-agent} \hfill \metainfo{Identificador}{mira-deck-academico} \hfill \metainfo{Ficha}{mira-deck-academico.md}}

\begin{descricao}
Gera decks MIRA de defesa acadêmica com animações didáticas e simula a banca
\end{descricao}

\begin{funcao}
No Core, este perfil apoia tarefas relacionadas a gera decks MIRA de defesa acadêmica com animações didáticas e simula a banca. O orquestrador pode encaminhar-lhe o trabalho na etapa correspondente; a resposta retorna ao fluxo principal para integração e conferência de fontes, critérios e permissões. A ficha documenta a função, mas não comprova que o perfil tenha sido chamado ou executado a tarefa.
\end{funcao}

\begin{arquitetura}
\textbf{Modo de execução declarado:} subagente · \textbf{Permissões:} edit: deny; bash: deny
\end{arquitetura}

\elemento{Mira new (v1.0.0)}{\metainfo{Categoria}{Orquestração} \hfill \metainfo{Tipo}{mira-agent} \hfill \metainfo{Identificador}{mira-new} \hfill \metainfo{Ficha}{mira-new.md}}

\begin{descricao}
Agente de entrada do MIRA. Coleta briefing, objetivo, formato, audiência e prepara o deck para as próximas etapas do pipeline.
\end{descricao}

\begin{funcao}
No Core, este perfil apoia tarefas relacionadas a agente de entrada do MIRA.. O orquestrador pode encaminhar-lhe o trabalho na etapa correspondente; a resposta retorna ao fluxo principal para integração e conferência de fontes, critérios e permissões. A ficha documenta a função, mas não comprova que o perfil tenha sido chamado ou executado a tarefa.
\end{funcao}

\begin{arquitetura}
\textbf{Modo de execução declarado:} subagente · \textbf{Permissões:} edit: deny; bash: deny
\end{arquitetura}

\elemento{Mira planner (v1.0.0)}{\metainfo{Categoria}{Orquestração} \hfill \metainfo{Tipo}{mira-agent} \hfill \metainfo{Identificador}{mira-planner} \hfill \metainfo{Ficha}{mira-planner.md}}

\begin{descricao}
Agente planejador do MIRA que transforma o briefing em um plano de slides, narrativa e checkpoints de aprovação.
\end{descricao}

\begin{funcao}
No Core, este perfil apoia tarefas relacionadas a agente planejador do MIRA que transforma o briefing em um plano de slides, narrativa e checkpoints de aprovação.. O orquestrador pode encaminhar-lhe o trabalho na etapa correspondente; a resposta retorna ao fluxo principal para integração e conferência de fontes, critérios e permissões. A ficha documenta a função, mas não comprova que o perfil tenha sido chamado ou executado a tarefa.
\end{funcao}

\begin{arquitetura}
\textbf{Modo de execução declarado:} subagente · \textbf{Permissões:} edit: deny; bash: deny
\end{arquitetura}

\clearpage
\section{Categoria: Produção acadêmica (52 perfis)}

\begin{descricao}
Reúne etapas especializadas da escrita acadêmica. Os perfis podem apoiar partes do fluxo MASWOS, da delimitação do problema à revisão editorial; citações e resultados continuam sujeitos a conferência (Módulo 6).
\end{descricao}

\elemento{00 editor chefe phd (v1.0.0)}{\metainfo{Categoria}{Produção acadêmica} \hfill \metainfo{Tipo}{maswos-agent} \hfill \metainfo{Identificador}{00\_editor\_chefe\_phd} \hfill \metainfo{Ficha}{00\_editor\_chefe\_phd.md}}

\begin{descricao}
Editor-Chefe PhD / Gerente de Qualis A1.
\end{descricao}

\begin{funcao}
No Core, este perfil apoia tarefas relacionadas a editor-Chefe PhD / Gerente de Qualis A1.. O orquestrador pode encaminhar-lhe o trabalho na etapa correspondente; a resposta retorna ao fluxo principal para integração e conferência de fontes, critérios e permissões. A ficha documenta a função, mas não comprova que o perfil tenha sido chamado ou executado a tarefa.
\end{funcao}

\begin{arquitetura}
\textbf{Modo de execução declarado:} subagente · \textbf{Permissões:} edit: deny; bash: deny
\end{arquitetura}

\elemento{01 agente diagnostico escopo (v1.0.0)}{\metainfo{Categoria}{Produção acadêmica} \hfill \metainfo{Tipo}{maswos-agent} \hfill \metainfo{Identificador}{01\_agente\_diagnostico\_escopo} \hfill \metainfo{Ficha}{01\_agente\_diagnostico\_escopo.md}}

\begin{descricao}
Transformar o tema do usuario em um problema de pesquisa claro, delimitado, auditavel e planejavel.
\end{descricao}

\begin{funcao}
No Core, este perfil apoia tarefas relacionadas a transformar o tema do usuario em um problema de pesquisa claro, delimitado, auditavel e planejavel.. O orquestrador pode encaminhar-lhe o trabalho na etapa correspondente; a resposta retorna ao fluxo principal para integração e conferência de fontes, critérios e permissões. A ficha documenta a função, mas não comprova que o perfil tenha sido chamado ou executado a tarefa.
\end{funcao}

\begin{arquitetura}
\textbf{Modo de execução declarado:} subagente · \textbf{Permissões:} edit: deny; bash: deny
\end{arquitetura}

\elemento{02 agente busca curadoria (v1.0.0)}{\metainfo{Categoria}{Produção acadêmica} \hfill \metainfo{Tipo}{maswos-agent} \hfill \metainfo{Identificador}{02\_agente\_busca\_curadoria} \hfill \metainfo{Ficha}{02\_agente\_busca\_curadoria.md}}

\begin{descricao}
Executar busca multipla, auditavel e suficientemente ampla para sustentar um artigo de alto nivel.
\end{descricao}

\begin{funcao}
No Core, este perfil apoia tarefas relacionadas a executar busca multipla, auditavel e suficientemente ampla para sustentar um artigo de alto nivel.. O orquestrador pode encaminhar-lhe o trabalho na etapa correspondente; a resposta retorna ao fluxo principal para integração e conferência de fontes, critérios e permissões. A ficha documenta a função, mas não comprova que o perfil tenha sido chamado ou executado a tarefa.
\end{funcao}

\begin{arquitetura}
\textbf{Modo de execução declarado:} subagente · \textbf{Permissões:} edit: deny; bash: deny
\end{arquitetura}

\elemento{03 agente evidencias citacoes (v1.0.0)}{\metainfo{Categoria}{Produção acadêmica} \hfill \metainfo{Tipo}{maswos-agent} \hfill \metainfo{Identificador}{03\_agente\_evidencias\_citacoes} \hfill \metainfo{Ficha}{03\_agente\_evidencias\_citacoes.md}}

\begin{descricao}
Converter fontes em evidencias localizadas, com funcao argumentativa e limites de uso explicitados.
\end{descricao}

\begin{funcao}
No Core, este perfil apoia tarefas relacionadas a converter fontes em evidencias localizadas, com funcao argumentativa e limites de uso explicitados.. O orquestrador pode encaminhar-lhe o trabalho na etapa correspondente; a resposta retorna ao fluxo principal para integração e conferência de fontes, critérios e permissões. A ficha documenta a função, mas não comprova que o perfil tenha sido chamado ou executado a tarefa.
\end{funcao}

\begin{arquitetura}
\textbf{Modo de execução declarado:} subagente · \textbf{Permissões:} edit: deny; bash: deny
\end{arquitetura}

\elemento{04 agente estrutura argumentativa (v1.0.0)}{\metainfo{Categoria}{Produção acadêmica} \hfill \metainfo{Tipo}{maswos-agent} \hfill \metainfo{Identificador}{04\_agente\_estrutura\_argumentativa} \hfill \metainfo{Ficha}{04\_agente\_estrutura\_argumentativa.md}}

\begin{descricao}
Projetar a espinha argumentativa do artigo para que cada secao execute uma funcao precisa.
\end{descricao}

\begin{funcao}
No Core, este perfil apoia tarefas relacionadas a projetar a espinha argumentativa do artigo para que cada secao execute uma funcao precisa.. O orquestrador pode encaminhar-lhe o trabalho na etapa correspondente; a resposta retorna ao fluxo principal para integração e conferência de fontes, critérios e permissões. A ficha documenta a função, mas não comprova que o perfil tenha sido chamado ou executado a tarefa.
\end{funcao}

\begin{arquitetura}
\textbf{Modo de execução declarado:} subagente · \textbf{Permissões:} edit: deny; bash: deny
\end{arquitetura}

\elemento{05 agente revisao literatura teoria (v1.0.0)}{\metainfo{Categoria}{Produção acadêmica} \hfill \metainfo{Tipo}{maswos-agent} \hfill \metainfo{Identificador}{05\_agente\_revisao\_literatura\_teoria} \hfill \metainfo{Ficha}{05\_agente\_revisao\_literatura\_teoria.md}}

\begin{descricao}
Construir uma revisao critica, dialogica e densa, em vez de uma lista de autores.
\end{descricao}

\begin{funcao}
No Core, este perfil apoia tarefas relacionadas a construir uma revisao critica, dialogica e densa, em vez de uma lista de autores.. O orquestrador pode encaminhar-lhe o trabalho na etapa correspondente; a resposta retorna ao fluxo principal para integração e conferência de fontes, critérios e permissões. A ficha documenta a função, mas não comprova que o perfil tenha sido chamado ou executado a tarefa.
\end{funcao}

\begin{arquitetura}
\textbf{Modo de execução declarado:} subagente · \textbf{Permissões:} edit: deny; bash: deny
\end{arquitetura}

\elemento{06 agente metodologia reprodutibilidade (v1.0.0)}{\metainfo{Categoria}{Produção acadêmica} \hfill \metainfo{Tipo}{maswos-agent} \hfill \metainfo{Identificador}{06\_agente\_metodologia\_reprodutibilidade} \hfill \metainfo{Ficha}{06\_agente\_metodologia\_reprodutibilidade.md}}

\begin{descricao}
Escrever uma metodologia replicavel, justificada e auditavel.
\end{descricao}

\begin{funcao}
No Core, este perfil apoia tarefas relacionadas a escrever uma metodologia replicavel, justificada e auditavel.. O orquestrador pode encaminhar-lhe o trabalho na etapa correspondente; a resposta retorna ao fluxo principal para integração e conferência de fontes, critérios e permissões. A ficha documenta a função, mas não comprova que o perfil tenha sido chamado ou executado a tarefa.
\end{funcao}

\begin{arquitetura}
\textbf{Modo de execução declarado:} subagente · \textbf{Permissões:} edit: deny; bash: deny
\end{arquitetura}

\elemento{07 agente estatistica analise (v1.0.0)}{\metainfo{Categoria}{Produção acadêmica} \hfill \metainfo{Tipo}{maswos-agent} \hfill \metainfo{Identificador}{07\_agente\_estatistica\_analise} \hfill \metainfo{Ficha}{07\_agente\_estatistica\_analise.md}}

\begin{descricao}
Validar a adequacao dos testes, a completude dos reportes e o limite inferencial do manuscrito.
\end{descricao}

\begin{funcao}
No Core, este perfil apoia tarefas relacionadas a validar a adequacao dos testes, a completude dos reportes e o limite inferencial do manuscrito.. O orquestrador pode encaminhar-lhe o trabalho na etapa correspondente; a resposta retorna ao fluxo principal para integração e conferência de fontes, critérios e permissões. A ficha documenta a função, mas não comprova que o perfil tenha sido chamado ou executado a tarefa.
\end{funcao}

\begin{arquitetura}
\textbf{Modo de execução declarado:} subagente · \textbf{Permissões:} edit: deny; bash: deny
\end{arquitetura}

\elemento{08 agente visualizacao evidencia grafica (v1.0.0)}{\metainfo{Categoria}{Produção acadêmica} \hfill \metainfo{Tipo}{maswos-agent} \hfill \metainfo{Identificador}{08\_agente\_visualizacao\_evidencia\_grafica} \hfill \metainfo{Ficha}{08\_agente\_visualizacao\_evidencia\_grafica.md}}

\begin{descricao}
Transformar resultados, metodo e comparacoes em visualizacoes com funcao argumentativa explicita.
\end{descricao}

\begin{funcao}
No Core, este perfil apoia tarefas relacionadas a transformar resultados, metodo e comparacoes em visualizacoes com funcao argumentativa explicita.. O orquestrador pode encaminhar-lhe o trabalho na etapa correspondente; a resposta retorna ao fluxo principal para integração e conferência de fontes, critérios e permissões. A ficha documenta a função, mas não comprova que o perfil tenha sido chamado ou executado a tarefa.
\end{funcao}

\begin{arquitetura}
\textbf{Modo de execução declarado:} subagente · \textbf{Permissões:} edit: deny; bash: deny
\end{arquitetura}

\elemento{09 agente resultados (v1.0.0)}{\metainfo{Categoria}{Produção acadêmica} \hfill \metainfo{Tipo}{maswos-agent} \hfill \metainfo{Identificador}{09\_agente\_resultados} \hfill \metainfo{Ficha}{09\_agente\_resultados.md}}

\begin{descricao}
Redigir os achados com precisao, ordem logica e neutralidade interpretativa.
\end{descricao}

\begin{funcao}
No Core, este perfil apoia tarefas relacionadas a redigir os achados com precisao, ordem logica e neutralidade interpretativa.. O orquestrador pode encaminhar-lhe o trabalho na etapa correspondente; a resposta retorna ao fluxo principal para integração e conferência de fontes, critérios e permissões. A ficha documenta a função, mas não comprova que o perfil tenha sido chamado ou executado a tarefa.
\end{funcao}

\begin{arquitetura}
\textbf{Modo de execução declarado:} subagente · \textbf{Permissões:} edit: deny; bash: deny
\end{arquitetura}

\elemento{10 agente discussao contribuicao (v1.0.0)}{\metainfo{Categoria}{Produção acadêmica} \hfill \metainfo{Tipo}{maswos-agent} \hfill \metainfo{Identificador}{10\_agente\_discussao\_contribuicao} \hfill \metainfo{Ficha}{10\_agente\_discussao\_contribuicao.md}}

\begin{descricao}
Interpretar os achados em dialogo real com a literatura, explicando significado, limite e contribuicao.
\end{descricao}

\begin{funcao}
No Core, este perfil apoia tarefas relacionadas a interpretar os achados em dialogo real com a literatura, explicando significado, limite e contribuicao.. O orquestrador pode encaminhar-lhe o trabalho na etapa correspondente; a resposta retorna ao fluxo principal para integração e conferência de fontes, critérios e permissões. A ficha documenta a função, mas não comprova que o perfil tenha sido chamado ou executado a tarefa.
\end{funcao}

\begin{arquitetura}
\textbf{Modo de execução declarado:} subagente · \textbf{Permissões:} edit: deny; bash: deny
\end{arquitetura}

\elemento{11 agente conclusao coerencia final (v1.0.0)}{\metainfo{Categoria}{Produção acadêmica} \hfill \metainfo{Tipo}{maswos-agent} \hfill \metainfo{Identificador}{11\_agente\_conclusao\_coerencia\_final} \hfill \metainfo{Ficha}{11\_agente\_conclusao\_coerencia\_final.md}}

\begin{descricao}
Fechar o artigo respondendo a pergunta central, aos objetivos e as hipoteses sem introduzir nada novo.
\end{descricao}

\begin{funcao}
No Core, este perfil apoia tarefas relacionadas a fechar o artigo respondendo a pergunta central, aos objetivos e as hipoteses sem introduzir nada novo.. O orquestrador pode encaminhar-lhe o trabalho na etapa correspondente; a resposta retorna ao fluxo principal para integração e conferência de fontes, critérios e permissões. A ficha documenta a função, mas não comprova que o perfil tenha sido chamado ou executado a tarefa.
\end{funcao}

\begin{arquitetura}
\textbf{Modo de execução declarado:} subagente · \textbf{Permissões:} edit: deny; bash: deny
\end{arquitetura}

\elemento{12 agente auditoria bibliografica abnt (v1.0.0)}{\metainfo{Categoria}{Produção acadêmica} \hfill \metainfo{Tipo}{maswos-agent} \hfill \metainfo{Identificador}{12\_agente\_auditoria\_bibliografica\_abnt} \hfill \metainfo{Ficha}{12\_agente\_auditoria\_bibliografica\_abnt.md}}

\begin{descricao}
Garantir consistencia total entre citacao no corpo, nota de rodape, referencia final e norma ABNT.
\end{descricao}

\begin{funcao}
No Core, este perfil apoia tarefas relacionadas a garantir consistencia total entre citacao no corpo, nota de rodape, referencia final e norma ABNT.. O orquestrador pode encaminhar-lhe o trabalho na etapa correspondente; a resposta retorna ao fluxo principal para integração e conferência de fontes, critérios e permissões. A ficha documenta a função, mas não comprova que o perfil tenha sido chamado ou executado a tarefa.
\end{funcao}

\begin{arquitetura}
\textbf{Modo de execução declarado:} subagente · \textbf{Permissões:} edit: deny; bash: deny
\end{arquitetura}

\elemento{13 agente qa qualis a1 (v1.0.0)}{\metainfo{Categoria}{Produção acadêmica} \hfill \metainfo{Tipo}{maswos-agent} \hfill \metainfo{Identificador}{13\_agente\_qa\_qualis\_a1} \hfill \metainfo{Ficha}{13\_agente\_qa\_qualis\_a1.md}}

\begin{descricao}
Avaliar se o manuscrito atende rigorosamente ao padrão MASWOS, o que significa EXIGIR
\end{descricao}

\begin{funcao}
No Core, este perfil apoia tarefas relacionadas a avaliar se o manuscrito atende rigorosamente ao padrão MASWOS, o que significa EXIGIR. O orquestrador pode encaminhar-lhe o trabalho na etapa correspondente; a resposta retorna ao fluxo principal para integração e conferência de fontes, critérios e permissões. A ficha documenta a função, mas não comprova que o perfil tenha sido chamado ou executado a tarefa.
\end{funcao}

\begin{arquitetura}
\textbf{Modo de execução declarado:} subagente · \textbf{Permissões:} edit: deny; bash: deny
\end{arquitetura}

\elemento{14 agente consistencia interna (v1.0.0)}{\metainfo{Categoria}{Produção acadêmica} \hfill \metainfo{Tipo}{maswos-agent} \hfill \metainfo{Identificador}{14\_agente\_consistencia\_interna} \hfill \metainfo{Ficha}{14\_agente\_consistencia\_interna.md}}

\begin{descricao}
Monitorar não apenas o contorno geral do artigo, mas aplicar auditoria micro e macro textual implacável (Padrão 10/10 MASWOS)
\end{descricao}

\begin{funcao}
No Core, este perfil apoia tarefas relacionadas a monitorar não apenas o contorno geral do artigo, mas aplicar auditoria micro e macro textual implacável (Padrão 10/10 MASWOS). O orquestrador pode encaminhar-lhe o trabalho na etapa correspondente; a resposta retorna ao fluxo principal para integração e conferência de fontes, critérios e permissões. A ficha documenta a função, mas não comprova que o perfil tenha sido chamado ou executado a tarefa.
\end{funcao}

\begin{arquitetura}
\textbf{Modo de execução declarado:} subagente · \textbf{Permissões:} edit: deny; bash: deny
\end{arquitetura}

\elemento{15 agente resumo abstract palavras chave (v1.0.0)}{\metainfo{Categoria}{Produção acadêmica} \hfill \metainfo{Tipo}{maswos-agent} \hfill \metainfo{Identificador}{15\_agente\_resumo\_abstract\_palavras\_chave} \hfill \metainfo{Ficha}{15\_agente\_resumo\_abstract\_palavras\_chave.md}}

\begin{descricao}
Produzir resumo em portugues, abstract em ingles e palavras-chave que sintetizem fielmente o manuscrito completo, sem inventar nada e sem perder densidade informativa.
\end{descricao}

\begin{funcao}
No Core, este perfil apoia tarefas relacionadas a produzir resumo em portugues, abstract em ingles e palavras-chave que sintetizem fielmente o manuscrito completo, sem inventar nada e sem perder densidade informativa.. O orquestrador pode encaminhar-lhe o trabalho na etapa correspondente; a resposta retorna ao fluxo principal para integração e conferência de fontes, critérios e permissões. A ficha documenta a função, mas não comprova que o perfil tenha sido chamado ou executado a tarefa.
\end{funcao}

\begin{arquitetura}
\textbf{Modo de execução declarado:} subagente · \textbf{Permissões:} edit: deny; bash: deny
\end{arquitetura}

\elemento{16 agente integracao editorial docx (v1.0.0)}{\metainfo{Categoria}{Produção acadêmica} \hfill \metainfo{Tipo}{maswos-agent} \hfill \metainfo{Identificador}{16\_agente\_integracao\_editorial\_docx} \hfill \metainfo{Ficha}{16\_agente\_integracao\_editorial\_docx.md}}

\begin{descricao}
Integrar todos os capitulos, apendices, figuras, referencias, resumo e metadados em um pacote editorial final coerente, pronto para revisao derradeira e conversao em DOCX.
\end{descricao}

\begin{funcao}
No Core, este perfil apoia tarefas relacionadas a integrar todos os capitulos, apendices, figuras, referencias, resumo e metadados em um pacote editorial final coerente, pronto para revisao derradeira e conversao em DOCX.. O orquestrador pode encaminhar-lhe o trabalho na etapa correspondente; a resposta retorna ao fluxo principal para integração e conferência de fontes, critérios e permissões. A ficha documenta a função, mas não comprova que o perfil tenha sido chamado ou executado a tarefa.
\end{funcao}

\begin{arquitetura}
\textbf{Modo de execução declarado:} subagente · \textbf{Permissões:} edit: deny; bash: deny
\end{arquitetura}

\elemento{17 agente framework reprodutivel ambientes (v1.0.0)}{\metainfo{Categoria}{Produção acadêmica} \hfill \metainfo{Tipo}{maswos-agent} \hfill \metainfo{Identificador}{17\_agente\_framework\_reprodutivel\_ambientes} \hfill \metainfo{Ficha}{17\_agente\_framework\_reprodutivel\_ambientes.md}}

\begin{descricao}
Transformar a parte computacional do artigo em um pacote reexecutavel, auditavel e explicito quanto a ambiente, dependencia, seed, fluxo e restricoes.
\end{descricao}

\begin{funcao}
No Core, este perfil apoia tarefas relacionadas a transformar a parte computacional do artigo em um pacote reexecutavel, auditavel e explicito quanto a ambiente, dependencia, seed, fluxo e restricoes.. O orquestrador pode encaminhar-lhe o trabalho na etapa correspondente; a resposta retorna ao fluxo principal para integração e conferência de fontes, critérios e permissões. A ficha documenta a função, mas não comprova que o perfil tenha sido chamado ou executado a tarefa.
\end{funcao}

\begin{arquitetura}
\textbf{Modo de execução declarado:} subagente · \textbf{Permissões:} edit: deny; bash: deny
\end{arquitetura}

\elemento{18 agente engenharia dados datasets proveniencia (v1.0.0)}{\metainfo{Categoria}{Produção acadêmica} \hfill \metainfo{Tipo}{maswos-agent} \hfill \metainfo{Identificador}{18\_agente\_engenharia\_dados\_datasets\_proveniencia} \hfill \metainfo{Ficha}{18\_agente\_engenharia\_dados\_datasets\_proveniencia.md}}

\begin{descricao}
Garantir que todo dado usado pelo artigo tenha origem, versao, papel analitico, restricao de uso e esquema documentalmente fechados.
\end{descricao}

\begin{funcao}
No Core, este perfil apoia tarefas relacionadas a garantir que todo dado usado pelo artigo tenha origem, versao, papel analitico, restricao de uso e esquema documentalmente fechados.. O orquestrador pode encaminhar-lhe o trabalho na etapa correspondente; a resposta retorna ao fluxo principal para integração e conferência de fontes, critérios e permissões. A ficha documenta a função, mas não comprova que o perfil tenha sido chamado ou executado a tarefa.
\end{funcao}

\begin{arquitetura}
\textbf{Modo de execução declarado:} subagente · \textbf{Permissões:} edit: deny; bash: deny
\end{arquitetura}

\elemento{19 agente auditoria codigo documentacao tecnica (v1.0.0)}{\metainfo{Categoria}{Produção acadêmica} \hfill \metainfo{Tipo}{maswos-agent} \hfill \metainfo{Identificador}{19\_agente\_auditoria\_codigo\_documentacao\_tecnica} \hfill \metainfo{Ficha}{19\_agente\_auditoria\_codigo\_documentacao\_tecnica.md}}

\begin{descricao}
Verificar se o codigo, os snippets, as bibliotecas e os pipelines tecnicos realmente correspondem ao que o artigo afirma e se estao ancorados em documentacao confiavel.
\end{descricao}

\begin{funcao}
No Core, este perfil apoia tarefas relacionadas a verificar se o codigo, os snippets, as bibliotecas e os pipelines tecnicos realmente correspondem ao que o artigo afirma e se estao ancorados em…. O orquestrador pode encaminhar-lhe o trabalho na etapa correspondente; a resposta retorna ao fluxo principal para integração e conferência de fontes, critérios e permissões. A ficha documenta a função, mas não comprova que o perfil tenha sido chamado ou executado a tarefa.
\end{funcao}

\begin{arquitetura}
\textbf{Modo de execução declarado:} subagente · \textbf{Permissões:} edit: deny; bash: deny
\end{arquitetura}

\elemento{20 agente estatistica avancada inferencia (v1.0.0)}{\metainfo{Categoria}{Produção acadêmica} \hfill \metainfo{Tipo}{maswos-agent} \hfill \metainfo{Identificador}{20\_agente\_estatistica\_avancada\_inferencia} \hfill \metainfo{Ficha}{20\_agente\_estatistica\_avancada\_inferencia.md}}

\begin{descricao}
Dar suporte inferencial de alto nivel a estudos quantitativos, incluindo cenarios frequentistas, bayesianos, causais, multilevel, temporais, espaciais, multivariados e de dados faltantes.
\end{descricao}

\begin{funcao}
No Core, este perfil apoia tarefas relacionadas a dar suporte inferencial de alto nivel a estudos quantitativos, incluindo cenarios frequentistas, bayesianos, causais, multilevel, temporais, espaciais, multivariados e de dados faltantes.. O orquestrador pode encaminhar-lhe o trabalho na etapa correspondente; a resposta retorna ao fluxo principal para integração e conferência de fontes, critérios e permissões. A ficha documenta a função, mas não comprova que o perfil tenha sido chamado ou executado a tarefa.
\end{funcao}

\begin{arquitetura}
\textbf{Modo de execução declarado:} subagente · \textbf{Permissões:} edit: deny; bash: deny
\end{arquitetura}

\elemento{21 agente matematica aplicada modelagem formal (v1.0.0)}{\metainfo{Categoria}{Produção acadêmica} \hfill \metainfo{Tipo}{maswos-agent} \hfill \metainfo{Identificador}{21\_agente\_matematica\_aplicada\_modelagem\_formal} \hfill \metainfo{Ficha}{21\_agente\_matematica\_aplicada\_modelagem\_formal.md}}

\begin{descricao}
Formalizar, derivar, verificar e delimitar modelos matematicos, equacoes, algoritmos numericos e estruturas simbolicas usadas no artigo.
\end{descricao}

\begin{funcao}
No Core, este perfil apoia tarefas relacionadas a formalizar, derivar, verificar e delimitar modelos matematicos, equacoes, algoritmos numericos e estruturas simbolicas usadas no artigo.. O orquestrador pode encaminhar-lhe o trabalho na etapa correspondente; a resposta retorna ao fluxo principal para integração e conferência de fontes, critérios e permissões. A ficha documenta a função, mas não comprova que o perfil tenha sido chamado ou executado a tarefa.
\end{funcao}

\begin{arquitetura}
\textbf{Modo de execução declarado:} subagente · \textbf{Permissões:} edit: deny; bash: deny
\end{arquitetura}

\elemento{22 agente ml dl datamining (v1.0.0)}{\metainfo{Categoria}{Produção acadêmica} \hfill \metainfo{Tipo}{maswos-agent} \hfill \metainfo{Identificador}{22\_agente\_ml\_dl\_datamining} \hfill \metainfo{Ficha}{22\_agente\_ml\_dl\_datamining.md}}

\begin{descricao}
Projetar, auditar e comparar pipelines de aprendizado de maquina e mineracao de dados com rigor de baseline, generalizacao, interpretabilidade e reproducibilidade.
\end{descricao}

\begin{funcao}
No Core, este perfil apoia tarefas relacionadas a projetar, auditar e comparar pipelines de aprendizado de maquina e mineracao de dados com rigor de baseline, generalizacao, interpretabilidade e reproducibilidade.. O orquestrador pode encaminhar-lhe o trabalho na etapa correspondente; a resposta retorna ao fluxo principal para integração e conferência de fontes, critérios e permissões. A ficha documenta a função, mas não comprova que o perfil tenha sido chamado ou executado a tarefa.
\end{funcao}

\begin{arquitetura}
\textbf{Modo de execução declarado:} subagente · \textbf{Permissões:} edit: deny; bash: deny
\end{arquitetura}

\elemento{23 agente bioinformatica omicas (v1.0.0)}{\metainfo{Categoria}{Produção acadêmica} \hfill \metainfo{Tipo}{maswos-agent} \hfill \metainfo{Identificador}{23\_agente\_bioinformatica\_omicas} \hfill \metainfo{Ficha}{23\_agente\_bioinformatica\_omicas.md}}

\begin{descricao}
Dar suporte especializado a pipelines de DNA, RNA, epigenomica, proteomica, metabolomica, single-cell e multiomicas com foco em rastreabilidade, normalizacao, controle de qualidade e interpretacao prudente.
\end{descricao}

\begin{funcao}
No Core, este perfil apoia tarefas relacionadas a dar suporte especializado a pipelines de DNA, RNA, epigenomica, proteomica, metabolomica, single-cell e multiomicas com foco em rastreabilidade, normalizacao, controle de qualidade e interpretacao…. O orquestrador pode encaminhar-lhe o trabalho na etapa correspondente; a resposta retorna ao fluxo principal para integração e conferência de fontes, critérios e permissões. A ficha documenta a função, mas não comprova que o perfil tenha sido chamado ou executado a tarefa.
\end{funcao}

\begin{arquitetura}
\textbf{Modo de execução declarado:} subagente · \textbf{Permissões:} edit: deny; bash: deny
\end{arquitetura}

\elemento{24 agente quimioinformatica modelagem molecular (v1.0.0)}{\metainfo{Categoria}{Produção acadêmica} \hfill \metainfo{Tipo}{maswos-agent} \hfill \metainfo{Identificador}{24\_agente\_quimioinformatica\_modelagem\_molecular} \hfill \metainfo{Ficha}{24\_agente\_quimioinformatica\_modelagem\_molecular.md}}

\begin{descricao}
Auditar e estruturar pipelines de chemometrics, descritores, QSAR/QSPR, docking, dinamica molecular, espectrometria e modelagem molecular aplicada ao artigo.
\end{descricao}

\begin{funcao}
No Core, este perfil apoia tarefas relacionadas a auditar e estruturar pipelines de chemometrics, descritores, QSAR/QSPR, docking, dinamica molecular, espectrometria e modelagem molecular aplicada ao artigo.. O orquestrador pode encaminhar-lhe o trabalho na etapa correspondente; a resposta retorna ao fluxo principal para integração e conferência de fontes, critérios e permissões. A ficha documenta a função, mas não comprova que o perfil tenha sido chamado ou executado a tarefa.
\end{funcao}

\begin{arquitetura}
\textbf{Modo de execução declarado:} subagente · \textbf{Permissões:} edit: deny; bash: deny
\end{arquitetura}

\elemento{25 agente ciencias sociais linguistica computacional (v1.0.0)}{\metainfo{Categoria}{Produção acadêmica} \hfill \metainfo{Tipo}{maswos-agent} \hfill \metainfo{Identificador}{25\_agente\_ciencias\_sociais\_linguistica\_computacional} \hfill \metainfo{Ficha}{25\_agente\_ciencias\_sociais\_linguistica\_computacional.md}}

\begin{descricao}
Dar suporte granular a estudos com surveys, psicometria, corpora, redes, NLP, estilometria, analise de discurso quantitativa e desenho observacional em ciencias sociais.
\end{descricao}

\begin{funcao}
No Core, este perfil apoia tarefas relacionadas a dar suporte granular a estudos com surveys, psicometria, corpora, redes, NLP, estilometria, analise de discurso quantitativa e desenho observacional em ciencias sociais.. O orquestrador pode encaminhar-lhe o trabalho na etapa correspondente; a resposta retorna ao fluxo principal para integração e conferência de fontes, critérios e permissões. A ficha documenta a função, mas não comprova que o perfil tenha sido chamado ou executado a tarefa.
\end{funcao}

\begin{arquitetura}
\textbf{Modo de execução declarado:} subagente · \textbf{Permissões:} edit: deny; bash: deny
\end{arquitetura}

\elemento{26 agente visao computacional multimodal (v1.0.0)}{\metainfo{Categoria}{Produção acadêmica} \hfill \metainfo{Tipo}{maswos-agent} \hfill \metainfo{Identificador}{26\_agente\_visao\_computacional\_multimodal} \hfill \metainfo{Ficha}{26\_agente\_visao\_computacional\_multimodal.md}}

\begin{descricao}
Estruturar e auditar pipelines de imagem, video, OCR, segmentacao, deteccao, classificacao visual e modelos multimodais com forte controle de dataset, pre-processamento e erro.
\end{descricao}

\begin{funcao}
No Core, este perfil apoia tarefas relacionadas a estruturar e auditar pipelines de imagem, video, OCR, segmentacao, deteccao, classificacao visual e modelos multimodais com forte controle de dataset, pre-processamento e erro.. O orquestrador pode encaminhar-lhe o trabalho na etapa correspondente; a resposta retorna ao fluxo principal para integração e conferência de fontes, critérios e permissões. A ficha documenta a função, mas não comprova que o perfil tenha sido chamado ou executado a tarefa.
\end{funcao}

\begin{arquitetura}
\textbf{Modo de execução declarado:} subagente · \textbf{Permissões:} edit: deny; bash: deny
\end{arquitetura}

\elemento{27 agente computacao quantica aplicada (v1.0.0)}{\metainfo{Categoria}{Produção acadêmica} \hfill \metainfo{Tipo}{maswos-agent} \hfill \metainfo{Identificador}{27\_agente\_computacao\_quantica\_aplicada} \hfill \metainfo{Ficha}{27\_agente\_computacao\_quantica\_aplicada.md}}

\begin{descricao}
Dar suporte especializado a artigos com circuitos quanticos, simuladores, modelos hibridos, kernels quanticos, variational algorithms e stacks como Qiskit, Cirq e PennyLane.
\end{descricao}

\begin{funcao}
No Core, este perfil apoia tarefas relacionadas a dar suporte especializado a artigos com circuitos quanticos, simuladores, modelos hibridos, kernels quanticos, variational algorithms e stacks como Qiskit, Cirq e PennyLane.. O orquestrador pode encaminhar-lhe o trabalho na etapa correspondente; a resposta retorna ao fluxo principal para integração e conferência de fontes, critérios e permissões. A ficha documenta a função, mas não comprova que o perfil tenha sido chamado ou executado a tarefa.
\end{funcao}

\begin{arquitetura}
\textbf{Modo de execução declarado:} subagente · \textbf{Permissões:} edit: deny; bash: deny
\end{arquitetura}

\elemento{28 agente benchmarking ablacao robustez (v1.0.0)}{\metainfo{Categoria}{Produção acadêmica} \hfill \metainfo{Tipo}{maswos-agent} \hfill \metainfo{Identificador}{28\_agente\_benchmarking\_ablacao\_robustez} \hfill \metainfo{Ficha}{28\_agente\_benchmarking\_ablacao\_robustez.md}}

\begin{descricao}
Verificar se o resultado computacional ou quantitativo resiste a comparacoes justas, seeds diferentes, ablations, sensibilidade, erro e cenarios adversos plausiveis.
\end{descricao}

\begin{funcao}
No Core, este perfil apoia tarefas relacionadas a verificar se o resultado computacional ou quantitativo resiste a comparacoes justas, seeds diferentes, ablations, sensibilidade, erro e cenarios adversos plausiveis.. O orquestrador pode encaminhar-lhe o trabalho na etapa correspondente; a resposta retorna ao fluxo principal para integração e conferência de fontes, critérios e permissões. A ficha documenta a função, mas não comprova que o perfil tenha sido chamado ou executado a tarefa.
\end{funcao}

\begin{arquitetura}
\textbf{Modo de execução declarado:} subagente · \textbf{Permissões:} edit: deny; bash: deny
\end{arquitetura}

\elemento{29 agente conformidade internacional (v1.0.0)}{\metainfo{Categoria}{Produção acadêmica} \hfill \metainfo{Tipo}{maswos-agent} \hfill \metainfo{Identificador}{29\_agente\_conformidade\_internacional} \hfill \metainfo{Ficha}{29\_agente\_conformidade\_internacional.md}}

\begin{descricao}
Garantir que o manuscrito cumpra rigorosamente as diretrizes internacionais exigidas por periódicos top-tier globais (Nature, Science, Lancet, IEEE, etc.) dependendo do tipo de estudo (revisão sistemática, ensaio...
\end{descricao}

\begin{funcao}
No Core, este perfil apoia tarefas relacionadas a garantir que o manuscrito cumpra rigorosamente as diretrizes internacionais exigidas por periódicos top-tier globais (Nature, Science, Lancet, IEEE, etc.) dependendo do tipo de estudo…. O orquestrador pode encaminhar-lhe o trabalho na etapa correspondente; a resposta retorna ao fluxo principal para integração e conferência de fontes, critérios e permissões. A ficha documenta a função, mas não comprova que o perfil tenha sido chamado ou executado a tarefa.
\end{funcao}

\begin{arquitetura}
\textbf{Modo de execução declarado:} subagente · \textbf{Permissões:} edit: deny; bash: deny
\end{arquitetura}

\elemento{30 agente traducao nativa proofreading (v1.0.0)}{\metainfo{Categoria}{Produção acadêmica} \hfill \metainfo{Tipo}{maswos-agent} \hfill \metainfo{Identificador}{30\_agente\_traducao\_nativa\_proofreading} \hfill \metainfo{Ficha}{30\_agente\_traducao\_nativa\_proofreading.md}}

\begin{descricao}
Prover fluência e estilo em língua inglesa no limite máximo exigido por bancas editoriais internacionais Tier-1.
\end{descricao}

\begin{funcao}
No Core, este perfil apoia tarefas relacionadas a prover fluência e estilo em língua inglesa no limite máximo exigido por bancas editoriais internacionais Tier-1.. O orquestrador pode encaminhar-lhe o trabalho na etapa correspondente; a resposta retorna ao fluxo principal para integração e conferência de fontes, critérios e permissões. A ficha documenta a função, mas não comprova que o perfil tenha sido chamado ou executado a tarefa.
\end{funcao}

\begin{arquitetura}
\textbf{Modo de execução declarado:} subagente · \textbf{Permissões:} edit: deny; bash: deny
\end{arquitetura}

\elemento{31 agente blind peer review emulado (v1.0.0)}{\metainfo{Categoria}{Produção acadêmica} \hfill \metainfo{Tipo}{maswos-agent} \hfill \metainfo{Identificador}{31\_agente\_blind\_peer\_review\_emulado} \hfill \metainfo{Ficha}{31\_agente\_blind\_peer\_review\_emulado.md}}

\begin{descricao}
Atuar como revisores contundentes (Reviewers 1, 2 e 3) e gerar comentários severos sobre falhas conceituais, metodológicas ou de interpretação antes de submeter ao periódico.
\end{descricao}

\begin{funcao}
No Core, este perfil apoia tarefas relacionadas a atuar como revisores contundentes (Reviewers 1, 2 e 3) e gerar comentários severos sobre falhas conceituais, metodológicas ou de interpretação antes de submeter ao…. O orquestrador pode encaminhar-lhe o trabalho na etapa correspondente; a resposta retorna ao fluxo principal para integração e conferência de fontes, critérios e permissões. A ficha documenta a função, mas não comprova que o perfil tenha sido chamado ou executado a tarefa.
\end{funcao}

\begin{arquitetura}
\textbf{Modo de execução declarado:} subagente · \textbf{Permissões:} edit: deny; bash: deny
\end{arquitetura}

\elemento{Agente de ética e ciência aberta (v1.0.0)}{\metainfo{Categoria}{Produção acadêmica} \hfill \metainfo{Tipo}{maswos-agent} \hfill \metainfo{Identificador}{32\_agente\_etica\_open\_science} \hfill \metainfo{Ficha}{32\_agente\_etica\_open\_science.md}}

\begin{descricao}
Protege os interesses legais e éticos da pesquisa e orienta a coleta e a guarda de dados pelos princípios FAIR — localizáveis, acessíveis, interoperáveis e reutilizáveis — e pela anonimização adequada.
\end{descricao}

\begin{funcao}
No Core, este perfil apoia tarefas relacionadas a protege os interesses legais e éticos da pesquisa e orienta a coleta e a guarda de dados pelos princípios FAIR — localizáveis, acessíveis, interoperáveis…. O orquestrador pode encaminhar-lhe o trabalho na etapa correspondente; a resposta retorna ao fluxo principal para integração e conferência de fontes, critérios e permissões. A ficha documenta a função, mas não comprova que o perfil tenha sido chamado ou executado a tarefa.
\end{funcao}

\begin{arquitetura}
\textbf{Modo de execução declarado:} subagente · \textbf{Permissões:} edit: deny; bash: deny
\end{arquitetura}

\elemento{33 agente automacao multi norma (v1.0.0)}{\metainfo{Categoria}{Produção acadêmica} \hfill \metainfo{Tipo}{maswos-agent} \hfill \metainfo{Identificador}{33\_agente\_automacao\_multi\_norma} \hfill \metainfo{Ficha}{33\_agente\_automacao\_multi\_norma.md}}

\begin{descricao}
Transmutar qualquer citação (seja via UUID, bibtex ou harvard style brutas) para o estilo exato e minucioso da norma global requerida pelo periódico.
\end{descricao}

\begin{funcao}
No Core, este perfil apoia tarefas relacionadas a transmutar qualquer citação (seja via UUID, bibtex ou harvard style brutas) para o estilo exato e minucioso da norma global requerida pelo periódico.. O orquestrador pode encaminhar-lhe o trabalho na etapa correspondente; a resposta retorna ao fluxo principal para integração e conferência de fontes, critérios e permissões. A ficha documenta a função, mas não comprova que o perfil tenha sido chamado ou executado a tarefa.
\end{funcao}

\begin{arquitetura}
\textbf{Modo de execução declarado:} subagente · \textbf{Permissões:} edit: deny; bash: deny
\end{arquitetura}

\elemento{34 agente identificacao conflitos similaridade (v1.0.0)}{\metainfo{Categoria}{Produção acadêmica} \hfill \metainfo{Tipo}{maswos-agent} \hfill \metainfo{Identificador}{34\_agente\_identificacao\_conflitos\_similaridade} \hfill \metainfo{Ficha}{34\_agente\_identificacao\_conflitos\_similaridade.md}}

\begin{descricao}
Emular o relatório frio, robótico e inexorável de sistemas de verificação léxica como iThenticate, Turnitin ou CrossCheck.
\end{descricao}

\begin{funcao}
No Core, este perfil apoia tarefas relacionadas a emular o relatório frio, robótico e inexorável de sistemas de verificação léxica como iThenticate, Turnitin ou CrossCheck.. O orquestrador pode encaminhar-lhe o trabalho na etapa correspondente; a resposta retorna ao fluxo principal para integração e conferência de fontes, critérios e permissões. A ficha documenta a função, mas não comprova que o perfil tenha sido chamado ou executado a tarefa.
\end{funcao}

\begin{arquitetura}
\textbf{Modo de execução declarado:} subagente · \textbf{Permissões:} edit: deny; bash: deny
\end{arquitetura}

\elemento{35 agente coleta datasets reais (v1.0.0)}{\metainfo{Categoria}{Produção acadêmica} \hfill \metainfo{Tipo}{maswos-agent} \hfill \metainfo{Identificador}{35\_agente\_coleta\_datasets\_reais} \hfill \metainfo{Ficha}{35\_agente\_coleta\_datasets\_reais.md}}

\begin{descricao}
Substituir absolutamente qualquer dataset sintético ou simulado por dados primários reais, coletados diretamente de APIs públicas oficiais, downloads de repositórios especializados, scraping estruturado de relatórios...
\end{descricao}

\begin{funcao}
No Core, este perfil apoia tarefas relacionadas a substituir absolutamente qualquer dataset sintético ou simulado por dados primários reais, coletados diretamente de APIs públicas oficiais, downloads de repositórios especializados, scraping estruturado de…. O orquestrador pode encaminhar-lhe o trabalho na etapa correspondente; a resposta retorna ao fluxo principal para integração e conferência de fontes, critérios e permissões. A ficha documenta a função, mas não comprova que o perfil tenha sido chamado ou executado a tarefa.
\end{funcao}

\begin{arquitetura}
\textbf{Modo de execução declarado:} subagente · \textbf{Permissões:} edit: deny; bash: deny
\end{arquitetura}

\elemento{36 agente exportacao latex pdf (v1.0.0)}{\metainfo{Categoria}{Produção acadêmica} \hfill \metainfo{Tipo}{maswos-agent} \hfill \metainfo{Identificador}{36\_agente\_exportacao\_latex\_pdf} \hfill \metainfo{Ficha}{36\_agente\_exportacao\_latex\_pdf.md}}

\begin{descricao}
Converter o manuscrito consolidado em formatos de submissão profissional: LaTeX (com classe adequada ao periódico), PDF compilado e DOCX.
\end{descricao}

\begin{funcao}
No Core, este perfil apoia tarefas relacionadas a converter o manuscrito consolidado em formatos de submissão profissional: LaTeX (com classe adequada ao periódico), PDF compilado e DOCX.. O orquestrador pode encaminhar-lhe o trabalho na etapa correspondente; a resposta retorna ao fluxo principal para integração e conferência de fontes, critérios e permissões. A ficha documenta a função, mas não comprova que o perfil tenha sido chamado ou executado a tarefa.
\end{funcao}

\begin{arquitetura}
\textbf{Modo de execução declarado:} subagente · \textbf{Permissões:} edit: deny; bash: deny
\end{arquitetura}

\elemento{37 agente apresentacao slides banca (v1.0.0)}{\metainfo{Categoria}{Produção acadêmica} \hfill \metainfo{Tipo}{maswos-agent} \hfill \metainfo{Identificador}{37\_agente\_apresentacao\_slides\_banca} \hfill \metainfo{Ficha}{37\_agente\_apresentacao\_slides\_banca.md}}

\begin{descricao}
Produzir uma apresentação acadêmica de altíssimo nível visual e argumentativo, pronta para defesa perante banca Qualis A1 ou conferência internacional, extraindo cirurgicamente os pontos-chave do manuscrito consolidado.
\end{descricao}

\begin{funcao}
No Core, este perfil apoia tarefas relacionadas a produzir uma apresentação acadêmica de altíssimo nível visual e argumentativo, pronta para defesa perante banca Qualis A1 ou conferência internacional, extraindo cirurgicamente os pontos-chave…. O orquestrador pode encaminhar-lhe o trabalho na etapa correspondente; a resposta retorna ao fluxo principal para integração e conferência de fontes, critérios e permissões. A ficha documenta a função, mas não comprova que o perfil tenha sido chamado ou executado a tarefa.
\end{funcao}

\begin{arquitetura}
\textbf{Modo de execução declarado:} subagente · \textbf{Permissões:} edit: deny; bash: deny
\end{arquitetura}

\elemento{38 agente montagem entrega final (v1.0.0)}{\metainfo{Categoria}{Produção acadêmica} \hfill \metainfo{Tipo}{maswos-agent} \hfill \metainfo{Identificador}{38\_agente\_montagem\_entrega\_final} \hfill \metainfo{Ficha}{38\_agente\_montagem\_entrega\_final.md}}

\begin{descricao}
Montar automaticamente, a partir dos fragmentos aprovados (manuscrito\_secoes/00 a 07), um documento completo, contínuo e pronto para submissão/defesa em múltiplos formatos.
\end{descricao}

\begin{funcao}
No Core, este perfil apoia tarefas relacionadas a montar automaticamente, a partir dos fragmentos aprovados (manuscrito\_secoes/00 a 07), um documento completo, contínuo e pronto para submissão/defesa em múltiplos formatos.. O orquestrador pode encaminhar-lhe o trabalho na etapa correspondente; a resposta retorna ao fluxo principal para integração e conferência de fontes, critérios e permissões. A ficha documenta a função, mas não comprova que o perfil tenha sido chamado ou executado a tarefa.
\end{funcao}

\begin{arquitetura}
\textbf{Modo de execução declarado:} subagente · \textbf{Permissões:} edit: deny; bash: deny
\end{arquitetura}

\elemento{39 agente metodologia multi paradigma (v1.0.0)}{\metainfo{Categoria}{Produção acadêmica} \hfill \metainfo{Tipo}{maswos-agent} \hfill \metainfo{Identificador}{39\_agente\_metodologia\_multi\_paradigma} \hfill \metainfo{Ficha}{39\_agente\_metodologia\_multi\_paradigma.md}}

\begin{descricao}
Garantir que a fundamentação e a execução metodológica do estudo estejam rigorosamente alinhadas ao paradigma epistemológico correto — seja ele quantitativo, qualitativo, misto, fenomenológico, etnográfico,...
\end{descricao}

\begin{funcao}
No Core, este perfil apoia tarefas relacionadas a garantir que a fundamentação e a execução metodológica do estudo estejam rigorosamente alinhadas ao paradigma epistemológico correto — seja ele quantitativo, qualitativo, misto, fenomenológico…. O orquestrador pode encaminhar-lhe o trabalho na etapa correspondente; a resposta retorna ao fluxo principal para integração e conferência de fontes, critérios e permissões. A ficha documenta a função, mas não comprova que o perfil tenha sido chamado ou executado a tarefa.
\end{funcao}

\begin{arquitetura}
\textbf{Modo de execução declarado:} subagente · \textbf{Permissões:} edit: deny; bash: deny
\end{arquitetura}

\elemento{40 agente marcos teoricos interpretacao (v1.0.0)}{\metainfo{Categoria}{Produção acadêmica} \hfill \metainfo{Tipo}{maswos-agent} \hfill \metainfo{Identificador}{40\_agente\_marcos\_teoricos\_interpretacao} \hfill \metainfo{Ficha}{40\_agente\_marcos\_teoricos\_interpretacao.md}}

\begin{descricao}
Garantir que o artigo/TCC esteja ancorado em um marco teórico coerente, que a interpretação de resultados siga os cânones da corrente adotada, e que a discussão dialogue com os teóricos corretos.
\end{descricao}

\begin{funcao}
No Core, este perfil apoia tarefas relacionadas a garantir que o artigo/TCC esteja ancorado em um marco teórico coerente, que a interpretação de resultados siga os cânones da corrente adotada, e que…. O orquestrador pode encaminhar-lhe o trabalho na etapa correspondente; a resposta retorna ao fluxo principal para integração e conferência de fontes, critérios e permissões. A ficha documenta a função, mas não comprova que o perfil tenha sido chamado ou executado a tarefa.
\end{funcao}

\begin{arquitetura}
\textbf{Modo de execução declarado:} subagente · \textbf{Permissões:} edit: deny; bash: deny
\end{arquitetura}

\elemento{41 agente gis geoprocessamento cartografia (v1.0.0)}{\metainfo{Categoria}{Produção acadêmica} \hfill \metainfo{Tipo}{maswos-agent} \hfill \metainfo{Identificador}{41\_agente\_gis\_geoprocessamento\_cartografia} \hfill \metainfo{Ficha}{41\_agente\_gis\_geoprocessamento\_cartografia.md}}

\begin{descricao}
Produzir, auditar e integrar elementos geoespaciais de alta qualidade ao manuscrito: mapas temáticos, cartas, plantas, modelos digitais de terreno, análise espacial, geovisualização e cartografia conforme normas...
\end{descricao}

\begin{funcao}
No Core, este perfil apoia tarefas relacionadas a produzir, auditar e integrar elementos geoespaciais de alta qualidade ao manuscrito: mapas temáticos, cartas, plantas, modelos digitais de terreno, análise espacial, geovisualização e cartografia…. O orquestrador pode encaminhar-lhe o trabalho na etapa correspondente; a resposta retorna ao fluxo principal para integração e conferência de fontes, critérios e permissões. A ficha documenta a função, mas não comprova que o perfil tenha sido chamado ou executado a tarefa.
\end{funcao}

\begin{arquitetura}
\textbf{Modo de execução declarado:} subagente · \textbf{Permissões:} edit: deny; bash: deny
\end{arquitetura}

\elemento{42 agente desenvolvedor cientista computacao (v1.0.0)}{\metainfo{Categoria}{Produção acadêmica} \hfill \metainfo{Tipo}{maswos-agent} \hfill \metainfo{Identificador}{42\_agente\_desenvolvedor\_cientista\_computacao} \hfill \metainfo{Ficha}{42\_agente\_desenvolvedor\_cientista\_computacao.md}}

\begin{descricao}
Projetar, gerar, auditar e otimizar TODO o código utilizado no manuscrito — desde scripts de coleta de dados até pipelines de ML, análise estatística, simulação e visualização.
\end{descricao}

\begin{funcao}
No Core, este perfil apoia tarefas relacionadas a projetar, gerar, auditar e otimizar TODO o código utilizado no manuscrito — desde scripts de coleta de dados até pipelines de ML, análise estatística…. O orquestrador pode encaminhar-lhe o trabalho na etapa correspondente; a resposta retorna ao fluxo principal para integração e conferência de fontes, critérios e permissões. A ficha documenta a função, mas não comprova que o perfil tenha sido chamado ou executado a tarefa.
\end{funcao}

\begin{arquitetura}
\textbf{Modo de execução declarado:} subagente · \textbf{Permissões:} edit: deny; bash: deny
\end{arquitetura}

\elemento{43 agente satelite bioinformatica omics (v1.0.0)}{\metainfo{Categoria}{Produção acadêmica} \hfill \metainfo{Tipo}{maswos-agent} \hfill \metainfo{Identificador}{43\_agente\_satelite\_bioinformatica\_omics} \hfill \metainfo{Ficha}{43\_agente\_satelite\_bioinformatica\_omics.md}}

\begin{descricao}
Coletar, tratar, processar e analisar dados de sensoriamento remoto (satélite, radar, LiDAR), dados biológicos (DNA, RNA, microarray, proteômica, metabolômica) e datasets de fontes especializadas — incluindo deep web...
\end{descricao}

\begin{funcao}
No Core, este perfil apoia tarefas relacionadas a coletar, tratar, processar e analisar dados de sensoriamento remoto (satélite, radar, LiDAR), dados biológicos (DNA, RNA, microarray, proteômica, metabolômica) e datasets de fontes especializadas…. O orquestrador pode encaminhar-lhe o trabalho na etapa correspondente; a resposta retorna ao fluxo principal para integração e conferência de fontes, critérios e permissões. A ficha documenta a função, mas não comprova que o perfil tenha sido chamado ou executado a tarefa.
\end{funcao}

\begin{arquitetura}
\textbf{Modo de execução declarado:} subagente · \textbf{Permissões:} edit: deny; bash: deny
\end{arquitetura}

\elemento{44 agente correcao textual qualis (v1.0.0)}{\metainfo{Categoria}{Produção acadêmica} \hfill \metainfo{Tipo}{maswos-agent} \hfill \metainfo{Identificador}{44\_agente\_correcao\_textual\_qualis} \hfill \metainfo{Ficha}{44\_agente\_correcao\_textual\_qualis.md}}

\begin{descricao}
Recebe textos reprovados ou com ressalvas dos agentes de validação (A13, A14) e os reescreve ativamente para atingir a nota 10/10.
\end{descricao}

\begin{funcao}
No Core, este perfil apoia tarefas relacionadas a recebe textos reprovados ou com ressalvas dos agentes de validação (A13, A14) e os reescreve ativamente para atingir a nota 10/10.. O orquestrador pode encaminhar-lhe o trabalho na etapa correspondente; a resposta retorna ao fluxo principal para integração e conferência de fontes, critérios e permissões. A ficha documenta a função, mas não comprova que o perfil tenha sido chamado ou executado a tarefa.
\end{funcao}

\begin{arquitetura}
\textbf{Modo de execução declarado:} subagente · \textbf{Permissões:} edit: deny; bash: deny
\end{arquitetura}

\elemento{45 agente refinamento argumentacao (v1.0.0)}{\metainfo{Categoria}{Produção acadêmica} \hfill \metainfo{Tipo}{maswos-agent} \hfill \metainfo{Identificador}{45\_agente\_refinamento\_argumentacao} \hfill \metainfo{Ficha}{45\_agente\_refinamento\_argumentacao.md}}

\begin{descricao}
Eleva a nota do critério 'Diálogo Crítico e Contribuição' (I2 e B1 da rubrica) para 10/10.
\end{descricao}

\begin{funcao}
No Core, este perfil apoia tarefas relacionadas a eleva a nota do critério 'Diálogo Crítico e Contribuição' (I2 e B1 da rubrica) para 10/10.. O orquestrador pode encaminhar-lhe o trabalho na etapa correspondente; a resposta retorna ao fluxo principal para integração e conferência de fontes, critérios e permissões. A ficha documenta a função, mas não comprova que o perfil tenha sido chamado ou executado a tarefa.
\end{funcao}

\begin{arquitetura}
\textbf{Modo de execução declarado:} subagente · \textbf{Permissões:} edit: deny; bash: deny
\end{arquitetura}

\elemento{Redator de acadêmico (v1.0.0)}{\metainfo{Categoria}{Produção acadêmica} \hfill \metainfo{Tipo}{maswos-agent} \hfill \metainfo{Identificador}{academic\_writer} \hfill \metainfo{Ficha}{academic\_writer.md}}

\begin{descricao}
Agente de redação acadêmica do Core. Produz manuscritos com estrutura IMRAD, citações rastreáveis e conformidade a ABNT (NBR 6023), APA ou Vancouver, sob o protocolo SDD/TDD do orquestrador marceloclaro. - academic\_writing — redação de manuscritos e capítulos em IMRAD - abnt — normalização de referências e citações segundo a NBR - imrad — estrutura Introdução, Métodos, Resultados, Discussão - qualis\_a1 — densidade argumentativa compatível com periódicos de alto impacto Não promete aceitação editorial, nota máxima ou indexação Qualis A1: essas conclusões pertencem ao periódico e à banca, não ao agente. Toda afirmação de mérito deve vir acompanhada de fonte verificável (R110, Módulo 6). - Prompt canônico: agents/
\end{descricao}

\begin{funcao}
No Core, este perfil apoia tarefas relacionadas a agente de redação acadêmica do Core.. O orquestrador pode encaminhar-lhe o trabalho na etapa correspondente; a resposta retorna ao fluxo principal para integração e conferência de fontes, critérios e permissões. A ficha documenta a função, mas não comprova que o perfil tenha sido chamado ou executado a tarefa.
\end{funcao}

\begin{arquitetura}
\textbf{Modo de execução declarado:} subagente · \textbf{Permissões:} edit: allow; bash: deny
\end{arquitetura}

\elemento{Dispatcher ATIVACAO (v1.0.0)}{\metainfo{Categoria}{Produção acadêmica} \hfill \metainfo{Tipo}{maswos-agent} \hfill \metainfo{Identificador}{DISPATCHER\_ATIVACAO} \hfill \metainfo{Ficha}{DISPATCHER\_ATIVACAO.md}}

\begin{descricao}
Localizacao: criador-artigo\textbackslash\{\}agents\textbackslash\{\}DISPATCHER\_ATIVACAO.md Tipo: Agente MASWOS (Criador de Artigos v2) Categoria: Pipeline Academico Qualis A1 Este agente faz parte do pipeline MASWOS para producao de artigos academicos. Para detalhes completos, consulte o arquivo original em criador-artigo\textbackslash\{\}agents\textbackslash\{\}DISPATCHER\_ATIVACAO.md.
\end{descricao}

\begin{funcao}
No Core, este perfil apoia tarefas relacionadas a localizacao: criador-artigo\textbackslash\{\}agents\textbackslash\{\}DISPATCHER\_ATIVACAO.md Tipo: Agente MASWOS (Criador de Artigos v2) Categoria: Pipeline Academico Qualis A1 Este agente faz parte do pipeline MASWOS para producao de…. O orquestrador pode encaminhar-lhe o trabalho na etapa correspondente; a resposta retorna ao fluxo principal para integração e conferência de fontes, critérios e permissões. A ficha documenta a função, mas não comprova que o perfil tenha sido chamado ou executado a tarefa.
\end{funcao}

\begin{arquitetura}
\textbf{Modo de execução declarado:} subagente · \textbf{Permissões:} edit: deny; bash: deny
\end{arquitetura}

\elemento{Mira copywriter (v1.0.0)}{\metainfo{Categoria}{Produção acadêmica} \hfill \metainfo{Tipo}{mira-agent} \hfill \metainfo{Identificador}{mira-copywriter} \hfill \metainfo{Ficha}{mira-copywriter.md}}

\begin{descricao}
Agente que refina títulos, textos curtos, legendas e mensagens visuais para maximizar clareza e impacto em apresentações MIRA.
\end{descricao}

\begin{funcao}
No Core, este perfil apoia tarefas relacionadas a agente que refina títulos, textos curtos, legendas e mensagens visuais para maximizar clareza e impacto em apresentações MIRA.. O orquestrador pode encaminhar-lhe o trabalho na etapa correspondente; a resposta retorna ao fluxo principal para integração e conferência de fontes, critérios e permissões. A ficha documenta a função, mas não comprova que o perfil tenha sido chamado ou executado a tarefa.
\end{funcao}

\begin{arquitetura}
\textbf{Modo de execução declarado:} subagente · \textbf{Permissões:} edit: deny; bash: deny
\end{arquitetura}

\elemento{Readme (v1.0.0)}{\metainfo{Categoria}{Produção acadêmica} \hfill \metainfo{Tipo}{maswos-agent} \hfill \metainfo{Identificador}{README} \hfill \metainfo{Ficha}{README.md}}

\begin{descricao}
Localizacao: criador-artigo\textbackslash\{\}agents\textbackslash\{\}README.md Tipo: Agente MASWOS (Criador de Artigos v2) Categoria: Pipeline Academico Qualis A1 Este agente faz parte do pipeline MASWOS para producao de artigos academicos. Para detalhes completos, consulte o arquivo original em criador-artigo\textbackslash\{\}agents\textbackslash\{\}README.md.
\end{descricao}

\begin{funcao}
No Core, este perfil apoia tarefas relacionadas a localizacao: criador-artigo\textbackslash\{\}agents\textbackslash\{\}README.md Tipo: Agente MASWOS (Criador de Artigos v2) Categoria: Pipeline Academico Qualis A1 Este agente faz parte do pipeline MASWOS para producao de…. O orquestrador pode encaminhar-lhe o trabalho na etapa correspondente; a resposta retorna ao fluxo principal para integração e conferência de fontes, critérios e permissões. A ficha documenta a função, mas não comprova que o perfil tenha sido chamado ou executado a tarefa.
\end{funcao}

\begin{arquitetura}
\textbf{Modo de execução declarado:} subagente · \textbf{Permissões:} edit: deny; bash: deny
\end{arquitetura}

\elemento{Synthetic University — Orquestrador Acadêmico Transversal (v1.0.0)}{\metainfo{Categoria}{Produção acadêmica} \hfill \metainfo{Tipo}{orchestrator} \hfill \metainfo{Identificador}{Synthetic University — Orquestrador Acadêmico Transversal} \hfill \metainfo{Ficha}{university\_synthetic.md}}

\begin{descricao}
Orquestrador Central da Universidade Sintética Transversal (SPEC-935). Coordena 10 faculdades, 40+ professores especialistas, motor combinatorial que testa 10.000+ combinações de conceitos, correlator interdisciplinar e gerador de teses PhD-level para descobrir novas correlações e combinações viáveis entre todos os domínios do conhecimento. 1. Ciências Humanas — Filosofia, Psicologia, Sociologia, Antropologia, Educação 2. Ciências Sociais Aplicadas — Economia, Política, Direito, Comunicação, RI 3. Engenharia — Software, Sistemas, Elétrica, Mecânica, Computação 4. Letras \& Linguística — Literatura, Linguística, Semiótica, Tradução, Filologia 5. História — História, Arqueologia, Historiografia, Patrimônio 6. Ciên
\end{descricao}

\begin{funcao}
No Core, este perfil apoia tarefas relacionadas a orquestrador Central da Universidade Sintética Transversal (SPEC-935).. O orquestrador pode encaminhar-lhe o trabalho na etapa correspondente; a resposta retorna ao fluxo principal para integração e conferência de fontes, critérios e permissões. A ficha documenta a função, mas não comprova que o perfil tenha sido chamado ou executado a tarefa.
\end{funcao}

\begin{arquitetura}
\textbf{Modo de execução declarado:} subagente
\end{arquitetura}

\elemento{Template HANDOFF (v1.0.0)}{\metainfo{Categoria}{Produção acadêmica} \hfill \metainfo{Tipo}{maswos-agent} \hfill \metainfo{Identificador}{TEMPLATE\_HANDOFF} \hfill \metainfo{Ficha}{TEMPLATE\_HANDOFF.md}}

\begin{descricao}
Localizacao: criador-artigo\textbackslash\{\}agents\textbackslash\{\}TEMPLATE\_HANDOFF.md Tipo: Agente MASWOS (Criador de Artigos v2) Categoria: Pipeline Academico Qualis A1 Este agente faz parte do pipeline MASWOS para producao de artigos academicos. Para detalhes completos, consulte o arquivo original em criador-artigo\textbackslash\{\}agents\textbackslash\{\}TEMPLATE\_HANDOFF.md.
\end{descricao}

\begin{funcao}
No Core, este perfil apoia tarefas relacionadas a localizacao: criador-artigo\textbackslash\{\}agents\textbackslash\{\}TEMPLATE\_HANDOFF.md Tipo: Agente MASWOS (Criador de Artigos v2) Categoria: Pipeline Academico Qualis A1 Este agente faz parte do pipeline MASWOS para producao de…. O orquestrador pode encaminhar-lhe o trabalho na etapa correspondente; a resposta retorna ao fluxo principal para integração e conferência de fontes, critérios e permissões. A ficha documenta a função, mas não comprova que o perfil tenha sido chamado ou executado a tarefa.
\end{funcao}

\begin{arquitetura}
\textbf{Modo de execução declarado:} subagente · \textbf{Permissões:} edit: deny; bash: deny
\end{arquitetura}

\clearpage
\section{Categoria: Pesquisa e evidências (14 perfis)}

\begin{descricao}
Agrupa perfis de busca, recuperação e síntese de evidências. No Core, suas fontes alimentam a análise e a escrita científica e devem permanecer rastreáveis (Módulo 6).
\end{descricao}

\elemento{Auditoria reprodutibilidade (v1.0.0)}{\metainfo{Categoria}{Pesquisa e evidências} \hfill \metainfo{Tipo}{mira-agent} \hfill \metainfo{Identificador}{auditoria-reprodutibilidade} \hfill \metainfo{Ficha}{48\_agente\_auditoria\_reprodutibilidade.md}}

\begin{descricao}
Audita reprodutibilidade com repro\_audit, ReproRepo e PRISMA-trAIce, sem alegar prova. Capacidades descritas: Confere seeds, hiperparâmetros, splits, dependências pinadas e proveniência via padrão repro\_audit; Cruza issues reais do GitHub no padrão ReproRepo para localizar região semântica da falha
\end{descricao}

\begin{funcao}
No Core, este perfil apoia tarefas relacionadas a audita reprodutibilidade com repro\_audit, ReproRepo e PRISMA-trAIce, sem alegar prova.. O orquestrador pode encaminhar-lhe o trabalho na etapa correspondente; a resposta retorna ao fluxo principal para integração e conferência de fontes, critérios e permissões. A ficha documenta a função, mas não comprova que o perfil tenha sido chamado ou executado a tarefa.
\end{funcao}

\begin{arquitetura}
\textbf{Modo de execução declarado:} subagente
\end{arquitetura}

\elemento{Bibtex crossref auditor (v1.0.0)}{\metainfo{Categoria}{Pesquisa e evidências} \hfill \metainfo{Tipo}{mira-agent} \hfill \metainfo{Identificador}{bibtex-crossref-auditor} \hfill \metainfo{Ficha}{bibtex-crossref-auditor.md}}

\begin{descricao}
Audita referências bibliográficas via Crossref/DOI antes de gravar no .bib
\end{descricao}

\begin{funcao}
No Core, este perfil apoia tarefas relacionadas a audita referências bibliográficas via Crossref/DOI antes de gravar no .bib. O orquestrador pode encaminhar-lhe o trabalho na etapa correspondente; a resposta retorna ao fluxo principal para integração e conferência de fontes, critérios e permissões. A ficha documenta a função, mas não comprova que o perfil tenha sido chamado ou executado a tarefa.
\end{funcao}

\begin{arquitetura}
\textbf{Modo de execução declarado:} subagente · \textbf{Permissões:} edit: deny; bash: deny
\end{arquitetura}

\elemento{Inferencia causal did iv rdd (v1.0.0)}{\metainfo{Categoria}{Pesquisa e evidências} \hfill \metainfo{Tipo}{mira-agent} \hfill \metainfo{Identificador}{inferencia-causal-did-iv-rdd} \hfill \metainfo{Ficha}{inferencia-causal-did-iv-rdd.md}}

\begin{descricao}
Desenha e executa DID, IV-2SLS e RDD com suposições declaradas e diagnósticos, sem alegar causalidade sem desenho
\end{descricao}

\begin{funcao}
No Core, este perfil apoia tarefas relacionadas a desenha e executa DID, IV-2SLS e RDD com suposições declaradas e diagnósticos, sem alegar causalidade sem desenho. O orquestrador pode encaminhar-lhe o trabalho na etapa correspondente; a resposta retorna ao fluxo principal para integração e conferência de fontes, critérios e permissões. A ficha documenta a função, mas não comprova que o perfil tenha sido chamado ou executado a tarefa.
\end{funcao}

\begin{arquitetura}
\textbf{Modo de execução declarado:} subagente · \textbf{Permissões:} edit: deny; bash: deny
\end{arquitetura}

\elemento{Laboratorio reproduzivel (v1.0.0)}{\metainfo{Categoria}{Pesquisa e evidências} \hfill \metainfo{Tipo}{mira-agent} \hfill \metainfo{Identificador}{laboratorio-reproduzivel} \hfill \metainfo{Ficha}{47\_agente\_laboratorio\_reproduzivel.md}}

\begin{descricao}
Ergue laboratório reproduzível no modelo rasilab com issues, versão e contêineres. Capacidades descritas: Organiza hipótese, desenho, dados e análise em issues rastreáveis com reabertura auditada; Prescreve Git + GitHub Packages + Dockerfile para ambiente reproduzível, sem executar build
\end{descricao}

\begin{funcao}
No Core, este perfil apoia tarefas relacionadas a ergue laboratório reproduzível no modelo rasilab com issues, versão e contêineres.. O orquestrador pode encaminhar-lhe o trabalho na etapa correspondente; a resposta retorna ao fluxo principal para integração e conferência de fontes, critérios e permissões. A ficha documenta a função, mas não comprova que o perfil tenha sido chamado ou executado a tarefa.
\end{funcao}

\begin{arquitetura}
\textbf{Modo de execução declarado:} subagente
\end{arquitetura}

\elemento{Curador da paisagem de agentes (v1.1.0)}{\metainfo{Categoria}{Pesquisa e evidências} \hfill \metainfo{Tipo}{curation-agent} \hfill \metainfo{Identificador}{landscape-curator} \hfill \metainfo{Ficha}{landscape-curator.md}}

\begin{descricao}
Curador da paisagem de agentes externos — consome manifests curados de múltiplas coleções (500-AI-Agents-Projects: 20 agentes auto-contidos, MIT, R482; awesome-llm-apps: 15 templates representativos, Apache-2.0, R521), cruza com o catálogo do Core (197 agent cards) por afinidade lexical auditável e gera LANDSCAPE\_REPORT.md honestos, sem copiar código de terceiros, sem prometer integração e com Observatório para entradas adversarial (veredito não-adotar/observar).
\end{descricao}

\begin{funcao}
No Core, este perfil apoia tarefas relacionadas a curador da paisagem de agentes externos — consome manifests curados de múltiplas coleções (500-AI-Agents-Projects: 20 agentes auto-contidos, MIT, R482. O orquestrador pode encaminhar-lhe o trabalho na etapa correspondente; a resposta retorna ao fluxo principal para integração e conferência de fontes, critérios e permissões. A ficha documenta a função, mas não comprova que o perfil tenha sido chamado ou executado a tarefa.
\end{funcao}

\begin{arquitetura}
\textbf{Modo de execução declarado:} subagente · \textbf{Permissões:} edit: deny; bash: deny
\end{arquitetura}

\elemento{Mira extract (v1.0.0)}{\metainfo{Categoria}{Pesquisa e evidências} \hfill \metainfo{Tipo}{mira-agent} \hfill \metainfo{Identificador}{mira-extract} \hfill \metainfo{Ficha}{mira-extract.md}}

\begin{descricao}
Extrai briefing e estrutura inicial a partir das fontes do MIRA. Capacidades descritas: Estrutura inicial a partir das fontes do mira
\end{descricao}

\begin{funcao}
No Core, este perfil apoia tarefas relacionadas a extrai briefing e estrutura inicial a partir das fontes do MIRA.. O orquestrador pode encaminhar-lhe o trabalho na etapa correspondente; a resposta retorna ao fluxo principal para integração e conferência de fontes, critérios e permissões. A ficha documenta a função, mas não comprova que o perfil tenha sido chamado ou executado a tarefa.
\end{funcao}

\begin{arquitetura}
\textbf{Modo de execução declarado:} subagente · \textbf{Permissões:} edit: deny; bash: deny
\end{arquitetura}

\elemento{Mira get videos (v1.0.0)}{\metainfo{Categoria}{Pesquisa e evidências} \hfill \metainfo{Tipo}{mira-agent} \hfill \metainfo{Identificador}{mira-get-videos} \hfill \metainfo{Ficha}{mira-get-videos.md}}

\begin{descricao}
Seleciona e organiza fundos em vídeo para apresentações MIRA. Capacidades descritas: Organiza fundos em vídeo para apresentações mira
\end{descricao}

\begin{funcao}
No Core, este perfil apoia tarefas relacionadas a seleciona e organiza fundos em vídeo para apresentações MIRA.. O orquestrador pode encaminhar-lhe o trabalho na etapa correspondente; a resposta retorna ao fluxo principal para integração e conferência de fontes, critérios e permissões. A ficha documenta a função, mas não comprova que o perfil tenha sido chamado ou executado a tarefa.
\end{funcao}

\begin{arquitetura}
\textbf{Modo de execução declarado:} subagente · \textbf{Permissões:} edit: deny; bash: deny
\end{arquitetura}

\elemento{Mira references (v1.0.0)}{\metainfo{Categoria}{Pesquisa e evidências} \hfill \metainfo{Tipo}{mira-agent} \hfill \metainfo{Identificador}{mira-references} \hfill \metainfo{Ficha}{mira-references.md}}

\begin{descricao}
Vincula fontes e referências ao tema do deck MIRA. Capacidades descritas: Referências ao tema do deck mira
\end{descricao}

\begin{funcao}
No Core, este perfil apoia tarefas relacionadas a vincula fontes e referências ao tema do deck MIRA.. O orquestrador pode encaminhar-lhe o trabalho na etapa correspondente; a resposta retorna ao fluxo principal para integração e conferência de fontes, critérios e permissões. A ficha documenta a função, mas não comprova que o perfil tenha sido chamado ou executado a tarefa.
\end{funcao}

\begin{arquitetura}
\textbf{Modo de execução declarado:} subagente · \textbf{Permissões:} edit: deny; bash: deny
\end{arquitetura}

\elemento{Mira survey (v1.0.0)}{\metainfo{Categoria}{Pesquisa e evidências} \hfill \metainfo{Tipo}{mira-agent} \hfill \metainfo{Identificador}{mira-survey} \hfill \metainfo{Ficha}{mira-survey.md}}

\begin{descricao}
Agente que projeta enquetes, perguntas ao vivo e mecanismos de participação da plateia para decks MIRA.
\end{descricao}

\begin{funcao}
No Core, este perfil apoia tarefas relacionadas a agente que projeta enquetes, perguntas ao vivo e mecanismos de participação da plateia para decks MIRA.. O orquestrador pode encaminhar-lhe o trabalho na etapa correspondente; a resposta retorna ao fluxo principal para integração e conferência de fontes, critérios e permissões. A ficha documenta a função, mas não comprova que o perfil tenha sido chamado ou executado a tarefa.
\end{funcao}

\begin{arquitetura}
\textbf{Modo de execução declarado:} subagente · \textbf{Permissões:} edit: deny; bash: deny
\end{arquitetura}

\elemento{Modelagem bayesiana mista (v1.0.0)}{\metainfo{Categoria}{Pesquisa e evidências} \hfill \metainfo{Tipo}{mira-agent} \hfill \metainfo{Identificador}{modelagem-bayesiana-mista} \hfill \metainfo{Ficha}{modelagem-bayesiana-mista.md}}

\begin{descricao}
Atualização bayesiana conjugada com prior explícito e modelos mistos leves (ICC, EP cluster-robustos)
\end{descricao}

\begin{funcao}
No Core, este perfil apoia tarefas relacionadas a atualização bayesiana conjugada com prior explícito e modelos mistos leves (ICC, EP cluster-robustos). O orquestrador pode encaminhar-lhe o trabalho na etapa correspondente; a resposta retorna ao fluxo principal para integração e conferência de fontes, critérios e permissões. A ficha documenta a função, mas não comprova que o perfil tenha sido chamado ou executado a tarefa.
\end{funcao}

\begin{arquitetura}
\textbf{Modo de execução declarado:} subagente · \textbf{Permissões:} edit: deny; bash: deny
\end{arquitetura}

\elemento{Pesquisador polimata (v1.0.0)}{\metainfo{Categoria}{Pesquisa e evidências} \hfill \metainfo{Tipo}{mira-agent} \hfill \metainfo{Identificador}{pesquisador-polimata} \hfill \metainfo{Ficha}{46\_agente\_pesquisador\_polimata.md}}

\begin{descricao}
Roteia tipo de raciocínio ao lab R711 adequado, sem executar código de terceiros. Capacidades descritas: Roteamento tipo de raciocínio para lab; Mapeia pergunta ao lab da allowlist R711 pelo tipo exigido, com licença e limite declarados
\end{descricao}

\begin{funcao}
No Core, este perfil apoia tarefas relacionadas a roteia tipo de raciocínio ao lab R711 adequado, sem executar código de terceiros.. O orquestrador pode encaminhar-lhe o trabalho na etapa correspondente; a resposta retorna ao fluxo principal para integração e conferência de fontes, critérios e permissões. A ficha documenta a função, mas não comprova que o perfil tenha sido chamado ou executado a tarefa.
\end{funcao}

\begin{arquitetura}
\textbf{Modo de execução declarado:} subagente
\end{arquitetura}

\elemento{Py PISearcher (v1.0.0)}{\metainfo{Categoria}{Pesquisa e evidências} \hfill \metainfo{Tipo}{especialista} \hfill \metainfo{Identificador}{PyPISearcher} \hfill \metainfo{Ficha}{pypi-searcher.md}}

\begin{descricao}
Especialista em descoberta, análise e recomendação de bibliotecas Python do Python Package Index (PyPI) para o ecossistema OpenCode v4.2. Utiliza duas estratégias de busca (scraping HTML + API JSON oficial) e integra-se ao DecisionNode para persistência de decisões arquiteturais. PyPISearcher research/discovery OpenCode v4.2 PT-BR formal (saída obrigatória) Alta — contexto em chinês, saída em português Ferramenta Primária: pypi\_search.py Utilize SEMPRE o script local skills/tooling/pypi\_search.py como primeira estratégia de busca. O script combina scraping da página de busca + enriquecimento via API JSON oficial do PyPI. python skills/tooling/pypi\_search.py "termo de busca" --limit N
\end{descricao}

\begin{funcao}
No Core, este perfil apoia tarefas relacionadas a especialista em descoberta, análise e recomendação de bibliotecas Python do Python Package Index (PyPI) para o ecossistema OpenCode v4.2.. O orquestrador pode encaminhar-lhe o trabalho na etapa correspondente; a resposta retorna ao fluxo principal para integração e conferência de fontes, critérios e permissões. A ficha documenta a função, mas não comprova que o perfil tenha sido chamado ou executado a tarefa.
\end{funcao}

\begin{arquitetura}
\textbf{Modo de execução declarado:} subagente
\end{arquitetura}

\elemento{Pesquisador de (v1.0.0)}{\metainfo{Categoria}{Pesquisa e evidências} \hfill \metainfo{Tipo}{especialista} \hfill \metainfo{Identificador}{researcher} \hfill \metainfo{Ficha}{researcher.md}}

\begin{descricao}
Agente de pesquisa profunda, síntese de literatura e verificação de fatos.
\end{descricao}

\begin{funcao}
No Core, este perfil apoia tarefas relacionadas a agente de pesquisa profunda, síntese de literatura e verificação de fatos.. O orquestrador pode encaminhar-lhe o trabalho na etapa correspondente; a resposta retorna ao fluxo principal para integração e conferência de fontes, critérios e permissões. A ficha documenta a função, mas não comprova que o perfil tenha sido chamado ou executado a tarefa.
\end{funcao}

\begin{arquitetura}
\textbf{Modo de execução declarado:} subagente · \textbf{Permissões:} edit: deny; bash: deny
\end{arquitetura}

\elemento{Tdah gap hunter (v1.0.0)}{\metainfo{Categoria}{Pesquisa e evidências} \hfill \metainfo{Tipo}{mira-agent} \hfill \metainfo{Identificador}{tdah-gap-hunter} \hfill \metainfo{Ficha}{tdah-gap-hunter.md}}

\begin{descricao}
Localiza o artigo-gap que norteia pesquisa sobre TDAH, brincar livre e natureza em pré-escolares
\end{descricao}

\begin{funcao}
No Core, este perfil apoia tarefas relacionadas a localiza o artigo-gap que norteia pesquisa sobre TDAH, brincar livre e natureza em pré-escolares. O orquestrador pode encaminhar-lhe o trabalho na etapa correspondente; a resposta retorna ao fluxo principal para integração e conferência de fontes, critérios e permissões. A ficha documenta a função, mas não comprova que o perfil tenha sido chamado ou executado a tarefa.
\end{funcao}

\begin{arquitetura}
\textbf{Modo de execução declarado:} subagente · \textbf{Permissões:} edit: deny; bash: deny
\end{arquitetura}

\clearpage
\section{Categoria: Auditoria e qualidade (13 perfis)}

\begin{descricao}
Reúne perfis de inspeção e revisão. Seus achados podem informar verificações de qualidade e rigor, mas a existência de um auditor não equivale a aprovação independente (Módulos 4 e 6).
\end{descricao}

\elemento{Auditor (v1.0.0)}{\metainfo{Categoria}{Auditoria e qualidade} \hfill \metainfo{Tipo}{especialista} \hfill \metainfo{Identificador}{auditor} \hfill \metainfo{Ficha}{auditor.md}}

\begin{descricao}
Agente auditor do ecossistema — sanidade, integração e métricas de qualidade.
\end{descricao}

\begin{funcao}
No Core, este perfil apoia tarefas relacionadas a agente auditor do ecossistema — sanidade, integração e métricas de qualidade.. O orquestrador pode encaminhar-lhe o trabalho na etapa correspondente; a resposta retorna ao fluxo principal para integração e conferência de fontes, critérios e permissões. A ficha documenta a função, mas não comprova que o perfil tenha sido chamado ou executado a tarefa.
\end{funcao}

\begin{arquitetura}
\textbf{Modo de execução declarado:} subagente · \textbf{Permissões:} edit: deny; bash: deny
\end{arquitetura}

\elemento{Preparação de livros para ajuste fino (v1.0.0)}{\metainfo{Categoria}{Auditoria e qualidade} \hfill \metainfo{Tipo}{especialista} \hfill \metainfo{Identificador}{book-finetuning} \hfill \metainfo{Ficha}{book-finetuning.md}}

\begin{descricao}
Prepara dados de livros para ajuste fino, com verificação de licença, origem, grupos e separação de conjuntos. A preparação dos dados não executa treinamento de modelo.
\end{descricao}

\begin{funcao}
No Core, este perfil apoia tarefas relacionadas a prepara dados de livros para ajuste fino, com verificação de licença, origem, grupos e separação de conjuntos.. O orquestrador pode encaminhar-lhe o trabalho na etapa correspondente; a resposta retorna ao fluxo principal para integração e conferência de fontes, critérios e permissões. A ficha documenta a função, mas não comprova que o perfil tenha sido chamado ou executado a tarefa.
\end{funcao}

\begin{arquitetura}
\textbf{Modo de execução declarado:} subagente · \textbf{Permissões:} read: allow; glob: allow; grep: allow; bash: \{'*': 'deny', 'python3 -m pytest tests/test\_r669\_requested\_agents.py -q': 'allow', 'python3 -m pytest tests/test\_r658\_finetuning\_data.py -q': 'allow'\}; edit: deny; webfetch: deny; task: deny
\end{arquitetura}

\elemento{Consulta de livros por MCP (v1.0.0)}{\metainfo{Categoria}{Auditoria e qualidade} \hfill \metainfo{Tipo}{especialista} \hfill \metainfo{Identificador}{book-mcp} \hfill \metainfo{Ficha}{book-mcp.md}}

\begin{descricao}
Consulta livros locais pela biblioteca do Core e sua interface MCP. Preserva origem e hashes das passagens e mantém a leitura dos documentos separada da execução de instruções neles contidas.
\end{descricao}

\begin{funcao}
No Core, este perfil apoia tarefas relacionadas a consulta livros locais pela biblioteca do Core e sua interface MCP.. O orquestrador pode encaminhar-lhe o trabalho na etapa correspondente; a resposta retorna ao fluxo principal para integração e conferência de fontes, critérios e permissões. A ficha documenta a função, mas não comprova que o perfil tenha sido chamado ou executado a tarefa.
\end{funcao}

\begin{arquitetura}
\textbf{Modo de execução declarado:} subagente · \textbf{Permissões:} read: allow; glob: allow; grep: allow; bash: \{'*': 'deny', 'python3 -m pytest tests/test\_r669\_requested\_agents.py -q': 'allow', 'python3 -m pytest tests/test\_r657\_book\_mcp.py -q': 'allow'\}; edit: deny; webfetch: deny; task: deny
\end{arquitetura}

\elemento{Auditor de evidências do Gemini Notebook (v1.0.0)}{\metainfo{Categoria}{Auditoria e qualidade} \hfill \metainfo{Tipo}{especialista} \hfill \metainfo{Identificador}{gemini-notebook-audit} \hfill \metainfo{Ficha}{gemini-notebook-audit.md}}

\begin{descricao}
Audita estados, evidências e proveniência do Gemini Notebook pelo orquestrador central. Confere autenticação, transporte, resultado da operação e arquivos, preservando pendências, resultados parciais e falhas sem promover confiança científica.
\end{descricao}

\begin{funcao}
No Core, este perfil apoia tarefas relacionadas a audita estados, evidências e proveniência do Gemini Notebook pelo orquestrador central.. O orquestrador pode encaminhar-lhe o trabalho na etapa correspondente; a resposta retorna ao fluxo principal para integração e conferência de fontes, critérios e permissões. A ficha documenta a função, mas não comprova que o perfil tenha sido chamado ou executado a tarefa.
\end{funcao}

\begin{arquitetura}
\textbf{Modo de execução declarado:} subagente · \textbf{Permissões:} read: allow; glob: allow; grep: allow; bash: \{'*': 'deny', 'python3 -m pytest tests/test\_r676\_gemini\_notebook\_specialists.py -q': 'allow', 'python3 -m pytest tests/test\_r672\_gemini\_notebook\_transport.py -q': 'allow', 'python3 -m pytest tests/test\_r672\_gemini\_notebook\_session.py -q': 'allow', 'python3 -m pytest tests/test\_r673\_gemini\_notebook\_catalog.py -q': 'allow', 'python3 -m pytest tests/test\_r674\_gemini\_notebook\_orchestration.py -q': 'allow'\}; edit: deny; webfetch: deny; task: deny
\end{arquitetura}

\elemento{Especialista em transporte do Gemini Notebook (v1.0.0)}{\metainfo{Categoria}{Auditoria e qualidade} \hfill \metainfo{Tipo}{especialista} \hfill \metainfo{Identificador}{gemini-notebook-transport} \hfill \metainfo{Ficha}{gemini-notebook-transport.md}}

\begin{descricao}
Verifica os executores oficiais CLI e MCP, as sessões persistentes e a associação de tarefas à sessão original. Examina limites, encerramento e downloads com arquivos reais; operações da conta seguem pelas entradas centrais e exigem autorização correspondente.
\end{descricao}

\begin{funcao}
No Core, este perfil apoia tarefas relacionadas a verifica os executores oficiais CLI e MCP, as sessões persistentes e a associação de tarefas à sessão original.. O orquestrador pode encaminhar-lhe o trabalho na etapa correspondente; a resposta retorna ao fluxo principal para integração e conferência de fontes, critérios e permissões. A ficha documenta a função, mas não comprova que o perfil tenha sido chamado ou executado a tarefa.
\end{funcao}

\begin{arquitetura}
\textbf{Modo de execução declarado:} subagente · \textbf{Permissões:} read: allow; glob: allow; grep: allow; bash: \{'*': 'deny', 'python3 -m pytest tests/test\_r676\_gemini\_notebook\_specialists.py -q': 'allow', 'python3 -m pytest tests/test\_r672\_gemini\_notebook\_transport.py -q': 'allow', 'python3 -m pytest tests/test\_r672\_gemini\_notebook\_session.py -q': 'allow', 'python3 -m pytest tests/test\_r673\_gemini\_notebook\_catalog.py -q': 'allow', 'python3 -m pytest tests/test\_r674\_gemini\_notebook\_orchestration.py -q': 'allow'\}; edit: allow; webfetch: deny; task: deny
\end{arquitetura}

\elemento{Auditor das fontes e contratos do Gemini Notebook (v1.0.0)}{\metainfo{Categoria}{Auditoria e qualidade} \hfill \metainfo{Tipo}{especialista} \hfill \metainfo{Identificador}{gemini-notebook-upstream} \hfill \metainfo{Ficha}{gemini-notebook-upstream.md}}

\begin{descricao}
Compara o fork solicitado, a versão instalada, os esquemas e as definições oficiais dos pipelines. Registra divergências e hashes das fontes públicas sem trocar perfis, atualizar pacotes ou executar operações da conta durante a descoberta.
\end{descricao}

\begin{funcao}
No Core, este perfil apoia tarefas relacionadas a compara o fork solicitado, a versão instalada, os esquemas e as definições oficiais dos pipelines.. O orquestrador pode encaminhar-lhe o trabalho na etapa correspondente; a resposta retorna ao fluxo principal para integração e conferência de fontes, critérios e permissões. A ficha documenta a função, mas não comprova que o perfil tenha sido chamado ou executado a tarefa.
\end{funcao}

\begin{arquitetura}
\textbf{Modo de execução declarado:} subagente · \textbf{Permissões:} read: allow; glob: allow; grep: allow; bash: \{'*': 'deny', 'python3 -m pytest tests/test\_r676\_gemini\_notebook\_specialists.py -q': 'allow', 'python3 -m pytest tests/test\_r672\_gemini\_notebook\_transport.py -q': 'allow', 'python3 -m pytest tests/test\_r672\_gemini\_notebook\_session.py -q': 'allow', 'python3 -m pytest tests/test\_r673\_gemini\_notebook\_catalog.py -q': 'allow', 'python3 -m pytest tests/test\_r674\_gemini\_notebook\_orchestration.py -q': 'allow'\}; edit: deny; webfetch: allow; task: deny
\end{arquitetura}

\elemento{Integração de hooks (v1.0.0)}{\metainfo{Categoria}{Auditoria e qualidade} \hfill \metainfo{Tipo}{especialista} \hfill \metainfo{Identificador}{hooks-integration} \hfill \metainfo{Ficha}{hooks-integration.md}}

\begin{descricao}
Verifica integração de hooks e contratos do SDK, com casos de falha e sucesso. A ativação de um evento não demonstra a conclusão da tarefa; a entrega exige resultado e evidência correspondente.
\end{descricao}

\begin{funcao}
No Core, este perfil apoia tarefas relacionadas a verifica integração de hooks e contratos do SDK, com casos de falha e sucesso.. O orquestrador pode encaminhar-lhe o trabalho na etapa correspondente; a resposta retorna ao fluxo principal para integração e conferência de fontes, critérios e permissões. A ficha documenta a função, mas não comprova que o perfil tenha sido chamado ou executado a tarefa.
\end{funcao}

\begin{arquitetura}
\textbf{Modo de execução declarado:} subagente · \textbf{Permissões:} read: allow; glob: allow; grep: allow; bash: \{'*': 'deny', 'python3 -m pytest tests/test\_r669\_requested\_agents.py -q': 'allow', 'python3 -m pytest tests/test\_r659\_sdk\_hooks\_contract.py -q': 'allow'\}; edit: allow; webfetch: deny; task: deny
\end{arquitetura}

\elemento{Arquitetura da biblioteca (v1.0.0)}{\metainfo{Categoria}{Auditoria e qualidade} \hfill \metainfo{Tipo}{especialista} \hfill \metainfo{Identificador}{library-architecture} \hfill \metainfo{Ficha}{library-architecture.md}}

\begin{descricao}
Verifica estrutura, índice, cobertura e recuperação da biblioteca. Encaminha alterações e tarefas ao orquestrador central e registra ausências sem tratá-las como capacidades já executadas.
\end{descricao}

\begin{funcao}
No Core, este perfil apoia tarefas relacionadas a verifica estrutura, índice, cobertura e recuperação da biblioteca.. O orquestrador pode encaminhar-lhe o trabalho na etapa correspondente; a resposta retorna ao fluxo principal para integração e conferência de fontes, critérios e permissões. A ficha documenta a função, mas não comprova que o perfil tenha sido chamado ou executado a tarefa.
\end{funcao}

\begin{arquitetura}
\textbf{Modo de execução declarado:} subagente · \textbf{Permissões:} read: allow; glob: allow; grep: allow; bash: \{'*': 'deny', 'python3 -m pytest tests/test\_r669\_requested\_agents.py -q': 'allow', 'python3 -m pytest tests/test\_r657\_book\_library.py -q': 'allow'\}; edit: deny; webfetch: deny; task: deny
\end{arquitetura}

\elemento{Execução externa de MiroFish e Hermes (v1.0.0)}{\metainfo{Categoria}{Auditoria e qualidade} \hfill \metainfo{Tipo}{especialista} \hfill \metainfo{Identificador}{live-mirofish-hermes} \hfill \metainfo{Ficha}{live-mirofish-hermes.md}}

\begin{descricao}
Executa Hermes e MiroFish/OASIS externos pelo método central scientific\_runtime\_run. Exige processo externo, resposta do modelo com tokens e artefato do runtime. Registra commits, hashes e limites; a população simulada permanece sintética.
\end{descricao}

\begin{funcao}
No Core, este perfil apoia tarefas relacionadas a executa Hermes e MiroFish/OASIS externos pelo método central scientific\_runtime\_run.. O orquestrador pode encaminhar-lhe o trabalho na etapa correspondente; a resposta retorna ao fluxo principal para integração e conferência de fontes, critérios e permissões. A ficha documenta a função, mas não comprova que o perfil tenha sido chamado ou executado a tarefa.
\end{funcao}

\begin{arquitetura}
\textbf{Modo de execução declarado:} subagente · \textbf{Permissões:} read: allow; glob: allow; grep: allow; bash: \{'*': 'deny', 'python3 -m pytest tests/test\_r669\_requested\_agents.py -q': 'allow', 'python3 -m pytest tests/test\_r667\_live\_scientific\_runtime.py -q': 'allow'\}; edit: deny; webfetch: deny; task: deny
\end{arquitetura}

\elemento{Integração de MCP e CLI (v1.0.0)}{\metainfo{Categoria}{Auditoria e qualidade} \hfill \metainfo{Tipo}{especialista} \hfill \metainfo{Identificador}{mcp-cli-integration} \hfill \metainfo{Ficha}{mcp-cli-integration.md}}

\begin{descricao}
Verifica esquemas, argumentos, inicializadores e ferramentas de MCP e CLI. Diferencia descoberta, instalação, autenticação, processo iniciado e operação concluída, preservando os erros observados.
\end{descricao}

\begin{funcao}
No Core, este perfil apoia tarefas relacionadas a verifica esquemas, argumentos, inicializadores e ferramentas de MCP e CLI.. O orquestrador pode encaminhar-lhe o trabalho na etapa correspondente; a resposta retorna ao fluxo principal para integração e conferência de fontes, critérios e permissões. A ficha documenta a função, mas não comprova que o perfil tenha sido chamado ou executado a tarefa.
\end{funcao}

\begin{arquitetura}
\textbf{Modo de execução declarado:} subagente · \textbf{Permissões:} read: allow; glob: allow; grep: allow; bash: \{'*': 'deny', 'python3 -m pytest tests/test\_r669\_requested\_agents.py -q': 'allow', 'python3 -m pytest tests/test\_r660\_mcp\_validation.py -q': 'allow'\}; edit: allow; webfetch: deny; task: deny
\end{arquitetura}

\elemento{Mira validator (v1.0.0)}{\metainfo{Categoria}{Auditoria e qualidade} \hfill \metainfo{Tipo}{mira-agent} \hfill \metainfo{Identificador}{mira-validator} \hfill \metainfo{Ficha}{mira-validator.md}}

\begin{descricao}
Agente auditor do MIRA que verifica coerência, completude, responsividade e conformidade da apresentação final.
\end{descricao}

\begin{funcao}
No Core, este perfil apoia tarefas relacionadas a agente auditor do MIRA que verifica coerência, completude, responsividade e conformidade da apresentação final.. O orquestrador pode encaminhar-lhe o trabalho na etapa correspondente; a resposta retorna ao fluxo principal para integração e conferência de fontes, critérios e permissões. A ficha documenta a função, mas não comprova que o perfil tenha sido chamado ou executado a tarefa.
\end{funcao}

\begin{arquitetura}
\textbf{Modo de execução declarado:} subagente · \textbf{Permissões:} edit: deny; bash: deny
\end{arquitetura}

\elemento{Auditor de capacidades científicas (v1.0.0)}{\metainfo{Categoria}{Auditoria e qualidade} \hfill \metainfo{Tipo}{especialista} \hfill \metainfo{Identificador}{scientific-capabilities-audit} \hfill \metainfo{Ficha}{scientific-capabilities-audit.md}}

\begin{descricao}
Audita capacidades científicas e lacunas pelo orquestrador central. Verifica fontes, métodos, reprodução e revisão computacional, preservando o escopo da evidência e a necessidade de avaliação externa.
\end{descricao}

\begin{funcao}
No Core, este perfil apoia tarefas relacionadas a audita capacidades científicas e lacunas pelo orquestrador central.. O orquestrador pode encaminhar-lhe o trabalho na etapa correspondente; a resposta retorna ao fluxo principal para integração e conferência de fontes, critérios e permissões. A ficha documenta a função, mas não comprova que o perfil tenha sido chamado ou executado a tarefa.
\end{funcao}

\begin{arquitetura}
\textbf{Modo de execução declarado:} subagente · \textbf{Permissões:} read: allow; glob: allow; grep: allow; bash: \{'*': 'deny', 'python3 -m pytest tests/test\_r669\_requested\_agents.py -q': 'allow', 'python3 -m pytest tests/test\_r663\_scientific\_provenance.py -q': 'allow'\}; edit: deny; webfetch: deny; task: deny
\end{arquitetura}

\elemento{Auditor de simulações e teoria dos jogos (v1.0.0)}{\metainfo{Categoria}{Auditoria e qualidade} \hfill \metainfo{Tipo}{especialista} \hfill \metainfo{Identificador}{simulation-game-audit} \hfill \metainfo{Ficha}{simulation-game-audit.md}}

\begin{descricao}
Audita simulações e teoria dos jogos pelo orquestrador central. Confere matrizes, equilíbrios e cenários, distinguindo execução de um cenário sintético de observação social ou previsão validada.
\end{descricao}

\begin{funcao}
No Core, este perfil apoia tarefas relacionadas a audita simulações e teoria dos jogos pelo orquestrador central.. O orquestrador pode encaminhar-lhe o trabalho na etapa correspondente; a resposta retorna ao fluxo principal para integração e conferência de fontes, critérios e permissões. A ficha documenta a função, mas não comprova que o perfil tenha sido chamado ou executado a tarefa.
\end{funcao}

\begin{arquitetura}
\textbf{Modo de execução declarado:} subagente · \textbf{Permissões:} read: allow; glob: allow; grep: allow; bash: \{'*': 'deny', 'python3 -m pytest tests/test\_r669\_requested\_agents.py -q': 'allow', 'python3 -m pytest tests/test\_r665\_evolutionary\_sequencing.py -q': 'allow'\}; edit: deny; webfetch: deny; task: deny
\end{arquitetura}

\clearpage
\section{Categoria: Apoio jurídico (4 perfis)}

\begin{descricao}
Agrupa perfis de apoio a documentos, pesquisa e organização jurídica. No Core, podem estruturar informações e apontar questões para revisão humana; não substituem julgamento profissional.
\end{descricao}

\elemento{Especialista em síntese de documentos jurídicos (v1.0.0)}{\metainfo{Categoria}{Apoio jurídico} \hfill \metainfo{Tipo}{specialist-agent} \hfill \metainfo{Identificador}{auxjuris\_document\_summarizer} \hfill \metainfo{Ficha}{auxjuris\_document\_summarizer.md}}

\begin{descricao}
Resume documentos jurídicos, destacando informações relevantes para facilitar a leitura e a análise posterior.
\end{descricao}

\begin{funcao}
No Core, este perfil apoia tarefas relacionadas a resume documentos jurídicos, destacando informações relevantes para facilitar a leitura e a análise posterior.. O orquestrador pode encaminhar-lhe o trabalho na etapa correspondente; a resposta retorna ao fluxo principal para integração e conferência de fontes, critérios e permissões. A ficha documenta a função, mas não comprova que o perfil tenha sido chamado ou executado a tarefa.
\end{funcao}

\begin{arquitetura}
\textbf{Modo de execução declarado:} subagente · \textbf{Permissões:} edit: deny; bash: deny
\end{arquitetura}

\elemento{Auxjuris email drafter (v1.0.0)}{\metainfo{Categoria}{Apoio jurídico} \hfill \metainfo{Tipo}{specialist-agent} \hfill \metainfo{Identificador}{auxjuris\_email\_drafter} \hfill \metainfo{Ficha}{auxjuris\_email\_drafter.md}}

\begin{descricao}
Agente de redação jurídica inspirado no AUXJURIS. Gera minutas de e-mails e comunicações profissionais para clientes, advogados, juízes e órgãos públicos com tom formal, claro e juridicamente apropriado.
\end{descricao}

\begin{funcao}
No Core, este perfil apoia tarefas relacionadas a agente de redação jurídica inspirado no AUXJURIS.. O orquestrador pode encaminhar-lhe o trabalho na etapa correspondente; a resposta retorna ao fluxo principal para integração e conferência de fontes, critérios e permissões. A ficha documenta a função, mas não comprova que o perfil tenha sido chamado ou executado a tarefa.
\end{funcao}

\begin{arquitetura}
\textbf{Modo de execução declarado:} subagente · \textbf{Permissões:} edit: deny; bash: deny
\end{arquitetura}

\elemento{Auxjuris jurídico assistant (v1.0.0)}{\metainfo{Categoria}{Apoio jurídico} \hfill \metainfo{Tipo}{specialist-agent} \hfill \metainfo{Identificador}{auxjuris\_legal\_assistant} \hfill \metainfo{Ficha}{auxjuris\_legal\_assistant.md}}

\begin{descricao}
Agente jurídico geral inspirado no AUXJURIS. Atua em análise de casos, subsunção legal, pesquisa de precedentes e orientação jurídica inicial com base em legislação brasileira, base de conhecimento jurídica e dados do Datajud.
\end{descricao}

\begin{funcao}
No Core, este perfil apoia tarefas relacionadas a agente jurídico geral inspirado no AUXJURIS.. O orquestrador pode encaminhar-lhe o trabalho na etapa correspondente; a resposta retorna ao fluxo principal para integração e conferência de fontes, critérios e permissões. A ficha documenta a função, mas não comprova que o perfil tenha sido chamado ou executado a tarefa.
\end{funcao}

\begin{arquitetura}
\textbf{Modo de execução declarado:} subagente · \textbf{Permissões:} edit: deny; bash: deny
\end{arquitetura}

\elemento{Auxjuris jurídico pesquisa (v1.0.0)}{\metainfo{Categoria}{Apoio jurídico} \hfill \metainfo{Tipo}{specialist-agent} \hfill \metainfo{Identificador}{auxjuris\_legal\_research} \hfill \metainfo{Ficha}{auxjuris\_legal\_research.md}}

\begin{descricao}
Agente de pesquisa jurídica inspirado no AUXJURIS. Analisa consultas, busca legislação, doutrina e jurisprudência relevante na base de conhecimento jurídica e nos dados processuais do Datajud, retornando trechos contextualizados com fontes.
\end{descricao}

\begin{funcao}
No Core, este perfil apoia tarefas relacionadas a agente de pesquisa jurídica inspirado no AUXJURIS.. O orquestrador pode encaminhar-lhe o trabalho na etapa correspondente; a resposta retorna ao fluxo principal para integração e conferência de fontes, critérios e permissões. A ficha documenta a função, mas não comprova que o perfil tenha sido chamado ou executado a tarefa.
\end{funcao}

\begin{arquitetura}
\textbf{Modo de execução declarado:} subagente · \textbf{Permissões:} edit: deny; bash: deny
\end{arquitetura}

\clearpage
\section{Categoria: Dados e proveniência (3 perfis)}

\begin{descricao}
Reúne perfis voltados a dados. No Core, podem apoiar preparação, análise e registro de proveniência antes que os resultados alimentem pesquisa ou decisão (Módulo 6).
\end{descricao}

\elemento{Especialista em Haystack RAG (v1.0.0)}{\metainfo{Categoria}{Dados e proveniência} \hfill \metainfo{Tipo}{integration} \hfill \metainfo{Identificador}{Haystack RAG Specialist} \hfill \metainfo{Ficha}{haystack-rag.md}}

\begin{descricao}
Integração do framework Haystack (deepset-ai/haystack, Apache-2.0) ao ecossistema: verificação de disponibilidade/versão (haystack-ai via PyPI), construção de pipeline RAG mínimo (InMemoryDocumentStore + SentenceTransformers Document/Text Embedder + InMemoryEmbeddingRetriever) e health check tolerante — sem exigir instalação para a suíte passar.
\end{descricao}

\begin{funcao}
No Core, este perfil apoia tarefas relacionadas a integração do framework Haystack (deepset-ai/haystack, Apache-2.0) ao ecossistema: verificação de disponibilidade/versão (haystack-ai via PyPI), construção de pipeline RAG mínimo (InMemoryDocumentStore + SentenceTransformers Document/Text Embedder…. O orquestrador pode encaminhar-lhe o trabalho na etapa correspondente; a resposta retorna ao fluxo principal para integração e conferência de fontes, critérios e permissões. A ficha documenta a função, mas não comprova que o perfil tenha sido chamado ou executado a tarefa.
\end{funcao}

\begin{arquitetura}
\textbf{Modo de execução declarado:} subagente
\end{arquitetura}

\elemento{Mira chart (v1.0.0)}{\metainfo{Categoria}{Dados e proveniência} \hfill \metainfo{Tipo}{mira-agent} \hfill \metainfo{Identificador}{mira-chart} \hfill \metainfo{Ficha}{mira-chart.md}}

\begin{descricao}
Gera gráficos a partir de dados para apresentações MIRA
\end{descricao}

\begin{funcao}
No Core, este perfil apoia tarefas relacionadas a gera gráficos a partir de dados para apresentações MIRA. O orquestrador pode encaminhar-lhe o trabalho na etapa correspondente; a resposta retorna ao fluxo principal para integração e conferência de fontes, critérios e permissões. A ficha documenta a função, mas não comprova que o perfil tenha sido chamado ou executado a tarefa.
\end{funcao}

\begin{arquitetura}
\textbf{Modo de execução declarado:} subagente · \textbf{Permissões:} edit: deny; bash: deny
\end{arquitetura}

\elemento{Mira chart race (v1.0.0)}{\metainfo{Categoria}{Dados e proveniência} \hfill \metainfo{Tipo}{mira-agent} \hfill \metainfo{Identificador}{mira-chart-race} \hfill \metainfo{Ficha}{mira-chart-race.md}}

\begin{descricao}
Agente especializado em chart races e comparações temporais animadas para apresentações MIRA.
\end{descricao}

\begin{funcao}
No Core, este perfil apoia tarefas relacionadas a agente especializado em chart races e comparações temporais animadas para apresentações MIRA.. O orquestrador pode encaminhar-lhe o trabalho na etapa correspondente; a resposta retorna ao fluxo principal para integração e conferência de fontes, critérios e permissões. A ficha documenta a função, mas não comprova que o perfil tenha sido chamado ou executado a tarefa.
\end{funcao}

\begin{arquitetura}
\textbf{Modo de execução declarado:} subagente · \textbf{Permissões:} edit: deny; bash: deny
\end{arquitetura}

\clearpage
\section{Categoria: Inferência e modelagem (1 perfis)}

\begin{descricao}
Contém perfis de inferência e modelagem. Suas contribuições se relacionam aos métodos de raciocínio do Core e precisam expor premissas e limites (Módulo 5).
\end{descricao}

\elemento{Colibri agente (v1.0.0)}{\metainfo{Categoria}{Inferência e modelagem} \hfill \metainfo{Tipo}{especialista} \hfill \metainfo{Identificador}{colibri-agent} \hfill \metainfo{Ficha}{colibri-agent.md}}

\begin{descricao}
ID: colibri-agent Tipo: Agente de inferência on-device (A2A registrado no Blackboard) MCP Server: integrations/colibri/colibri\_mcp\_server.py Bridge: integrations/colibri/bridge.py::ColibriBridge Repositório: https://github.com/MarceloClaro/colibri (fork de JustVugg/colibri) Confiança inicial: 0.90 (Trust Engine) --- Agente especializado em execução local de modelos MoE via runtime Colibri — motor de inferência em C puro, zero dependências. Suporta dois modelos: ; Modelo ; Params ; Experts ; RAM ; Status ; ; ---; ---; ---; ---; ---; ; OLMoE-1B-7B ; 1B dense + 7B sparse ; 64 (8 ativos) ; \~{}8 GB ; Já convertido e validado ; ; GLM-5.2 ; 744B MoE ; 19.456 ; \~{}25 GB mínimo ; Requer hardware com GPU ; ; Capacidade ; Des
\end{descricao}

\begin{funcao}
No Core, este perfil apoia tarefas relacionadas a iD: colibri-agent Tipo: Agente de inferência on-device (A2A registrado no Blackboard) MCP Server: integrations/colibri/colibri\_mcp\_server.py Bridge: integrations/colibri/bridge.py::ColibriBridge Repositório: https://github.com/MarceloClaro/colibri (fork de JustVugg/colibri) Confiança inicial: 0.90…. O orquestrador pode encaminhar-lhe o trabalho na etapa correspondente; a resposta retorna ao fluxo principal para integração e conferência de fontes, critérios e permissões. A ficha documenta a função, mas não comprova que o perfil tenha sido chamado ou executado a tarefa.
\end{funcao}

\begin{arquitetura}
\textbf{Modo de execução declarado:} subagente · \textbf{Permissões:} edit: deny; bash: deny
\end{arquitetura}

\clearpage
\section{Categoria: Engenharia (16 perfis)}

\begin{descricao}
Agrupa perfis que apoiam análise, implementação e manutenção de software. No Core, o trabalho deve ser ligado a especificações e verificações quando aplicável (Módulos 1 e 4).
\end{descricao}

\elemento{Abnt latex modular (v1.0.0)}{\metainfo{Categoria}{Engenharia} \hfill \metainfo{Tipo}{mira-agent} \hfill \metainfo{Identificador}{abnt-latex-modular} \hfill \metainfo{Ficha}{abnt-latex-modular.md}}

\begin{descricao}
Constrói artigos em LaTeX modular ABNT com abnTeX2/BibTeX e compilação verificada. Capacidades descritas: Capacidade especializada no ciclo pdflatex-bibtex com estilo ABNT autor-data
\end{descricao}

\begin{funcao}
No Core, este perfil apoia tarefas relacionadas a constrói artigos em LaTeX modular ABNT com abnTeX2/BibTeX e compilação verificada.. O orquestrador pode encaminhar-lhe o trabalho na etapa correspondente; a resposta retorna ao fluxo principal para integração e conferência de fontes, critérios e permissões. A ficha documenta a função, mas não comprova que o perfil tenha sido chamado ou executado a tarefa.
\end{funcao}

\begin{arquitetura}
\textbf{Modo de execução declarado:} subagente · \textbf{Permissões:} edit: deny; bash: deny
\end{arquitetura}

\elemento{Programador de (v1.0.0)}{\metainfo{Categoria}{Engenharia} \hfill \metainfo{Tipo}{especialista} \hfill \metainfo{Identificador}{coder} \hfill \metainfo{Ficha}{coder.md}}

\begin{descricao}
Agente de implementação de código Python, refatoração e depuração.
\end{descricao}

\begin{funcao}
No Core, este perfil apoia tarefas relacionadas a agente de implementação de código Python, refatoração e depuração.. O orquestrador pode encaminhar-lhe o trabalho na etapa correspondente; a resposta retorna ao fluxo principal para integração e conferência de fontes, critérios e permissões. A ficha documenta a função, mas não comprova que o perfil tenha sido chamado ou executado a tarefa.
\end{funcao}

\begin{arquitetura}
\textbf{Modo de execução declarado:} subagente · \textbf{Permissões:} edit: allow; bash: allow
\end{arquitetura}

\elemento{Docx abnt converter (v1.0.0)}{\metainfo{Categoria}{Engenharia} \hfill \metainfo{Tipo}{mira-agent} \hfill \metainfo{Identificador}{docx-abnt-converter} \hfill \metainfo{Ficha}{docx-abnt-converter.md}}

\begin{descricao}
Converte artigos LaTeX em DOCX formatado ABNT com citações reconstruídas do .aux
\end{descricao}

\begin{funcao}
No Core, este perfil apoia tarefas relacionadas a converte artigos LaTeX em DOCX formatado ABNT com citações reconstruídas do .aux. O orquestrador pode encaminhar-lhe o trabalho na etapa correspondente; a resposta retorna ao fluxo principal para integração e conferência de fontes, critérios e permissões. A ficha documenta a função, mas não comprova que o perfil tenha sido chamado ou executado a tarefa.
\end{funcao}

\begin{arquitetura}
\textbf{Modo de execução declarado:} subagente · \textbf{Permissões:} edit: deny; bash: deny
\end{arquitetura}

\elemento{Agente de inferência local LiteRT-LM (v1.0.0)}{\metainfo{Categoria}{Engenharia} \hfill \metainfo{Tipo}{especialista} \hfill \metainfo{Identificador}{litert-lm-agent} \hfill \metainfo{Ficha}{litert-lm-agent.md}}

\begin{descricao}
Opera modelos de linguagem no próprio dispositivo por meio do LiteRT-LM: lista e importa modelos, executa prompts, mantém conversas interativas, consulta metadados e pode iniciar um servidor compatível com a API da OpenAI.
\end{descricao}

\begin{funcao}
No Core, este perfil apoia tarefas relacionadas a opera modelos de linguagem no próprio dispositivo por meio do LiteRT-LM: lista e importa modelos, executa prompts, mantém conversas interativas, consulta metadados e pode…. O orquestrador pode encaminhar-lhe o trabalho na etapa correspondente; a resposta retorna ao fluxo principal para integração e conferência de fontes, critérios e permissões. A ficha documenta a função, mas não comprova que o perfil tenha sido chamado ou executado a tarefa.
\end{funcao}

\begin{arquitetura}
\textbf{Modo de execução declarado:} subagente · \textbf{Permissões:} edit: deny; bash: deny
\end{arquitetura}

\elemento{Mira 3d (v1.0.0)}{\metainfo{Categoria}{Engenharia} \hfill \metainfo{Tipo}{mira-agent} \hfill \metainfo{Identificador}{mira-3d} \hfill \metainfo{Ficha}{mira-3d.md}}

\begin{descricao}
Agente dedicado a incorporar ou orientar o uso de elementos 3D em slides e metáforas visuais do MIRA.
\end{descricao}

\begin{funcao}
No Core, este perfil apoia tarefas relacionadas a agente dedicado a incorporar ou orientar o uso de elementos 3D em slides e metáforas visuais do MIRA.. O orquestrador pode encaminhar-lhe o trabalho na etapa correspondente; a resposta retorna ao fluxo principal para integração e conferência de fontes, critérios e permissões. A ficha documenta a função, mas não comprova que o perfil tenha sido chamado ou executado a tarefa.
\end{funcao}

\begin{arquitetura}
\textbf{Modo de execução declarado:} subagente · \textbf{Permissões:} edit: deny; bash: deny
\end{arquitetura}

\elemento{Mira animated metaphor (v1.0.0)}{\metainfo{Categoria}{Engenharia} \hfill \metainfo{Tipo}{mira-agent} \hfill \metainfo{Identificador}{mira-animated-metaphor} \hfill \metainfo{Ficha}{mira-animated-metaphor.md}}

\begin{descricao}
Agente especializado em destilar a alma do conceito e convertê-la numa metáfora visual animada e didática.
\end{descricao}

\begin{funcao}
No Core, este perfil apoia tarefas relacionadas a agente especializado em destilar a alma do conceito e convertê-la numa metáfora visual animada e didática.. O orquestrador pode encaminhar-lhe o trabalho na etapa correspondente; a resposta retorna ao fluxo principal para integração e conferência de fontes, critérios e permissões. A ficha documenta a função, mas não comprova que o perfil tenha sido chamado ou executado a tarefa.
\end{funcao}

\begin{arquitetura}
\textbf{Modo de execução declarado:} subagente · \textbf{Permissões:} edit: deny; bash: deny
\end{arquitetura}

\elemento{Mira animator (v1.0.0)}{\metainfo{Categoria}{Engenharia} \hfill \metainfo{Tipo}{mira-agent} \hfill \metainfo{Identificador}{mira-animator} \hfill \metainfo{Ficha}{mira-animator.md}}

\begin{descricao}
Agente que usa o MiraEngine para transformar conceitos em cards HTML animados com Regra Zero e movimento perpétuo.
\end{descricao}

\begin{funcao}
No Core, este perfil apoia tarefas relacionadas a agente que usa o MiraEngine para transformar conceitos em cards HTML animados com Regra Zero e movimento perpétuo.. O orquestrador pode encaminhar-lhe o trabalho na etapa correspondente; a resposta retorna ao fluxo principal para integração e conferência de fontes, critérios e permissões. A ficha documenta a função, mas não comprova que o perfil tenha sido chamado ou executado a tarefa.
\end{funcao}

\begin{arquitetura}
\textbf{Modo de execução declarado:} subagente · \textbf{Permissões:} edit: deny; bash: deny
\end{arquitetura}

\elemento{Mira construtor (v1.0.0)}{\metainfo{Categoria}{Engenharia} \hfill \metainfo{Tipo}{mira-agent} \hfill \metainfo{Identificador}{mira-builder} \hfill \metainfo{Ficha}{mira-builder.md}}

\begin{descricao}
Monta o deck MIRA em cards e seções navegáveis. Capacidades descritas: Monta o deck mira em cards
\end{descricao}

\begin{funcao}
No Core, este perfil apoia tarefas relacionadas a monta o deck MIRA em cards e seções navegáveis.. O orquestrador pode encaminhar-lhe o trabalho na etapa correspondente; a resposta retorna ao fluxo principal para integração e conferência de fontes, critérios e permissões. A ficha documenta a função, mas não comprova que o perfil tenha sido chamado ou executado a tarefa.
\end{funcao}

\begin{arquitetura}
\textbf{Modo de execução declarado:} subagente · \textbf{Permissões:} edit: deny; bash: deny
\end{arquitetura}

\elemento{Mira imagem (v1.0.0)}{\metainfo{Categoria}{Engenharia} \hfill \metainfo{Tipo}{mira-agent} \hfill \metainfo{Identificador}{mira-image} \hfill \metainfo{Ficha}{mira-image.md}}

\begin{descricao}
Agente que posiciona imagens existentes do usuário dentro do deck, preservando composição, legibilidade e coerência visual.
\end{descricao}

\begin{funcao}
No Core, este perfil apoia tarefas relacionadas a agente que posiciona imagens existentes do usuário dentro do deck, preservando composição, legibilidade e coerência visual.. O orquestrador pode encaminhar-lhe o trabalho na etapa correspondente; a resposta retorna ao fluxo principal para integração e conferência de fontes, critérios e permissões. A ficha documenta a função, mas não comprova que o perfil tenha sido chamado ou executado a tarefa.
\end{funcao}

\begin{arquitetura}
\textbf{Modo de execução declarado:} subagente · \textbf{Permissões:} edit: deny; bash: deny
\end{arquitetura}

\elemento{Mira imagem template (v1.0.0)}{\metainfo{Categoria}{Engenharia} \hfill \metainfo{Tipo}{mira-agent} \hfill \metainfo{Identificador}{mira-image-template} \hfill \metainfo{Ficha}{mira-image-template.md}}

\begin{descricao}
Agente que deriva um template visual consistente a partir de uma imagem de referência para reutilização no MIRA.
\end{descricao}

\begin{funcao}
No Core, este perfil apoia tarefas relacionadas a agente que deriva um template visual consistente a partir de uma imagem de referência para reutilização no MIRA.. O orquestrador pode encaminhar-lhe o trabalho na etapa correspondente; a resposta retorna ao fluxo principal para integração e conferência de fontes, critérios e permissões. A ficha documenta a função, mas não comprova que o perfil tenha sido chamado ou executado a tarefa.
\end{funcao}

\begin{arquitetura}
\textbf{Modo de execução declarado:} subagente · \textbf{Permissões:} edit: deny; bash: deny
\end{arquitetura}

\elemento{Mira qrcode (v1.0.0)}{\metainfo{Categoria}{Engenharia} \hfill \metainfo{Tipo}{mira-agent} \hfill \metainfo{Identificador}{mira-qrcode} \hfill \metainfo{Ficha}{mira-qrcode.md}}

\begin{descricao}
Agente utilitário do MIRA que gera e posiciona QR codes em slides para referências, downloads e interações externas.
\end{descricao}

\begin{funcao}
No Core, este perfil apoia tarefas relacionadas a agente utilitário do MIRA que gera e posiciona QR codes em slides para referências, downloads e interações externas.. O orquestrador pode encaminhar-lhe o trabalho na etapa correspondente; a resposta retorna ao fluxo principal para integração e conferência de fontes, critérios e permissões. A ficha documenta a função, mas não comprova que o perfil tenha sido chamado ou executado a tarefa.
\end{funcao}

\begin{arquitetura}
\textbf{Modo de execução declarado:} subagente · \textbf{Permissões:} edit: deny; bash: deny
\end{arquitetura}

\elemento{Mira size animator (v1.0.0)}{\metainfo{Categoria}{Engenharia} \hfill \metainfo{Tipo}{mira-agent} \hfill \metainfo{Identificador}{mira-size-animator} \hfill \metainfo{Ficha}{mira-size-animator.md}}

\begin{descricao}
Ajusta percepção de escala e tamanho nas animações MIRA
\end{descricao}

\begin{funcao}
No Core, este perfil apoia tarefas relacionadas a ajusta percepção de escala e tamanho nas animações MIRA. O orquestrador pode encaminhar-lhe o trabalho na etapa correspondente; a resposta retorna ao fluxo principal para integração e conferência de fontes, critérios e permissões. A ficha documenta a função, mas não comprova que o perfil tenha sido chamado ou executado a tarefa.
\end{funcao}

\begin{arquitetura}
\textbf{Modo de execução declarado:} subagente · \textbf{Permissões:} edit: deny; bash: deny
\end{arquitetura}

\elemento{Mira squared (v1.0.0)}{\metainfo{Categoria}{Engenharia} \hfill \metainfo{Tipo}{mira-agent} \hfill \metainfo{Identificador}{mira-squared} \hfill \metainfo{Ficha}{mira-squared.md}}

\begin{descricao}
Agente responsivo do MIRA que adapta apresentações para formato quadrado 1:1 preservando hierarquia visual.
\end{descricao}

\begin{funcao}
No Core, este perfil apoia tarefas relacionadas a agente responsivo do MIRA que adapta apresentações para formato quadrado 1:1 preservando hierarquia visual.. O orquestrador pode encaminhar-lhe o trabalho na etapa correspondente; a resposta retorna ao fluxo principal para integração e conferência de fontes, critérios e permissões. A ficha documenta a função, mas não comprova que o perfil tenha sido chamado ou executado a tarefa.
\end{funcao}

\begin{arquitetura}
\textbf{Modo de execução declarado:} subagente · \textbf{Permissões:} edit: deny; bash: deny
\end{arquitetura}

\elemento{Mira thirds (v1.0.0)}{\metainfo{Categoria}{Engenharia} \hfill \metainfo{Tipo}{mira-agent} \hfill \metainfo{Identificador}{mira-thirds} \hfill \metainfo{Ficha}{mira-thirds.md}}

\begin{descricao}
Agente responsivo do MIRA que aplica enquadramento segundo a regra dos terços para melhorar composição e foco visual.
\end{descricao}

\begin{funcao}
No Core, este perfil apoia tarefas relacionadas a agente responsivo do MIRA que aplica enquadramento segundo a regra dos terços para melhorar composição e foco visual.. O orquestrador pode encaminhar-lhe o trabalho na etapa correspondente; a resposta retorna ao fluxo principal para integração e conferência de fontes, critérios e permissões. A ficha documenta a função, mas não comprova que o perfil tenha sido chamado ou executado a tarefa.
\end{funcao}

\begin{arquitetura}
\textbf{Modo de execução declarado:} subagente · \textbf{Permissões:} edit: deny; bash: deny
\end{arquitetura}

\elemento{Mira vertical (v1.0.0)}{\metainfo{Categoria}{Engenharia} \hfill \metainfo{Tipo}{mira-agent} \hfill \metainfo{Identificador}{mira-vertical} \hfill \metainfo{Ficha}{mira-vertical.md}}

\begin{descricao}
Agente responsivo do MIRA que reorganiza o deck para consumo vertical em celular e vídeo curto.
\end{descricao}

\begin{funcao}
No Core, este perfil apoia tarefas relacionadas a agente responsivo do MIRA que reorganiza o deck para consumo vertical em celular e vídeo curto.. O orquestrador pode encaminhar-lhe o trabalho na etapa correspondente; a resposta retorna ao fluxo principal para integração e conferência de fontes, critérios e permissões. A ficha documenta a função, mas não comprova que o perfil tenha sido chamado ou executado a tarefa.
\end{funcao}

\begin{arquitetura}
\textbf{Modo de execução declarado:} subagente · \textbf{Permissões:} edit: deny; bash: deny
\end{arquitetura}

\elemento{Mira visuals (v1.0.0)}{\metainfo{Categoria}{Engenharia} \hfill \metainfo{Tipo}{mira-agent} \hfill \metainfo{Identificador}{mira-visuals} \hfill \metainfo{Ficha}{mira-visuals.md}}

\begin{descricao}
Agente visual do MIRA que produz ilustrações estáticas, painéis e infográficos alinhados ao conteúdo do deck.
\end{descricao}

\begin{funcao}
No Core, este perfil apoia tarefas relacionadas a agente visual do MIRA que produz ilustrações estáticas, painéis e infográficos alinhados ao conteúdo do deck.. O orquestrador pode encaminhar-lhe o trabalho na etapa correspondente; a resposta retorna ao fluxo principal para integração e conferência de fontes, critérios e permissões. A ficha documenta a função, mas não comprova que o perfil tenha sido chamado ou executado a tarefa.
\end{funcao}

\begin{arquitetura}
\textbf{Modo de execução declarado:} subagente · \textbf{Permissões:} edit: deny; bash: deny
\end{arquitetura}

\clearpage
\section{Categoria: Testes (1 perfis)}

\begin{descricao}
Reúne perfis orientados a testes. No Core, podem converter critérios em verificações e produzir evidência para o fluxo de qualidade; aprovação depende dos critérios e resultados observados (Módulo 4).
\end{descricao}

\elemento{Executor do ambiente de avaliação (v1.0.0)}{\metainfo{Categoria}{Testes} \hfill \metainfo{Tipo}{utility} \hfill \metainfo{Identificador}{Eval Runner} \hfill \metainfo{Ficha}{eval-runner.md}}

\begin{descricao}
Executa o harness de avaliação; a ficha orienta que este perfil não seja chamado diretamente.
\end{descricao}

\begin{funcao}
No Core, este perfil apoia tarefas relacionadas a executa o harness de avaliação. O orquestrador pode encaminhar-lhe o trabalho na etapa correspondente; a resposta retorna ao fluxo principal para integração e conferência de fontes, critérios e permissões. A ficha documenta a função, mas não comprova que o perfil tenha sido chamado ou executado a tarefa.
\end{funcao}

\begin{arquitetura}
\textbf{Modo de execução declarado:} subagente
\end{arquitetura}

\clearpage
\section{Categoria: Escrita e tradução (13 perfis)}

\begin{descricao}
Agrupa perfis de escrita, edição e tradução. No Core, apoiam tarefas de produção textual; fontes, terminologia e decisões editoriais devem ser conferidas segundo o objetivo do trabalho.
\end{descricao}

\elemento{Author voice guardian (v1.0.0)}{\metainfo{Categoria}{Escrita e tradução} \hfill \metainfo{Tipo}{literary-agent} \hfill \metainfo{Identificador}{author-voice-guardian} \hfill \metainfo{Ficha}{author-voice-guardian.md}}

\begin{descricao}
Guarda de voz autoral que audita traduções contra o perfil de voz da obra — marcadores regionais e orais, modernismos proibidos e deriva de registro pragmático — sempre com revisão humana para alto risco.
\end{descricao}

\begin{funcao}
No Core, este perfil apoia tarefas relacionadas a guarda de voz autoral que audita traduções contra o perfil de voz da obra — marcadores regionais e orais, modernismos proibidos e deriva de…. O orquestrador pode encaminhar-lhe o trabalho na etapa correspondente; a resposta retorna ao fluxo principal para integração e conferência de fontes, critérios e permissões. A ficha documenta a função, mas não comprova que o perfil tenha sido chamado ou executado a tarefa.
\end{funcao}

\begin{arquitetura}
\textbf{Modo de execução declarado:} subagente · \textbf{Permissões:} read: allow; glob: allow; grep: allow; edit: deny; bash: deny; task: deny; webfetch: deny; websearch: deny
\end{arquitetura}

\elemento{Back translation verifier (v1.0.0)}{\metainfo{Categoria}{Escrita e tradução} \hfill \metainfo{Tipo}{literary-agent} \hfill \metainfo{Identificador}{back-translation-verifier} \hfill \metainfo{Ficha}{back-translation-verifier.md}}

\begin{descricao}
Verificador determinístico de retrotradução que compara original e retrotradução quanto a números, entidades, negação, pragmática, comprimento e termos preservados — sem jamais aprovar equivalência.
\end{descricao}

\begin{funcao}
No Core, este perfil apoia tarefas relacionadas a verificador determinístico de retrotradução que compara original e retrotradução quanto a números, entidades, negação, pragmática, comprimento e termos preservados — sem jamais aprovar equivalência.. O orquestrador pode encaminhar-lhe o trabalho na etapa correspondente; a resposta retorna ao fluxo principal para integração e conferência de fontes, critérios e permissões. A ficha documenta a função, mas não comprova que o perfil tenha sido chamado ou executado a tarefa.
\end{funcao}

\begin{arquitetura}
\textbf{Modo de execução declarado:} subagente · \textbf{Permissões:} read: allow; glob: allow; grep: allow; edit: deny; bash: deny; task: deny; webfetch: deny; websearch: deny
\end{arquitetura}

\elemento{Cultural episteme agente (v1.0.0)}{\metainfo{Categoria}{Escrita e tradução} \hfill \metainfo{Tipo}{literary-agent} \hfill \metainfo{Identificador}{cultural-episteme-agent} \hfill \metainfo{Ficha}{cultural-episteme-agent.md}}

\begin{descricao}
Agente de Epistemes Culturais e Equivalência Interpretativa que audita traduções literárias quanto a voz, história, pragmática, símbolos e riscos de domesticação ou exotização, sempre com revisão humana.
\end{descricao}

\begin{funcao}
No Core, este perfil apoia tarefas relacionadas a agente de Epistemes Culturais e Equivalência Interpretativa que audita traduções literárias quanto a voz, história, pragmática, símbolos e riscos de domesticação ou exotização, sempre…. O orquestrador pode encaminhar-lhe o trabalho na etapa correspondente; a resposta retorna ao fluxo principal para integração e conferência de fontes, critérios e permissões. A ficha documenta a função, mas não comprova que o perfil tenha sido chamado ou executado a tarefa.
\end{funcao}

\begin{arquitetura}
\textbf{Modo de execução declarado:} subagente · \textbf{Permissões:} read: allow; glob: allow; grep: allow; edit: deny; bash: deny; task: deny; webfetch: deny; websearch: deny
\end{arquitetura}

\elemento{Literário character psychology phd (v1.0.0)}{\metainfo{Categoria}{Escrita e tradução} \hfill \metainfo{Tipo}{literary-agent} \hfill \metainfo{Identificador}{literary-character-psychology-phd} \hfill \metainfo{Ficha}{literary-character-psychology-phd.md}}

\begin{descricao}
Especialista PhD em personagens literários, psicologia narrativa, agência, desejo, conflito interno, transformação e relações dramáticas. Capacidades descritas: Analisa e desenvolve psicologia de personagens literários: agência, desejo, conflito interno, transformação.; Projeta relações entre personagens com base em tensão dramática e verossimilhança psicológica.
\end{descricao}

\begin{funcao}
No Core, este perfil apoia tarefas relacionadas a especialista PhD em personagens literários, psicologia narrativa, agência, desejo, conflito interno, transformação e relações dramáticas.. O orquestrador pode encaminhar-lhe o trabalho na etapa correspondente; a resposta retorna ao fluxo principal para integração e conferência de fontes, critérios e permissões. A ficha documenta a função, mas não comprova que o perfil tenha sido chamado ou executado a tarefa.
\end{funcao}

\begin{arquitetura}
\textbf{Modo de execução declarado:} subagente · \textbf{Permissões:} edit: deny; bash: deny
\end{arquitetura}

\elemento{Literário ethics trauma phd (v1.0.0)}{\metainfo{Categoria}{Escrita e tradução} \hfill \metainfo{Tipo}{literary-agent} \hfill \metainfo{Identificador}{literary-ethics-trauma-phd} \hfill \metainfo{Ficha}{literary-ethics-trauma-phd.md}}

\begin{descricao}
Especialista PhD em ética literária da representação, trauma, alteridade, violência institucional, memória histórica e anti-exploração estética.
\end{descricao}

\begin{funcao}
No Core, este perfil apoia tarefas relacionadas a especialista PhD em ética literária da representação, trauma, alteridade, violência institucional, memória histórica e anti-exploração estética.. O orquestrador pode encaminhar-lhe o trabalho na etapa correspondente; a resposta retorna ao fluxo principal para integração e conferência de fontes, critérios e permissões. A ficha documenta a função, mas não comprova que o perfil tenha sido chamado ou executado a tarefa.
\end{funcao}

\begin{arquitetura}
\textbf{Modo de execução declarado:} subagente · \textbf{Permissões:} edit: deny; bash: deny
\end{arquitetura}

\elemento{Literário innovation editorial phd (v1.0.0)}{\metainfo{Categoria}{Escrita e tradução} \hfill \metainfo{Tipo}{literary-agent} \hfill \metainfo{Identificador}{literary-innovation-editorial-phd} \hfill \metainfo{Ficha}{literary-innovation-editorial-phd.md}}

\begin{descricao}
Especialista PhD em inovação formal literária, materialidade editorial, paratextos, hipertexto impresso, design narrativo e contribuição potencial.
\end{descricao}

\begin{funcao}
No Core, este perfil apoia tarefas relacionadas a especialista PhD em inovação formal literária, materialidade editorial, paratextos, hipertexto impresso, design narrativo e contribuição potencial.. O orquestrador pode encaminhar-lhe o trabalho na etapa correspondente; a resposta retorna ao fluxo principal para integração e conferência de fontes, critérios e permissões. A ficha documenta a função, mas não comprova que o perfil tenha sido chamado ou executado a tarefa.
\end{funcao}

\begin{arquitetura}
\textbf{Modo de execução declarado:} subagente · \textbf{Permissões:} edit: deny; bash: deny
\end{arquitetura}

\elemento{Arquiteto de narratologia literária (v1.0.0)}{\metainfo{Categoria}{Escrita e tradução} \hfill \metainfo{Tipo}{literary-agent} \hfill \metainfo{Identificador}{literary-narratology-architect-phd} \hfill \metainfo{Ficha}{literary-narratology-architect-phd.md}}

\begin{descricao}
Especialista PhD em narratologia para arquitetura narrativa, enredo, temporalidade, focalização, rotas, partes e coerência estrutural. Capacidades descritas: Projeta estrutura macro da narrativa: enredo, temporalidade, focalização, partes.; Garante consistência entre timelines, pontos de vista e arcos narrativos.
\end{descricao}

\begin{funcao}
No Core, este perfil apoia tarefas relacionadas a especialista PhD em narratologia para arquitetura narrativa, enredo, temporalidade, focalização, rotas, partes e coerência estrutural.. O orquestrador pode encaminhar-lhe o trabalho na etapa correspondente; a resposta retorna ao fluxo principal para integração e conferência de fontes, critérios e permissões. A ficha documenta a função, mas não comprova que o perfil tenha sido chamado ou executado a tarefa.
\end{funcao}

\begin{arquitetura}
\textbf{Modo de execução declarado:} subagente · \textbf{Permissões:} edit: deny; bash: deny
\end{arquitetura}

\elemento{Literário orquestrador phd (v1.0.0)}{\metainfo{Categoria}{Escrita e tradução} \hfill \metainfo{Tipo}{literary-agent} \hfill \metainfo{Identificador}{literary-orchestrator-phd} \hfill \metainfo{Ficha}{literary-orchestrator-phd.md}}

\begin{descricao}
Orquestrador PhD de projetos literários para coordenar criação, estudo, crítica, scanners, pesquisa, revisão ética e entregáveis editoriais.
\end{descricao}

\begin{funcao}
No Core, este perfil apoia tarefas relacionadas a orquestrador PhD de projetos literários para coordenar criação, estudo, crítica, scanners, pesquisa, revisão ética e entregáveis editoriais.. O orquestrador pode encaminhar-lhe o trabalho na etapa correspondente; a resposta retorna ao fluxo principal para integração e conferência de fontes, critérios e permissões. A ficha documenta a função, mas não comprova que o perfil tenha sido chamado ou executado a tarefa.
\end{funcao}

\begin{arquitetura}
\textbf{Modo de execução declarado:} subagente · \textbf{Permissões:} edit: deny; bash: deny
\end{arquitetura}

\elemento{Literário pesquisa scholar phd (v1.0.0)}{\metainfo{Categoria}{Escrita e tradução} \hfill \metainfo{Tipo}{literary-agent} \hfill \metainfo{Identificador}{literary-research-scholar-phd} \hfill \metainfo{Ficha}{literary-research-scholar-phd.md}}

\begin{descricao}
Pesquisador PhD de busca e pesquisa literária para corpus comparativo, bibliografia, teoria, fontes, citações, lacunas e rigor internacional.
\end{descricao}

\begin{funcao}
No Core, este perfil apoia tarefas relacionadas a pesquisador PhD de busca e pesquisa literária para corpus comparativo, bibliografia, teoria, fontes, citações, lacunas e rigor internacional.. O orquestrador pode encaminhar-lhe o trabalho na etapa correspondente; a resposta retorna ao fluxo principal para integração e conferência de fontes, critérios e permissões. A ficha documenta a função, mas não comprova que o perfil tenha sido chamado ou executado a tarefa.
\end{funcao}

\begin{arquitetura}
\textbf{Modo de execução declarado:} subagente · \textbf{Permissões:} edit: deny; bash: deny
\end{arquitetura}

\elemento{Literário smoke minimal (v1.0.0)}{\metainfo{Categoria}{Escrita e tradução} \hfill \metainfo{Tipo}{literary-agent} \hfill \metainfo{Identificador}{literary-smoke-minimal} \hfill \metainfo{Ficha}{literary-smoke-minimal.md}}

\begin{descricao}
Agente mínimo de smoke test literário para isolar falhas de runtime, slug, model routing e registry dos agentes literary-.
\end{descricao}

\begin{funcao}
No Core, este perfil apoia tarefas relacionadas a agente mínimo de smoke test literário para isolar falhas de runtime, slug, model routing e registry dos agentes literary-.. O orquestrador pode encaminhar-lhe o trabalho na etapa correspondente; a resposta retorna ao fluxo principal para integração e conferência de fontes, critérios e permissões. A ficha documenta a função, mas não comprova que o perfil tenha sido chamado ou executado a tarefa.
\end{funcao}

\begin{arquitetura}
\textbf{Modo de execução declarado:} subagente · \textbf{Permissões:} edit: deny; bash: deny
\end{arquitetura}

\elemento{Literário style voice phd (v1.0.0)}{\metainfo{Categoria}{Escrita e tradução} \hfill \metainfo{Tipo}{literary-agent} \hfill \metainfo{Identificador}{literary-style-voice-phd} \hfill \metainfo{Ficha}{literary-style-voice-phd.md}}

\begin{descricao}
Especialista PhD em estilo literário, voz, ritmo, léxico, registro, dicção, musicalidade, revisão de prosa e assinatura discursiva.
\end{descricao}

\begin{funcao}
No Core, este perfil apoia tarefas relacionadas a especialista PhD em estilo literário, voz, ritmo, léxico, registro, dicção, musicalidade, revisão de prosa e assinatura discursiva.. O orquestrador pode encaminhar-lhe o trabalho na etapa correspondente; a resposta retorna ao fluxo principal para integração e conferência de fontes, critérios e permissões. A ficha documenta a função, mas não comprova que o perfil tenha sido chamado ou executado a tarefa.
\end{funcao}

\begin{arquitetura}
\textbf{Modo de execução declarado:} subagente · \textbf{Permissões:} edit: deny; bash: deny
\end{arquitetura}

\elemento{Literário symbolic imagery phd (v1.0.0)}{\metainfo{Categoria}{Escrita e tradução} \hfill \metainfo{Tipo}{literary-agent} \hfill \metainfo{Identificador}{literary-symbolic-imagery-phd} \hfill \metainfo{Ficha}{literary-symbolic-imagery-phd.md}}

\begin{descricao}
Especialista PhD em símbolos, motivos recorrentes, imagens, campos sensoriais, metáforas, arquétipos e coesão simbólica literária. Capacidades descritas: Mapeia símbolos, motivos recorrentes e arquétipos na obra literária.; Garante consistência entre metáforas, campos sensoriais e universo simbólico.
\end{descricao}

\begin{funcao}
No Core, este perfil apoia tarefas relacionadas a especialista PhD em símbolos, motivos recorrentes, imagens, campos sensoriais, metáforas, arquétipos e coesão simbólica literária.. O orquestrador pode encaminhar-lhe o trabalho na etapa correspondente; a resposta retorna ao fluxo principal para integração e conferência de fontes, critérios e permissões. A ficha documenta a função, mas não comprova que o perfil tenha sido chamado ou executado a tarefa.
\end{funcao}

\begin{arquitetura}
\textbf{Modo de execução declarado:} subagente · \textbf{Permissões:} edit: deny; bash: deny
\end{arquitetura}

\elemento{Terminology grafo agente (v1.0.0)}{\metainfo{Categoria}{Escrita e tradução} \hfill \metainfo{Tipo}{literary-agent} \hfill \metainfo{Identificador}{terminology-graph-agent} \hfill \metainfo{Ficha}{terminology-graph-agent.md}}

\begin{descricao}
Grafo terminológico trilíngue (PT-BR/EN/ZH-CN) que consome deltas do CulturalEpistemeAgent, exige aprovação humana por termo e detecta conflito terminológico e deriva simbólica com gate fail-closed.
\end{descricao}

\begin{funcao}
No Core, este perfil apoia tarefas relacionadas a grafo terminológico trilíngue (PT-BR/EN/ZH-CN) que consome deltas do CulturalEpistemeAgent, exige aprovação humana por termo e detecta conflito terminológico e deriva simbólica com gate fail-closed.. O orquestrador pode encaminhar-lhe o trabalho na etapa correspondente; a resposta retorna ao fluxo principal para integração e conferência de fontes, critérios e permissões. A ficha documenta a função, mas não comprova que o perfil tenha sido chamado ou executado a tarefa.
\end{funcao}

\begin{arquitetura}
\textbf{Modo de execução declarado:} subagente · \textbf{Permissões:} read: allow; glob: allow; grep: allow; edit: deny; bash: deny; task: deny; webfetch: deny; websearch: deny
\end{arquitetura}

\clearpage
\section{Categoria: Ferramentas e apoio transversal (128 perfis)}

\begin{descricao}
Categoria transversal com perfis variados de repositório, documentação, automação e suporte. Leia a descrição individual para saber a tarefa específica; não presuma que todos executem a mesma função.
\end{descricao}

\elemento{Gestor de decisões arquiteturais (v1.0.0)}{\metainfo{Categoria}{Ferramentas e apoio transversal} \hfill \metainfo{Tipo}{especialista} \hfill \metainfo{Identificador}{adr-manager} \hfill \metainfo{Ficha}{adr-manager.md}}

\begin{descricao}
Registra decisões arquiteturais em formato ADR, incluindo contexto, alternativas consideradas e consequências.
\end{descricao}

\begin{funcao}
No Core, este perfil apoia tarefas relacionadas a registra decisões arquiteturais em formato ADR, incluindo contexto, alternativas consideradas e consequências.. O orquestrador pode encaminhar-lhe o trabalho na etapa correspondente; a resposta retorna ao fluxo principal para integração e conferência de fontes, critérios e permissões. A ficha documenta a função, mas não comprova que o perfil tenha sido chamado ou executado a tarefa.
\end{funcao}

\begin{arquitetura}
\textbf{Modo de execução declarado:} subagente · \textbf{Permissões:} edit: deny; bash: deny
\end{arquitetura}

\elemento{Moderador de fórum e debates multiagente (v1.0.0)}{\metainfo{Categoria}{Ferramentas e apoio transversal} \hfill \metainfo{Tipo}{especialista} \hfill \metainfo{Identificador}{Agente Reversa: Agent Forum / Debate Moderator} \hfill \metainfo{Ficha}{reversa-agent-forum.md}}

\begin{descricao}
Modera debates entre agentes especializados em quatro etapas — abertura, discussão, síntese e conclusão — com controle das falas e consolidação do resultado.
\end{descricao}

\begin{funcao}
No Core, este perfil apoia tarefas relacionadas a modera debates entre agentes especializados em quatro etapas — abertura, discussão, síntese e conclusão — com controle das falas e consolidação do resultado.. O orquestrador pode encaminhar-lhe o trabalho na etapa correspondente; a resposta retorna ao fluxo principal para integração e conferência de fontes, critérios e permissões. A ficha documenta a função, mas não comprova que o perfil tenha sido chamado ou executado a tarefa.
\end{funcao}

\begin{arquitetura}
\textbf{Modo de execução declarado:} subagente
\end{arquitetura}

\elemento{Pipeline de extração e análise de documentos (v1.0.0)}{\metainfo{Categoria}{Ferramentas e apoio transversal} \hfill \metainfo{Tipo}{especialista} \hfill \metainfo{Identificador}{Agente Reversa: Document IR Report Pipeline} \hfill \metainfo{Ficha}{reversa-document-ir.md}}

\begin{descricao}
Coordena um fluxo de extração de informação de documentos, da análise do material à produção de um relatório estruturado.
\end{descricao}

\begin{funcao}
No Core, este perfil apoia tarefas relacionadas a coordena um fluxo de extração de informação de documentos, da análise do material à produção de um relatório estruturado.. O orquestrador pode encaminhar-lhe o trabalho na etapa correspondente; a resposta retorna ao fluxo principal para integração e conferência de fontes, critérios e permissões. A ficha documenta a função, mas não comprova que o perfil tenha sido chamado ou executado a tarefa.
\end{funcao}

\begin{arquitetura}
\textbf{Modo de execução declarado:} subagente
\end{arquitetura}

\elemento{Gestor de Agente Reversa: Process Lifecycle (v1.0.0)}{\metainfo{Categoria}{Ferramentas e apoio transversal} \hfill \metainfo{Tipo}{especialista} \hfill \metainfo{Identificador}{Agente Reversa: Process Lifecycle Manager} \hfill \metainfo{Ficha}{reversa-process-lifecycle.md}}

\begin{descricao}
Agente gerenciador de ciclo de vida de processos background. Inspirado pelo SimulationRunner do MiroFish-Offline (simulation\_runner.py). Inicia, monitora, pausa, retoma e finaliza processos com tracking cross-platform e ingestão de logs em tempo real. Use via: "processo", "runner", "background", /process-lifecycle.
\end{descricao}

\begin{funcao}
No Core, este perfil apoia tarefas relacionadas a agente gerenciador de ciclo de vida de processos background.. O orquestrador pode encaminhar-lhe o trabalho na etapa correspondente; a resposta retorna ao fluxo principal para integração e conferência de fontes, critérios e permissões. A ficha documenta a função, mas não comprova que o perfil tenha sido chamado ou executado a tarefa.
\end{funcao}

\begin{arquitetura}
\textbf{Modo de execução declarado:} subagente
\end{arquitetura}

\elemento{Antigravity cli}{\metainfo{Categoria}{Ferramentas e apoio transversal} \hfill \metainfo{Tipo}{integration} \hfill \metainfo{Identificador}{antigravity-cli} \hfill \metainfo{Ficha}{antigravity-cli.md}}

\begin{descricao}
Runner direto do Antigravity CLI agy (google-antigravity, oficial): status/version/run/doctor no padrão M7, com detecção de falha silenciosa. Complementa bridge + MCP + executor (não substitui).
\end{descricao}

\begin{funcao}
No Core, este perfil apoia tarefas relacionadas a runner direto do Antigravity CLI agy (google-antigravity, oficial): status/version/run/doctor no padrão M7, com detecção de falha silenciosa.. O orquestrador pode encaminhar-lhe o trabalho na etapa correspondente; a resposta retorna ao fluxo principal para integração e conferência de fontes, critérios e permissões. A ficha documenta a função, mas não comprova que o perfil tenha sido chamado ou executado a tarefa.
\end{funcao}

\begin{arquitetura}
\textbf{Modo de execução declarado:} subagente · \textbf{Permissões:} edit: deny; bash: deny
\end{arquitetura}

\elemento{Orquestrador de Antigravity}{\metainfo{Categoria}{Ferramentas e apoio transversal} \hfill \metainfo{Tipo}{especialista} \hfill \metainfo{Identificador}{AntigravityOrchestrator} \hfill \metainfo{Ficha}{antigravity-orchestrator.md}}

\begin{descricao}
Orquestrador especializado que delega tarefas ao Antigravity (Google DeepMind Advanced Agentic Coding), expondo e coordenando suas capacidades exclusivas no pipeline do OpenCode Ecosystem.
\end{descricao}

\begin{funcao}
No Core, este perfil apoia tarefas relacionadas a orquestrador especializado que delega tarefas ao Antigravity (Google DeepMind Advanced Agentic Coding), expondo e coordenando suas capacidades exclusivas no pipeline do OpenCode Ecosystem.. O orquestrador pode encaminhar-lhe o trabalho na etapa correspondente; a resposta retorna ao fluxo principal para integração e conferência de fontes, critérios e permissões. A ficha documenta a função, mas não comprova que o perfil tenha sido chamado ou executado a tarefa.
\end{funcao}

\begin{arquitetura}
\textbf{Modo de execução declarado:} subagente
\end{arquitetura}

\elemento{Arquiteto de software (v1.0.0)}{\metainfo{Categoria}{Ferramentas e apoio transversal} \hfill \metainfo{Tipo}{especialista} \hfill \metainfo{Identificador}{architect} \hfill \metainfo{Ficha}{architect.md}}

\begin{descricao}
Projeta arquitetura de software e toma decisoes de design
\end{descricao}

\begin{funcao}
No Core, este perfil apoia tarefas relacionadas a projeta arquitetura de software e toma decisoes de design. O orquestrador pode encaminhar-lhe o trabalho na etapa correspondente; a resposta retorna ao fluxo principal para integração e conferência de fontes, critérios e permissões. A ficha documenta a função, mas não comprova que o perfil tenha sido chamado ou executado a tarefa.
\end{funcao}

\begin{arquitetura}
\textbf{Modo de execução declarado:} subagente · \textbf{Permissões:} edit: deny; bash: deny
\end{arquitetura}

\elemento{Analisador de arquitetura de software (v1.0.0)}{\metainfo{Categoria}{Ferramentas e apoio transversal} \hfill \metainfo{Tipo}{especialista} \hfill \metainfo{Identificador}{architecture-analyzer} \hfill \metainfo{Ficha}{architecture-analyzer.md}}

\begin{descricao}
Analisa arquitetura orientada a domínio e identifica contextos delimitados, fronteiras entre módulos e relações de domínio para fluxos em várias etapas.
\end{descricao}

\begin{funcao}
No Core, este perfil apoia tarefas relacionadas a analisa arquitetura orientada a domínio e identifica contextos delimitados, fronteiras entre módulos e relações de domínio para fluxos em várias etapas.. O orquestrador pode encaminhar-lhe o trabalho na etapa correspondente; a resposta retorna ao fluxo principal para integração e conferência de fontes, critérios e permissões. A ficha documenta a função, mas não comprova que o perfil tenha sido chamado ou executado a tarefa.
\end{funcao}

\begin{arquitetura}
\textbf{Modo de execução declarado:} subagente · \textbf{Permissões:} edit: deny; bash: deny
\end{arquitetura}

\elemento{Autoevolve (v1.0.0)}{\metainfo{Categoria}{Ferramentas e apoio transversal} \hfill \metainfo{Tipo}{especialista} \hfill \metainfo{Identificador}{autoevolve} \hfill \metainfo{Ficha}{autoevolve.md}}

\begin{descricao}
AutoEvolve — engine de evolução autônoma do ecossistema OpenCode v5.1. Roteia subcomandos (/evolve status; discover; install; verify; update; learn), executa pipeline SENSE→DISCOVER→INSTALL→VERIFY→EVOLVE→LEARN, integra SPEC-025+SPEC-026+SPEC-027.
\end{descricao}

\begin{funcao}
No Core, este perfil apoia tarefas relacionadas a autoEvolve — engine de evolução autônoma do ecossistema OpenCode v5.1.. O orquestrador pode encaminhar-lhe o trabalho na etapa correspondente; a resposta retorna ao fluxo principal para integração e conferência de fontes, critérios e permissões. A ficha documenta a função, mas não comprova que o perfil tenha sido chamado ou executado a tarefa.
\end{funcao}

\begin{arquitetura}
\textbf{Modo de execução declarado:} subagente · \textbf{Permissões:} read: allow; grep: allow; glob: allow; bash: allow; edit: allow; webfetch: allow
\end{arquitetura}

\elemento{Awesome mcp servers}{\metainfo{Categoria}{Ferramentas e apoio transversal} \hfill \metainfo{Tipo}{curation} \hfill \metainfo{Identificador}{awesome-mcp-servers} \hfill \metainfo{Ficha}{awesome-mcp-servers.md}}

\begin{descricao}
Curadoria executável da lista wong2/awesome-mcp-servers (SPEC-935-R652): referência-primeiro (fetch, sequential-thinking, filesystem registrados), SaaS com trust gate, adiamentos motivados.
\end{descricao}

\begin{funcao}
No Core, este perfil apoia tarefas relacionadas a curadoria executável da lista wong2/awesome-mcp-servers (SPEC-935-R652): referência-primeiro (fetch, sequential-thinking, filesystem registrados), SaaS com trust gate, adiamentos motivados.. O orquestrador pode encaminhar-lhe o trabalho na etapa correspondente; a resposta retorna ao fluxo principal para integração e conferência de fontes, critérios e permissões. A ficha documenta a função, mas não comprova que o perfil tenha sido chamado ou executado a tarefa.
\end{funcao}

\begin{arquitetura}
\textbf{Modo de execução declarado:} subagente · \textbf{Permissões:} edit: deny; bash: deny
\end{arquitetura}

\elemento{Executor de tarefas em lote (v1.0.0)}{\metainfo{Categoria}{Ferramentas e apoio transversal} \hfill \metainfo{Tipo}{especialista} \hfill \metainfo{Identificador}{batch-executor} \hfill \metainfo{Ficha}{batch-executor.md}}

\begin{descricao}
Executa tarefas em lotes paralelos, coordena delegações simultâneas ao CoderAgent e acompanha a conclusão do lote.
\end{descricao}

\begin{funcao}
No Core, este perfil apoia tarefas relacionadas a executa tarefas em lotes paralelos, coordena delegações simultâneas ao CoderAgent e acompanha a conclusão do lote.. O orquestrador pode encaminhar-lhe o trabalho na etapa correspondente; a resposta retorna ao fluxo principal para integração e conferência de fontes, critérios e permissões. A ficha documenta a função, mas não comprova que o perfil tenha sido chamado ou executado a tarefa.
\end{funcao}

\begin{arquitetura}
\textbf{Modo de execução declarado:} subagente · \textbf{Permissões:} edit: deny; bash: deny
\end{arquitetura}

\elemento{Orquestrador de bernstein}{\metainfo{Categoria}{Ferramentas e apoio transversal} \hfill \metainfo{Tipo}{especialista} \hfill \metainfo{Identificador}{bernstein-orchestrator} \hfill \metainfo{Ficha}{bernstein-orchestrator.md}}

\begin{descricao}
Bernstein é o maestro do ecossistema OpenCode. Ele orquestra agentes CLI coding (Claude, Codex, Gemini, Qwen) em pipelines multi-agente com:
\end{descricao}

\begin{funcao}
No Core, este perfil apoia tarefas relacionadas a bernstein é o maestro do ecossistema OpenCode.. O orquestrador pode encaminhar-lhe o trabalho na etapa correspondente; a resposta retorna ao fluxo principal para integração e conferência de fontes, critérios e permissões. A ficha documenta a função, mas não comprova que o perfil tenha sido chamado ou executado a tarefa.
\end{funcao}

\begin{arquitetura}
\textbf{Modo de execução declarado:} subagente · \textbf{Permissões:} edit: deny; bash: deny
\end{arquitetura}

\elemento{Agente de compilação e validação (v1.0.0)}{\metainfo{Categoria}{Ferramentas e apoio transversal} \hfill \metainfo{Tipo}{especialista} \hfill \metainfo{Identificador}{build-agent} \hfill \metainfo{Ficha}{build-agent.md}}

\begin{descricao}
Confere os tipos e valida a compilação do projeto.
\end{descricao}

\begin{funcao}
No Core, este perfil apoia tarefas relacionadas a confere os tipos e valida a compilação do projeto.. O orquestrador pode encaminhar-lhe o trabalho na etapa correspondente; a resposta retorna ao fluxo principal para integração e conferência de fontes, critérios e permissões. A ficha documenta a função, mas não comprova que o perfil tenha sido chamado ou executado a tarefa.
\end{funcao}

\begin{arquitetura}
\textbf{Modo de execução declarado:} subagente · \textbf{Permissões:} edit: deny; bash: deny
\end{arquitetura}

\elemento{Catálogo de Skills nuvem (Antigravity Backup) (v1.0.0)}{\metainfo{Categoria}{Ferramentas e apoio transversal} \hfill \metainfo{Tipo}{especialista} \hfill \metainfo{Identificador}{Catálogo de Skills Cloud (Antigravity Backup)} \hfill \metainfo{Ficha}{skills-cloud-antigravity.md}}

\begin{descricao}
56 skills de infraestrutura Google Cloud Platform (AlloyDB, Cloud SQL, BigQuery, Dataflow, Cloud Composer, GCS Security, Firestore, Spanner) com 269 scripts operacionais importados do backup Antigravity 2026-07-11. Licença Apache 2.0 (Google).
\end{descricao}

\begin{funcao}
No Core, este perfil apoia tarefas relacionadas a 56 skills de infraestrutura Google Cloud Platform (AlloyDB, Cloud SQL, BigQuery, Dataflow, Cloud Composer, GCS Security, Firestore, Spanner) com 269 scripts operacionais importados do…. O orquestrador pode encaminhar-lhe o trabalho na etapa correspondente; a resposta retorna ao fluxo principal para integração e conferência de fontes, critérios e permissões. A ficha documenta a função, mas não comprova que o perfil tenha sido chamado ou executado a tarefa.
\end{funcao}

\begin{arquitetura}
\textbf{Modo de execução declarado:} subagente
\end{arquitetura}

\elemento{Claude agente sdk python}{\metainfo{Categoria}{Ferramentas e apoio transversal} \hfill \metainfo{Tipo}{integration} \hfill \metainfo{Identificador}{claude-agent-sdk-python} \hfill \metainfo{Ficha}{claude-agent-sdk-python.md}}

\begin{descricao}
Ponte do Claude Agent SDK Python (anthropics, MIT + Commercial Terms): query()/ClaudeSDKClient, in-process MCP, hooks; monta options sem executar.
\end{descricao}

\begin{funcao}
No Core, este perfil apoia tarefas relacionadas a ponte do Claude Agent SDK Python (anthropics, MIT + Commercial Terms): query()/ClaudeSDKClient, in-process MCP, hooks. O orquestrador pode encaminhar-lhe o trabalho na etapa correspondente; a resposta retorna ao fluxo principal para integração e conferência de fontes, critérios e permissões. A ficha documenta a função, mas não comprova que o perfil tenha sido chamado ou executado a tarefa.
\end{funcao}

\begin{arquitetura}
\textbf{Modo de execução declarado:} subagente · \textbf{Permissões:} edit: deny; bash: deny
\end{arquitetura}

\elemento{Claude código harness}{\metainfo{Categoria}{Ferramentas e apoio transversal} \hfill \metainfo{Tipo}{curation} \hfill \metainfo{Identificador}{claude-code-harness} \hfill \metainfo{Ficha}{claude-code-harness.md}}

\begin{descricao}
Curadoria do Claude Code Harness (Chachamaru127, MIT): loop Plan→Work→Review→Release, 5 verb skills, rota OpenCode internal-compatible. Referência metodológica; não executa o harness.
\end{descricao}

\begin{funcao}
No Core, este perfil apoia tarefas relacionadas a curadoria do Claude Code Harness (Chachamaru127, MIT): loop Plan→Work→Review→Release, 5 verb skills, rota OpenCode internal-compatible.. O orquestrador pode encaminhar-lhe o trabalho na etapa correspondente; a resposta retorna ao fluxo principal para integração e conferência de fontes, critérios e permissões. A ficha documenta a função, mas não comprova que o perfil tenha sido chamado ou executado a tarefa.
\end{funcao}

\begin{arquitetura}
\textbf{Modo de execução declarado:} subagente · \textbf{Permissões:} edit: deny; bash: deny
\end{arquitetura}

\elemento{Claude plugins official}{\metainfo{Categoria}{Ferramentas e apoio transversal} \hfill \metainfo{Tipo}{curation} \hfill \metainfo{Identificador}{claude-plugins-official} \hfill \metainfo{Ficha}{claude-plugins-official.md}}

\begin{descricao}
Curadoria do marketplace Claude Code Plugins (anthropics, licença por plugin): /plugins + /external\_plugins, install por marketplace, gate de trust obrigatório (Anthropic não audita conteúdo dos plugins).
\end{descricao}

\begin{funcao}
No Core, este perfil apoia tarefas relacionadas a curadoria do marketplace Claude Code Plugins (anthropics, licença por plugin): /plugins + /external\_plugins, install por marketplace, gate de trust obrigatório (Anthropic não audita conteúdo…. O orquestrador pode encaminhar-lhe o trabalho na etapa correspondente; a resposta retorna ao fluxo principal para integração e conferência de fontes, critérios e permissões. A ficha documenta a função, mas não comprova que o perfil tenha sido chamado ou executado a tarefa.
\end{funcao}

\begin{arquitetura}
\textbf{Modo de execução declarado:} subagente · \textbf{Permissões:} edit: deny; bash: deny
\end{arquitetura}

\elemento{Especialista em nuvem Alloy DB (v1.0.0)}{\metainfo{Categoria}{Ferramentas e apoio transversal} \hfill \metainfo{Tipo}{especialista} \hfill \metainfo{Identificador}{Cloud AlloyDB Specialist} \hfill \metainfo{Ficha}{cloud-alloydb-specialist.md}}

\begin{descricao}
Especialista em AlloyDB Omni e AlloyDB PostgreSQL — administração, saúde, monitoramento, otimização, performance, replicação e operações Kubernetes. Baseado em 16 skills do Antigravity Backup (SPEC-935-R130).
\end{descricao}

\begin{funcao}
No Core, este perfil apoia tarefas relacionadas a especialista em AlloyDB Omni e AlloyDB PostgreSQL — administração, saúde, monitoramento, otimização, performance, replicação e operações Kubernetes.. O orquestrador pode encaminhar-lhe o trabalho na etapa correspondente; a resposta retorna ao fluxo principal para integração e conferência de fontes, critérios e permissões. A ficha documenta a função, mas não comprova que o perfil tenha sido chamado ou executado a tarefa.
\end{funcao}

\begin{arquitetura}
\textbf{Modo de execução declarado:} subagente
\end{arquitetura}

\elemento{Especialista em nuvem Big Query (v1.0.0)}{\metainfo{Categoria}{Ferramentas e apoio transversal} \hfill \metainfo{Tipo}{especialista} \hfill \metainfo{Identificador}{Cloud BigQuery Specialist} \hfill \metainfo{Ficha}{cloud-bigquery-specialist.md}}

\begin{descricao}
Especialista em BigQuery — SQL, ML/AI, BigFrames, Graph Analytics, Data Transfer Service, Dataform e dbt. Baseado em 4 skills do Antigravity Backup (SPEC-935-R130).
\end{descricao}

\begin{funcao}
No Core, este perfil apoia tarefas relacionadas a especialista em BigQuery — SQL, ML/AI, BigFrames, Graph Analytics, Data Transfer Service, Dataform e dbt.. O orquestrador pode encaminhar-lhe o trabalho na etapa correspondente; a resposta retorna ao fluxo principal para integração e conferência de fontes, critérios e permissões. A ficha documenta a função, mas não comprova que o perfil tenha sido chamado ou executado a tarefa.
\end{funcao}

\begin{arquitetura}
\textbf{Modo de execução declarado:} subagente
\end{arquitetura}

\elemento{Nuvem dados Infra Generalist (v1.0.0)}{\metainfo{Categoria}{Ferramentas e apoio transversal} \hfill \metainfo{Tipo}{especialista} \hfill \metainfo{Identificador}{Cloud Data Infra Generalist} \hfill \metainfo{Ficha}{cloud-data-infra-generalist.md}}

\begin{descricao}
Generalista em infraestrutura de dados GCP — Firestore, Spanner, descoberta de assets, building data apps, ML best practices, notebooks e migrações Airflow. Baseado em 8 skills do Antigravity Backup (SPEC-935-R130).
\end{descricao}

\begin{funcao}
No Core, este perfil apoia tarefas relacionadas a generalista em infraestrutura de dados GCP — Firestore, Spanner, descoberta de assets, building data apps, ML best practices, notebooks e migrações Airflow.. O orquestrador pode encaminhar-lhe o trabalho na etapa correspondente; a resposta retorna ao fluxo principal para integração e conferência de fontes, critérios e permissões. A ficha documenta a função, mas não comprova que o perfil tenha sido chamado ou executado a tarefa.
\end{funcao}

\begin{arquitetura}
\textbf{Modo de execução declarado:} subagente
\end{arquitetura}

\elemento{Especialista em nuvem dados Pipelines (v1.0.0)}{\metainfo{Categoria}{Ferramentas e apoio transversal} \hfill \metainfo{Tipo}{especialista} \hfill \metainfo{Identificador}{Cloud Data Pipelines Specialist} \hfill \metainfo{Ficha}{cloud-data-pipelines-specialist.md}}

\begin{descricao}
Especialista em pipelines de dados GCP — Dataflow, Cloud Composer/Airflow, Dataproc/Spark, orquestração de pipelines, provisionamento de recursos, Dataform e federação lakehouse. Baseado em 8 skills do Antigravity Backup.
\end{descricao}

\begin{funcao}
No Core, este perfil apoia tarefas relacionadas a especialista em pipelines de dados GCP — Dataflow, Cloud Composer/Airflow, Dataproc/Spark, orquestração de pipelines, provisionamento de recursos, Dataform e federação lakehouse.. O orquestrador pode encaminhar-lhe o trabalho na etapa correspondente; a resposta retorna ao fluxo principal para integração e conferência de fontes, critérios e permissões. A ficha documenta a função, mas não comprova que o perfil tenha sido chamado ou executado a tarefa.
\end{funcao}

\begin{arquitetura}
\textbf{Modo de execução declarado:} subagente
\end{arquitetura}

\elemento{Especialista em nuvem segurança (v1.0.0)}{\metainfo{Categoria}{Ferramentas e apoio transversal} \hfill \metainfo{Tipo}{especialista} \hfill \metainfo{Identificador}{Cloud Security Specialist} \hfill \metainfo{Ficha}{cloud-security-specialist.md}}

\begin{descricao}
Especialista em segurança GCP — avaliação de postura GCS, prevenção de perda de dados, verificação de autenticação e análise de combinações tóxicas de permissões. Baseado em 3 skills do Antigravity Backup (SPEC-935-R130).
\end{descricao}

\begin{funcao}
No Core, este perfil apoia tarefas relacionadas a especialista em segurança GCP — avaliação de postura GCS, prevenção de perda de dados, verificação de autenticação e análise de combinações tóxicas de permissões.. O orquestrador pode encaminhar-lhe o trabalho na etapa correspondente; a resposta retorna ao fluxo principal para integração e conferência de fontes, critérios e permissões. A ficha documenta a função, mas não comprova que o perfil tenha sido chamado ou executado a tarefa.
\end{funcao}

\begin{arquitetura}
\textbf{Modo de execução declarado:} subagente
\end{arquitetura}

\elemento{Especialista em nuvem SQL My SQL (v1.0.0)}{\metainfo{Categoria}{Ferramentas e apoio transversal} \hfill \metainfo{Tipo}{especialista} \hfill \metainfo{Identificador}{Cloud SQL MySQL Specialist} \hfill \metainfo{Ficha}{cloud-sql-mysql-specialist.md}}

\begin{descricao}
Especialista em Cloud SQL MySQL — administração, dados, ciclo de vida e monitoramento. Baseado em 4 skills do Antigravity Backup (SPEC-935-R130).
\end{descricao}

\begin{funcao}
No Core, este perfil apoia tarefas relacionadas a especialista em Cloud SQL MySQL — administração, dados, ciclo de vida e monitoramento.. O orquestrador pode encaminhar-lhe o trabalho na etapa correspondente; a resposta retorna ao fluxo principal para integração e conferência de fontes, critérios e permissões. A ficha documenta a função, mas não comprova que o perfil tenha sido chamado ou executado a tarefa.
\end{funcao}

\begin{arquitetura}
\textbf{Modo de execução declarado:} subagente
\end{arquitetura}

\elemento{Especialista em nuvem SQL Postgre SQL (v1.0.0)}{\metainfo{Categoria}{Ferramentas e apoio transversal} \hfill \metainfo{Tipo}{especialista} \hfill \metainfo{Identificador}{Cloud SQL PostgreSQL Specialist} \hfill \metainfo{Ficha}{cloud-sql-postgres-specialist.md}}

\begin{descricao}
Especialista em Cloud SQL PostgreSQL — administração, dados, saúde, ciclo de vida, monitoramento, replicação, vector assist e configuração. Baseado em 9 skills do Antigravity Backup (SPEC-935-R130).
\end{descricao}

\begin{funcao}
No Core, este perfil apoia tarefas relacionadas a especialista em Cloud SQL PostgreSQL — administração, dados, saúde, ciclo de vida, monitoramento, replicação, vector assist e configuração.. O orquestrador pode encaminhar-lhe o trabalho na etapa correspondente; a resposta retorna ao fluxo principal para integração e conferência de fontes, critérios e permissões. A ficha documenta a função, mas não comprova que o perfil tenha sido chamado ou executado a tarefa.
\end{funcao}

\begin{arquitetura}
\textbf{Modo de execução declarado:} subagente
\end{arquitetura}

\elemento{Especialista em nuvem SQL SQL Server (v1.0.0)}{\metainfo{Categoria}{Ferramentas e apoio transversal} \hfill \metainfo{Tipo}{especialista} \hfill \metainfo{Identificador}{Cloud SQL SQL Server Specialist} \hfill \metainfo{Ficha}{cloud-sql-sqlserver-specialist.md}}

\begin{descricao}
Especialista em Cloud SQL SQL Server — administração, dados, ciclo de vida e monitoramento. Baseado em 4 skills do Antigravity Backup (SPEC-935-R130).
\end{descricao}

\begin{funcao}
No Core, este perfil apoia tarefas relacionadas a especialista em Cloud SQL SQL Server — administração, dados, ciclo de vida e monitoramento.. O orquestrador pode encaminhar-lhe o trabalho na etapa correspondente; a resposta retorna ao fluxo principal para integração e conferência de fontes, critérios e permissões. A ficha documenta a função, mas não comprova que o perfil tenha sido chamado ou executado a tarefa.
\end{funcao}

\begin{arquitetura}
\textbf{Modo de execução declarado:} subagente
\end{arquitetura}

\elemento{Revisor de código (v1.0.0)}{\metainfo{Categoria}{Ferramentas e apoio transversal} \hfill \metainfo{Tipo}{especialista} \hfill \metainfo{Identificador}{code-reviewer} \hfill \metainfo{Ficha}{code-reviewer.md}}

\begin{descricao}
Voce e revisor de codigo senior. Foco: identificar problemas sem alterar codigo. - Corretude: bugs logicos, casos de borda, race conditions - Seguranca: injecao, XSS, autenticacao, exposicao de dados - Performance: loops ineficientes, memory leaks - Manutenibilidade: nomes claros, funcoes pequenas - Padroes: consistencia com codigo existente Arquivo/Linha -> Severidade -> Problema -> Sugestao
\end{descricao}

\begin{funcao}
No Core, este perfil apoia tarefas relacionadas a voce e revisor de codigo senior.. O orquestrador pode encaminhar-lhe o trabalho na etapa correspondente; a resposta retorna ao fluxo principal para integração e conferência de fontes, critérios e permissões. A ficha documenta a função, mas não comprova que o perfil tenha sido chamado ou executado a tarefa.
\end{funcao}

\begin{arquitetura}
\textbf{Modo de execução declarado:} subagente · \textbf{Permissões:} edit: deny; bash: deny
\end{arquitetura}

\elemento{Analisador de código-fonte (v1.0.0)}{\metainfo{Categoria}{Ferramentas e apoio transversal} \hfill \metainfo{Tipo}{especialista} \hfill \metainfo{Identificador}{codebase-analyzer} \hfill \metainfo{Ficha}{codebase-analyzer.md}}

\begin{descricao}
Analisa detalhes da implementação e ajuda a localizar e explicar componentes específicos do código-fonte.
\end{descricao}

\begin{funcao}
No Core, este perfil apoia tarefas relacionadas a analisa detalhes da implementação e ajuda a localizar e explicar componentes específicos do código-fonte.. O orquestrador pode encaminhar-lhe o trabalho na etapa correspondente; a resposta retorna ao fluxo principal para integração e conferência de fontes, critérios e permissões. A ficha documenta a função, mas não comprova que o perfil tenha sido chamado ou executado a tarefa.
\end{funcao}

\begin{arquitetura}
\textbf{Modo de execução declarado:} subagente · \textbf{Permissões:} read: allow; grep: allow; glob: allow; list: allow; bash: deny; edit: deny; todoread: deny; todowrite: deny; webfetch: deny
\end{arquitetura}

\elemento{Localizador de arquivos e componentes (v1.0.0)}{\metainfo{Categoria}{Ferramentas e apoio transversal} \hfill \metainfo{Tipo}{especialista} \hfill \metainfo{Identificador}{codebase-locator} \hfill \metainfo{Ficha}{codebase-locator.md}}

\begin{descricao}
Localiza arquivos, diretórios e componentes pertinentes a uma funcionalidade ou tarefa no repositório.
\end{descricao}

\begin{funcao}
No Core, este perfil apoia tarefas relacionadas a localiza arquivos, diretórios e componentes pertinentes a uma funcionalidade ou tarefa no repositório.. O orquestrador pode encaminhar-lhe o trabalho na etapa correspondente; a resposta retorna ao fluxo principal para integração e conferência de fontes, critérios e permissões. A ficha documenta a função, mas não comprova que o perfil tenha sido chamado ou executado a tarefa.
\end{funcao}

\begin{arquitetura}
\textbf{Modo de execução declarado:} subagente · \textbf{Permissões:} edit: deny; bash: deny
\end{arquitetura}

\elemento{Localizador de padrões no código (v1.0.0)}{\metainfo{Categoria}{Ferramentas e apoio transversal} \hfill \metainfo{Tipo}{especialista} \hfill \metainfo{Identificador}{codebase-pattern-finder} \hfill \metainfo{Ficha}{codebase-pattern-finder.md}}

\begin{descricao}
Localiza implementações semelhantes, exemplos de uso e padrões existentes que possam orientar uma mudança.
\end{descricao}

\begin{funcao}
No Core, este perfil apoia tarefas relacionadas a localiza implementações semelhantes, exemplos de uso e padrões existentes que possam orientar uma mudança.. O orquestrador pode encaminhar-lhe o trabalho na etapa correspondente; a resposta retorna ao fluxo principal para integração e conferência de fontes, critérios e permissões. A ficha documenta a função, mas não comprova que o perfil tenha sido chamado ou executado a tarefa.
\end{funcao}

\begin{arquitetura}
\textbf{Modo de execução declarado:} subagente · \textbf{Permissões:} read: allow; grep: allow; glob: allow; list: allow; bash: deny; edit: deny; todoread: deny; todowrite: deny; webfetch: deny
\end{arquitetura}

\elemento{Programador agente (v1.0.0)}{\metainfo{Categoria}{Ferramentas e apoio transversal} \hfill \metainfo{Tipo}{especialista} \hfill \metainfo{Identificador}{coder-agent} \hfill \metainfo{Ficha}{coder-agent.md}}

\begin{descricao}
Executes coding subtasks in sequence, ensuring completion as specified. Capacidades descritas: Executes coding subtasks in sequence, ensuring completion as specified
\end{descricao}

\begin{funcao}
No Core, este perfil apoia tarefas relacionadas a executes coding subtasks in sequence, ensuring completion as specified.. O orquestrador pode encaminhar-lhe o trabalho na etapa correspondente; a resposta retorna ao fluxo principal para integração e conferência de fontes, critérios e permissões. A ficha documenta a função, mas não comprova que o perfil tenha sido chamado ou executado a tarefa.
\end{funcao}

\begin{arquitetura}
\textbf{Modo de execução declarado:} subagente · \textbf{Permissões:} edit: deny; bash: deny
\end{arquitetura}

\elemento{Colab cli}{\metainfo{Categoria}{Ferramentas e apoio transversal} \hfill \metainfo{Tipo}{integration} \hfill \metainfo{Identificador}{colab-cli} \hfill \metainfo{Ficha}{colab-cli.md}}

\begin{descricao}
Executor externo orquestrável: colab (googlecolab/google-colab-cli, Apache-2.0). Provisiona runtimes CPU/GPU/TPU, executa código e jobs efêmeros, gerencia arquivos do Colab via terminal.
\end{descricao}

\begin{funcao}
No Core, este perfil apoia tarefas relacionadas a executor externo orquestrável: colab (googlecolab/google-colab-cli, Apache-2.0).. O orquestrador pode encaminhar-lhe o trabalho na etapa correspondente; a resposta retorna ao fluxo principal para integração e conferência de fontes, critérios e permissões. A ficha documenta a função, mas não comprova que o perfil tenha sido chamado ou executado a tarefa.
\end{funcao}

\begin{arquitetura}
\textbf{Modo de execução declarado:} subagente · \textbf{Permissões:} edit: deny; bash: deny
\end{arquitetura}

\elemento{Colab mcp}{\metainfo{Categoria}{Ferramentas e apoio transversal} \hfill \metainfo{Tipo}{integration} \hfill \metainfo{Identificador}{colab-mcp} \hfill \metainfo{Ficha}{colab-mcp.md}}

\begin{descricao}
Ponte orquestrável do Colab MCP Server (googlecolab/colab-mcp, Apache-2.0): codificação assistida interativa dentro do notebook Colab, complemento do fluxo terminal do colab-cli. Expõe presença, versão e config stdio.
\end{descricao}

\begin{funcao}
No Core, este perfil apoia tarefas relacionadas a ponte orquestrável do Colab MCP Server (googlecolab/colab-mcp, Apache-2.0): codificação assistida interativa dentro do notebook Colab, complemento do fluxo terminal do colab-cli.. O orquestrador pode encaminhar-lhe o trabalho na etapa correspondente; a resposta retorna ao fluxo principal para integração e conferência de fontes, critérios e permissões. A ficha documenta a função, mas não comprova que o perfil tenha sido chamado ou executado a tarefa.
\end{funcao}

\begin{arquitetura}
\textbf{Modo de execução declarado:} subagente · \textbf{Permissões:} edit: deny; bash: deny
\end{arquitetura}

\elemento{Gestor de contexto do projeto (v1.0.0)}{\metainfo{Categoria}{Ferramentas e apoio transversal} \hfill \metainfo{Tipo}{especialista} \hfill \metainfo{Identificador}{context-manager} \hfill \metainfo{Ficha}{context-manager.md}}

\begin{descricao}
Descobre, cataloga, valida e mantém a estrutura de contexto do projeto, acompanhando suas dependências.
\end{descricao}

\begin{funcao}
No Core, este perfil apoia tarefas relacionadas a descobre, cataloga, valida e mantém a estrutura de contexto do projeto, acompanhando suas dependências.. O orquestrador pode encaminhar-lhe o trabalho na etapa correspondente; a resposta retorna ao fluxo principal para integração e conferência de fontes, critérios e permissões. A ficha documenta a função, mas não comprova que o perfil tenha sido chamado ou executado a tarefa.
\end{funcao}

\begin{arquitetura}
\textbf{Modo de execução declarado:} subagente · \textbf{Permissões:} edit: deny; bash: deny
\end{arquitetura}

\elemento{Recuperador de contexto do repositório (v1.0.0)}{\metainfo{Categoria}{Ferramentas e apoio transversal} \hfill \metainfo{Tipo}{especialista} \hfill \metainfo{Identificador}{context-retriever} \hfill \metainfo{Ficha}{context-retriever.md}}

\begin{descricao}
Localiza arquivos de contexto, padrões e guias relevantes em um repositório.
\end{descricao}

\begin{funcao}
No Core, este perfil apoia tarefas relacionadas a localiza arquivos de contexto, padrões e guias relevantes em um repositório.. O orquestrador pode encaminhar-lhe o trabalho na etapa correspondente; a resposta retorna ao fluxo principal para integração e conferência de fontes, critérios e permissões. A ficha documenta a função, mas não comprova que o perfil tenha sido chamado ou executado a tarefa.
\end{funcao}

\begin{arquitetura}
\textbf{Modo de execução declarado:} subagente · \textbf{Permissões:} edit: deny; bash: deny
\end{arquitetura}

\elemento{Explorador de arquivos de contexto (v1.0.0)}{\metainfo{Categoria}{Ferramentas e apoio transversal} \hfill \metainfo{Tipo}{especialista} \hfill \metainfo{Identificador}{ContextScout} \hfill \metainfo{Ficha}{contextscout.md}}

\begin{descricao}
Localiza e prioriza arquivos de contexto do projeto e aponta quando a documentação externa de uma biblioteca pode ser necessária.
\end{descricao}

\begin{funcao}
No Core, este perfil apoia tarefas relacionadas a localiza e prioriza arquivos de contexto do projeto e aponta quando a documentação externa de uma biblioteca pode ser necessária.. O orquestrador pode encaminhar-lhe o trabalho na etapa correspondente; a resposta retorna ao fluxo principal para integração e conferência de fontes, critérios e permissões. A ficha documenta a função, mas não comprova que o perfil tenha sido chamado ou executado a tarefa.
\end{funcao}

\begin{arquitetura}
\textbf{Modo de execução declarado:} subagente
\end{arquitetura}

\elemento{Gestor de contratos de API (v1.0.0)}{\metainfo{Categoria}{Ferramentas e apoio transversal} \hfill \metainfo{Tipo}{especialista} \hfill \metainfo{Identificador}{contract-manager} \hfill \metainfo{Ficha}{contract-manager.md}}

\begin{descricao}
Gerencia contratos de API para apoiar desenvolvimento paralelo com desenho contract-first e suporte a OpenAPI/Swagger.
\end{descricao}

\begin{funcao}
No Core, este perfil apoia tarefas relacionadas a gerencia contratos de API para apoiar desenvolvimento paralelo com desenho contract-first e suporte a OpenAPI/Swagger.. O orquestrador pode encaminhar-lhe o trabalho na etapa correspondente; a resposta retorna ao fluxo principal para integração e conferência de fontes, critérios e permissões. A ficha documenta a função, mas não comprova que o perfil tenha sido chamado ou executado a tarefa.
\end{funcao}

\begin{arquitetura}
\textbf{Modo de execução declarado:} subagente · \textbf{Permissões:} edit: deny; bash: deny
\end{arquitetura}

\elemento{Core hooks}{\metainfo{Categoria}{Ferramentas e apoio transversal} \hfill \metainfo{Tipo}{integration} \hfill \metainfo{Identificador}{core-hooks} \hfill \metainfo{Ficha}{core-hooks.md}}

\begin{descricao}
Hooks do Core (SPEC-935-R653): engine fail-closed (HookMatcher, eventos Pre/PostToolUse e sessão) + política deny-list auditada para Bash, fiada ao SDK livre.
\end{descricao}

\begin{funcao}
No Core, este perfil apoia tarefas relacionadas a hooks do Core (SPEC-935-R653): engine fail-closed (HookMatcher, eventos Pre/PostToolUse e sessão) + política deny-list auditada para Bash, fiada ao SDK livre.. O orquestrador pode encaminhar-lhe o trabalho na etapa correspondente; a resposta retorna ao fluxo principal para integração e conferência de fontes, critérios e permissões. A ficha documenta a função, mas não comprova que o perfil tenha sido chamado ou executado a tarefa.
\end{funcao}

\begin{arquitetura}
\textbf{Modo de execução declarado:} subagente · \textbf{Permissões:} edit: deny; bash: deny
\end{arquitetura}

\elemento{Dados knowledge hub (v1.0.0)}{\metainfo{Categoria}{Ferramentas e apoio transversal} \hfill \metainfo{Tipo}{especialista} \hfill \metainfo{Identificador}{data-knowledge-hub} \hfill \metainfo{Ficha}{data-knowledge-hub.md}}

\begin{descricao}
ID: data-knowledge-hub Tipo: Subagente local (com fallback online/offline) Fonte: skills/tooling/data\_knowledge\_hub/ Hub unificado de descoberta de dados e conhecimento que integra 16 fontes em 5 domínios, com roteamento automático por tipo de consulta e fallback offline com dados mockados. ; Fonte ; Dados ; API ; ; ---; ---; ---; ; yfinance ; Ações B3 (PETR4, VALE3, ITUB4), índices, câmbio ; yfinance ; ; BCB/SGS ; 2.000+ séries (IPCA, SELIC, PIB, câmbio) ; api.bcb.gov.br/sgs ; ; FRED ; 800K+ séries macroeconômicas (GDP, unemployment, CPI) ; api.stlouisfed.org/fred ; ; World Bank ; Indicadores globais (PIB, população, educação) ; api.worldbank.org/v2 ; ; Alpha Vantage ; Ações, câmbio, cripto, indicadores técnic
\end{descricao}

\begin{funcao}
No Core, este perfil apoia tarefas relacionadas a iD: data-knowledge-hub Tipo: Subagente local (com fallback online/offline) Fonte: skills/tooling/data\_knowledge\_hub/ Hub unificado de descoberta de dados e conhecimento que integra 16 fontes em 5…. O orquestrador pode encaminhar-lhe o trabalho na etapa correspondente; a resposta retorna ao fluxo principal para integração e conferência de fontes, critérios e permissões. A ficha documenta a função, mas não comprova que o perfil tenha sido chamado ou executado a tarefa.
\end{funcao}

\begin{arquitetura}
\textbf{Modo de execução declarado:} subagente · \textbf{Permissões:} edit: deny; bash: deny
\end{arquitetura}

\elemento{Debugger (v1.0.0)}{\metainfo{Categoria}{Ferramentas e apoio transversal} \hfill \metainfo{Tipo}{especialista} \hfill \metainfo{Identificador}{debugger} \hfill \metainfo{Ficha}{debugger.md}}

\begin{descricao}
Investiga e diagnostica bugs com acesso a bash e logs. Capacidades descritas: Diagnostica bugs com acesso a bash
\end{descricao}

\begin{funcao}
No Core, este perfil apoia tarefas relacionadas a investiga e diagnostica bugs com acesso a bash e logs.. O orquestrador pode encaminhar-lhe o trabalho na etapa correspondente; a resposta retorna ao fluxo principal para integração e conferência de fontes, critérios e permissões. A ficha documenta a função, mas não comprova que o perfil tenha sido chamado ou executado a tarefa.
\end{funcao}

\begin{arquitetura}
\textbf{Modo de execução declarado:} subagente · \textbf{Permissões:} edit: deny; bash: deny
\end{arquitetura}

\elemento{Especialista em DevOps (v1.0.0)}{\metainfo{Categoria}{Ferramentas e apoio transversal} \hfill \metainfo{Tipo}{especialista} \hfill \metainfo{Identificador}{devops-specialist} \hfill \metainfo{Ficha}{devops-specialist.md}}

\begin{descricao}
Apoia CI/CD, infraestrutura como código e automação de implantação.
\end{descricao}

\begin{funcao}
No Core, este perfil apoia tarefas relacionadas a apoia CI/CD, infraestrutura como código e automação de implantação.. O orquestrador pode encaminhar-lhe o trabalho na etapa correspondente; a resposta retorna ao fluxo principal para integração e conferência de fontes, critérios e permissões. A ficha documenta a função, mas não comprova que o perfil tenha sido chamado ou executado a tarefa.
\end{funcao}

\begin{arquitetura}
\textbf{Modo de execução declarado:} subagente · \textbf{Permissões:} edit: deny; bash: deny
\end{arquitetura}

\elemento{Redator de docs (v1.0.0)}{\metainfo{Categoria}{Ferramentas e apoio transversal} \hfill \metainfo{Tipo}{especialista} \hfill \metainfo{Identificador}{docs-writer} \hfill \metainfo{Ficha}{docs-writer.md}}

\begin{descricao}
Voce e um escritor tecnico. Crie documentacao clara e abrangente. - Clareza: explicacoes diretas sem jargao - Estrutura: cabecalhos, listas, tabelas, codigo formatado - Exemplos: sempre inclua exemplos de codigo funcionais - Consistencia: mesmo tom e formato em toda a doc 1. README.md: visao geral, instalacao, uso 2. API docs: endpoints, parametros, respostas 3. Guia de contribuicao: setup, padroes, PRs 4. Arquitetura: decisoes de design, diagramas 5. Changelog: versoes, breaking changes
\end{descricao}

\begin{funcao}
No Core, este perfil apoia tarefas relacionadas a voce e um escritor tecnico.. O orquestrador pode encaminhar-lhe o trabalho na etapa correspondente; a resposta retorna ao fluxo principal para integração e conferência de fontes, critérios e permissões. A ficha documenta a função, mas não comprova que o perfil tenha sido chamado ou executado a tarefa.
\end{funcao}

\begin{arquitetura}
\textbf{Modo de execução declarado:} subagente · \textbf{Permissões:} edit: allow; bash: deny
\end{arquitetura}

\elemento{Especialista em documentação (v1.0.0)}{\metainfo{Categoria}{Ferramentas e apoio transversal} \hfill \metainfo{Tipo}{especialista} \hfill \metainfo{Identificador}{documentation} \hfill \metainfo{Ficha}{documentation.md}}

\begin{descricao}
Aplica as convenções documentais do repositório: consulta o contexto do projeto, propõe alterações antes de editar e prioriza textos que possam ser lidos rapidamente.
\end{descricao}

\begin{funcao}
No Core, este perfil apoia tarefas relacionadas a aplica as convenções documentais do repositório: consulta o contexto do projeto, propõe alterações antes de editar e prioriza textos que possam ser lidos rapidamente.. O orquestrador pode encaminhar-lhe o trabalho na etapa correspondente; a resposta retorna ao fluxo principal para integração e conferência de fontes, critérios e permissões. A ficha documenta a função, mas não comprova que o perfil tenha sido chamado ou executado a tarefa.
\end{funcao}

\begin{arquitetura}
\textbf{Modo de execução declarado:} subagente · \textbf{Permissões:} edit: deny; bash: deny
\end{arquitetura}

\elemento{Pesquisador de documentação externa (v1.0.0)}{\metainfo{Categoria}{Ferramentas e apoio transversal} \hfill \metainfo{Tipo}{especialista} \hfill \metainfo{Identificador}{externalscout} \hfill \metainfo{Ficha}{externalscout.md}}

\begin{descricao}
Busca documentação atual e específica por versão de bibliotecas e frameworks usando Context7 e outras fontes; filtra e retorna material pertinente.
\end{descricao}

\begin{funcao}
No Core, este perfil apoia tarefas relacionadas a busca documentação atual e específica por versão de bibliotecas e frameworks usando Context7 e outras fontes. O orquestrador pode encaminhar-lhe o trabalho na etapa correspondente; a resposta retorna ao fluxo principal para integração e conferência de fontes, critérios e permissões. A ficha documenta a função, mas não comprova que o perfil tenha sido chamado ou executado a tarefa.
\end{funcao}

\begin{arquitetura}
\textbf{Modo de execução declarado:} subagente · \textbf{Permissões:} edit: deny; bash: deny
\end{arquitetura}

\elemento{Especialista em interfaces front-end (v1.0.0)}{\metainfo{Categoria}{Ferramentas e apoio transversal} \hfill \metainfo{Tipo}{especialista} \hfill \metainfo{Identificador}{frontend-specialist} \hfill \metainfo{Ficha}{frontend-specialist.md}}

\begin{descricao}
Apoia o desenho de interfaces, sistemas de design, temas e animações.
\end{descricao}

\begin{funcao}
No Core, este perfil apoia tarefas relacionadas a apoia o desenho de interfaces, sistemas de design, temas e animações.. O orquestrador pode encaminhar-lhe o trabalho na etapa correspondente; a resposta retorna ao fluxo principal para integração e conferência de fontes, critérios e permissões. A ficha documenta a função, mas não comprova que o perfil tenha sido chamado ou executado a tarefa.
\end{funcao}

\begin{arquitetura}
\textbf{Modo de execução declarado:} subagente · \textbf{Permissões:} edit: deny; bash: deny
\end{arquitetura}

\elemento{Gametheory local (v1.0.0)}{\metainfo{Categoria}{Ferramentas e apoio transversal} \hfill \metainfo{Tipo}{especialista} \hfill \metainfo{Identificador}{gametheory-local} \hfill \metainfo{Ficha}{gametheory-local.md}}

\begin{descricao}
ID: gametheory-local Tipo: Subagente local (sem LLM) Fonte: skills/tooling/game\_theory\_local.py Cálculo de teoria dos jogos determinístico usando nashpy + numpy. Substitui completamente o módulo debate\_strategies.py que fazia 146 chamadas de reasoning via LLM para calcular equilíbrios de Nash e valores de Shapley. ; Capacidade ; Descrição ; Substitui ; ; ---; ---; ---; ; nash\_equilibrium ; Equilíbrio de Nash para 2 jogadores (suporte enum) ; debate\_strategies.nash\_via\_llm ; ; stackelberg\_equilibrium ; Equilíbrio líder-seguidor ; otimização manual ; ; shapley\_value ; Valor de Shapley justo por agente ; distribuição baseada em LLM ; ; pareto\_frontier ; Fronteira de Pareto multi-objetivo ; trade-off via LLM ; - Na
\end{descricao}

\begin{funcao}
No Core, este perfil apoia tarefas relacionadas a iD: gametheory-local Tipo: Subagente local (sem LLM) Fonte: skills/tooling/game\_theory\_local.py Cálculo de teoria dos jogos determinístico usando nashpy + numpy.. O orquestrador pode encaminhar-lhe o trabalho na etapa correspondente; a resposta retorna ao fluxo principal para integração e conferência de fontes, critérios e permissões. A ficha documenta a função, mas não comprova que o perfil tenha sido chamado ou executado a tarefa.
\end{funcao}

\begin{arquitetura}
\textbf{Modo de execução declarado:} subagente · \textbf{Permissões:} edit: deny; bash: deny
\end{arquitetura}

\elemento{Gemini cli}{\metainfo{Categoria}{Ferramentas e apoio transversal} \hfill \metainfo{Tipo}{integration} \hfill \metainfo{Identificador}{gemini-cli} \hfill \metainfo{Ficha}{gemini-cli.md}}

\begin{descricao}
Executor externo orquestrável: gemini (pacote @google/gemini-cli, repositório google-gemini/gemini-cli). Roda localmente, chama a API Gemini (key própria GEMINI\_API\_KEY ou Vertex AI). - gemini\_available() / gemini\_version() — presença e versão semântica. - gemini\_run(prompt, model, output\_format, timeout) — headless -p com --output-format json; stream-json, sem shell (lista), nunca lança. - doctor\_check() — pass se instalado (com versão), warn se ausente. - install\_instructions() — npm/npx/Homebrew + variáveis de auth. - CLI /gemini status; run; doctor; install com flags -m/--model, --output-format, --timeout e --help. - Execução externa: resultados NÃO são "verificados" sem validação do Core. - Sem GEMINI\_API
\end{descricao}

\begin{funcao}
No Core, este perfil apoia tarefas relacionadas a executor externo orquestrável: gemini (pacote @google/gemini-cli, repositório google-gemini/gemini-cli).. O orquestrador pode encaminhar-lhe o trabalho na etapa correspondente; a resposta retorna ao fluxo principal para integração e conferência de fontes, critérios e permissões. A ficha documenta a função, mas não comprova que o perfil tenha sido chamado ou executado a tarefa.
\end{funcao}

\begin{arquitetura}
\textbf{Modo de execução declarado:} subagente · \textbf{Permissões:} edit: deny; bash: deny
\end{arquitetura}

\elemento{Gemini notebooklm bridge}{\metainfo{Categoria}{Ferramentas e apoio transversal} \hfill \metainfo{Tipo}{integration} \hfill \metainfo{Identificador}{gemini-notebooklm-bridge} \hfill \metainfo{Ficha}{gemini-notebooklm-bridge.md}}

\begin{descricao}
Integra CLI, MCP e skill oficiais do Gemini Notebook à orquestração MarceloClaro, preservando schemas, perfis, estados operacionais e proveniência de artefatos.
\end{descricao}

\begin{funcao}
No Core, este perfil apoia tarefas relacionadas a integra CLI, MCP e skill oficiais do Gemini Notebook à orquestração MarceloClaro, preservando schemas, perfis, estados operacionais e proveniência de artefatos.. O orquestrador pode encaminhar-lhe o trabalho na etapa correspondente; a resposta retorna ao fluxo principal para integração e conferência de fontes, critérios e permissões. A ficha documenta a função, mas não comprova que o perfil tenha sido chamado ou executado a tarefa.
\end{funcao}

\begin{arquitetura}
\textbf{Modo de execução declarado:} subagente · \textbf{Permissões:} edit: deny; bash: deny
\end{arquitetura}

\elemento{Gestor de git (v1.0.0)}{\metainfo{Categoria}{Ferramentas e apoio transversal} \hfill \metainfo{Tipo}{especialista} \hfill \metainfo{Identificador}{git-manager} \hfill \metainfo{Ficha}{git-manager.md}}

\begin{descricao}
Voce e gerente de git. Commits atomicos e bem descritos. - Commits atomicos: uma mudanca logica por commit - Conventional Commits: type(scope): descricao - Tipos: feat, fix, refactor, test, docs, chore, perf, style - Sempre verificar git status e git diff antes de commitar - NUNCA force push em branches compartilhadas - NUNCA commitar secrets, .env, node\_modules - feat(auth): adiciona login com Google OAuth - fix(api): corrige race condition no endpoint /users - refactor(db): extrai logica de query para repository
\end{descricao}

\begin{funcao}
No Core, este perfil apoia tarefas relacionadas a voce e gerente de git.. O orquestrador pode encaminhar-lhe o trabalho na etapa correspondente; a resposta retorna ao fluxo principal para integração e conferência de fontes, critérios e permissões. A ficha documenta a função, mas não comprova que o perfil tenha sido chamado ou executado a tarefa.
\end{funcao}

\begin{arquitetura}
\textbf{Modo de execução declarado:} subagente · \textbf{Permissões:} bash: allow; edit: deny
\end{arquitetura}

\elemento{Goose cli}{\metainfo{Categoria}{Ferramentas e apoio transversal} \hfill \metainfo{Tipo}{especialista} \hfill \metainfo{Identificador}{goose-cli} \hfill \metainfo{Ficha}{goose-cli.md}}

\begin{descricao}
Você é o agente proxy do Goose, agente de IA nativo open source da Agentic
\end{descricao}

\begin{funcao}
No Core, este perfil apoia tarefas relacionadas a você é o agente proxy do Goose, agente de IA nativo open source da Agentic. O orquestrador pode encaminhar-lhe o trabalho na etapa correspondente; a resposta retorna ao fluxo principal para integração e conferência de fontes, critérios e permissões. A ficha documenta a função, mas não comprova que o perfil tenha sido chamado ou executado a tarefa.
\end{funcao}

\begin{arquitetura}
\textbf{Modo de execução declarado:} subagente · \textbf{Permissões:} edit: deny; bash: deny
\end{arquitetura}

\elemento{Honest critic agente (v1.0.0)}{\metainfo{Categoria}{Ferramentas e apoio transversal} \hfill \metainfo{Tipo}{especialista} \hfill \metainfo{Identificador}{honest-critic-agent} \hfill \metainfo{Ficha}{honest-critic-agent.md}}

\begin{descricao}
Crítico antioverclaim que separa cobertura/processo de mérito de qualidade e recusa nota de topo sem validação externa (Honest Evaluation Engine, R142)
\end{descricao}

\begin{funcao}
No Core, este perfil apoia tarefas relacionadas a crítico antioverclaim que separa cobertura/processo de mérito de qualidade e recusa nota de topo sem validação externa (Honest Evaluation Engine, R142). O orquestrador pode encaminhar-lhe o trabalho na etapa correspondente; a resposta retorna ao fluxo principal para integração e conferência de fontes, critérios e permissões. A ficha documenta a função, mas não comprova que o perfil tenha sido chamado ou executado a tarefa.
\end{funcao}

\begin{arquitetura}
\textbf{Modo de execução declarado:} subagente · \textbf{Permissões:} edit: deny; bash: deny
\end{arquitetura}

\elemento{Especialista em imagens (v1.0.0)}{\metainfo{Categoria}{Ferramentas e apoio transversal} \hfill \metainfo{Tipo}{especialista} \hfill \metainfo{Identificador}{Image Specialist} \hfill \metainfo{Ficha}{image-specialist.md}}

\begin{descricao}
Apoia edição e análise de imagens com ferramentas Gemini AI.
\end{descricao}

\begin{funcao}
No Core, este perfil apoia tarefas relacionadas a apoia edição e análise de imagens com ferramentas Gemini AI.. O orquestrador pode encaminhar-lhe o trabalho na etapa correspondente; a resposta retorna ao fluxo principal para integração e conferência de fontes, critérios e permissões. A ficha documenta a função, mas não comprova que o perfil tenha sido chamado ou executado a tarefa.
\end{funcao}

\begin{arquitetura}
\textbf{Modo de execução declarado:} subagente
\end{arquitetura}

\elemento{Jinja2 templates (v1.0.0)}{\metainfo{Categoria}{Ferramentas e apoio transversal} \hfill \metainfo{Tipo}{especialista} \hfill \metainfo{Identificador}{jinja2-templates} \hfill \metainfo{Ficha}{jinja2-templates.md}}

\begin{descricao}
ID: jinja2-templates Tipo: Subagente local (sem LLM) Fonte: skills/tooling/jinja2\_templates/ Motor de templates Jinja2 para geração de documentos, relatórios e notebooks sem chamadas LLM. Substitui f-strings hardcoded e chamadas de LLM para geração de texto estruturado. ; Template ; Uso ; Substitui ; ; ---; ---; ---; ; fichamento.md.j2 ; Resenha crítica de artigo ; research/fichamento.py LLM ; ; quality\_report.md.j2 ; Relatório de qualidade ; nano\_orchestration/quality\_checker.py LLM ; ; manifest.md.j2 ; Manifesto de pesquisa ; research/hub.py LLM ; ; nanoblock.md.j2 ; Bloco de manuscrito acadêmico ; nano\_orchestration/writer.py LLM ; ; debate\_synthesis.md.j2 ; Síntese de debate multi-agente ; gametheory/modera
\end{descricao}

\begin{funcao}
No Core, este perfil apoia tarefas relacionadas a iD: jinja2-templates Tipo: Subagente local (sem LLM) Fonte: skills/tooling/jinja2\_templates/ Motor de templates Jinja2 para geração de documentos, relatórios e notebooks sem chamadas LLM.. O orquestrador pode encaminhar-lhe o trabalho na etapa correspondente; a resposta retorna ao fluxo principal para integração e conferência de fontes, critérios e permissões. A ficha documenta a função, mas não comprova que o perfil tenha sido chamado ou executado a tarefa.
\end{funcao}

\begin{arquitetura}
\textbf{Modo de execução declarado:} subagente · \textbf{Permissões:} edit: deny; bash: deny
\end{arquitetura}

\elemento{Kaggle cli}{\metainfo{Categoria}{Ferramentas e apoio transversal} \hfill \metainfo{Tipo}{integration} \hfill \metainfo{Identificador}{kaggle-cli} \hfill \metainfo{Ficha}{kaggle-cli.md}}

\begin{descricao}
Executor externo orquestrável: kaggle (Kaggle/kaggle-cli) — competições, datasets, kernels, models, quota via passthrough com auth do operador.
\end{descricao}

\begin{funcao}
No Core, este perfil apoia tarefas relacionadas a executor externo orquestrável: kaggle (Kaggle/kaggle-cli) — competições, datasets, kernels, models, quota via passthrough com auth do operador.. O orquestrador pode encaminhar-lhe o trabalho na etapa correspondente; a resposta retorna ao fluxo principal para integração e conferência de fontes, critérios e permissões. A ficha documenta a função, mas não comprova que o perfil tenha sido chamado ou executado a tarefa.
\end{funcao}

\begin{arquitetura}
\textbf{Modo de execução declarado:} subagente · \textbf{Permissões:} edit: deny; bash: deny
\end{arquitetura}

\elemento{Kdp Cover engenheiro Ph D (v1.0.0)}{\metainfo{Categoria}{Ferramentas e apoio transversal} \hfill \metainfo{Tipo}{especialista} \hfill \metainfo{Identificador}{KDP Cover Engineer PhD} \hfill \metainfo{Ficha}{kdp-cover-engineer-phd.md}}

\begin{descricao}
Especialista PhD Amazon KDP em capa completa, contracapa, lombada, wrap, bleed, template, barcode e PDF de capa.
\end{descricao}

\begin{funcao}
No Core, este perfil apoia tarefas relacionadas a especialista PhD Amazon KDP em capa completa, contracapa, lombada, wrap, bleed, template, barcode e PDF de capa.. O orquestrador pode encaminhar-lhe o trabalho na etapa correspondente; a resposta retorna ao fluxo principal para integração e conferência de fontes, critérios e permissões. A ficha documenta a função, mas não comprova que o perfil tenha sido chamado ou executado a tarefa.
\end{funcao}

\begin{arquitetura}
\textbf{Modo de execução declarado:} subagente
\end{arquitetura}

\elemento{Kdp e Book e Pub Ph D (v1.0.0)}{\metainfo{Categoria}{Ferramentas e apoio transversal} \hfill \metainfo{Tipo}{especialista} \hfill \metainfo{Identificador}{KDP eBook ePub PhD} \hfill \metainfo{Ficha}{kdp-ebook-epub-phd.md}}

\begin{descricao}
Especialista PhD Amazon KDP em ePub, Kindle, KPF, sumário navegável, metadados digitais e conversão LaTeX/Markdown. Capacidades descritas: Converte LaTeX/Markdown para ePub válido com sumário navegável.; Gera KPF (Kindle Package Format) compatível com Amazon.
\end{descricao}

\begin{funcao}
No Core, este perfil apoia tarefas relacionadas a especialista PhD Amazon KDP em ePub, Kindle, KPF, sumário navegável, metadados digitais e conversão LaTeX/Markdown.. O orquestrador pode encaminhar-lhe o trabalho na etapa correspondente; a resposta retorna ao fluxo principal para integração e conferência de fontes, critérios e permissões. A ficha documenta a função, mas não comprova que o perfil tenha sido chamado ou executado a tarefa.
\end{funcao}

\begin{arquitetura}
\textbf{Modo de execução declarado:} subagente
\end{arquitetura}

\elemento{Kdp Final QA Ph D (v1.0.0)}{\metainfo{Categoria}{Ferramentas e apoio transversal} \hfill \metainfo{Tipo}{especialista} \hfill \metainfo{Identificador}{KDP Final QA PhD} \hfill \metainfo{Ficha}{kdp-final-qa-phd.md}}

\begin{descricao}
Gate PhD Amazon KDP de QA final para pacote de upload, checklist, evidências, riscos residuais e instruções finais.
\end{descricao}

\begin{funcao}
No Core, este perfil apoia tarefas relacionadas a gate PhD Amazon KDP de QA final para pacote de upload, checklist, evidências, riscos residuais e instruções finais.. O orquestrador pode encaminhar-lhe o trabalho na etapa correspondente; a resposta retorna ao fluxo principal para integração e conferência de fontes, critérios e permissões. A ficha documenta a função, mas não comprova que o perfil tenha sido chamado ou executado a tarefa.
\end{funcao}

\begin{arquitetura}
\textbf{Modo de execução declarado:} subagente
\end{arquitetura}

\elemento{Kdp Interior Layout Ph D (v1.0.0)}{\metainfo{Categoria}{Ferramentas e apoio transversal} \hfill \metainfo{Tipo}{especialista} \hfill \metainfo{Identificador}{KDP Interior Layout PhD} \hfill \metainfo{Ficha}{kdp-interior-layout-phd.md}}

\begin{descricao}
Especialista PhD Amazon KDP em miolo, trim size, margens internas/externas, sangria, LaTeX e PDF pronto para impressão.
\end{descricao}

\begin{funcao}
No Core, este perfil apoia tarefas relacionadas a especialista PhD Amazon KDP em miolo, trim size, margens internas/externas, sangria, LaTeX e PDF pronto para impressão.. O orquestrador pode encaminhar-lhe o trabalho na etapa correspondente; a resposta retorna ao fluxo principal para integração e conferência de fontes, critérios e permissões. A ficha documenta a função, mas não comprova que o perfil tenha sido chamado ou executado a tarefa.
\end{funcao}

\begin{arquitetura}
\textbf{Modo de execução declarado:} subagente
\end{arquitetura}

\elemento{Kdp Metadata \& ISBN Ph D (v1.0.0)}{\metainfo{Categoria}{Ferramentas e apoio transversal} \hfill \metainfo{Tipo}{especialista} \hfill \metainfo{Identificador}{KDP Metadata \& ISBN PhD} \hfill \metainfo{Ficha}{kdp-metadata-isbn-phd.md}}

\begin{descricao}
Especialista PhD Amazon KDP em ISBN, copyright, ficha catalográfica, metadados bibliográficos e consistência editorial. Capacidades descritas: Gerencia atribuição e consistência de ISBN entre formatos (impresso, digital, capa).; Produz ficha catalográfica CIP e metadados bibliográficos conforme legislação brasileira.
\end{descricao}

\begin{funcao}
No Core, este perfil apoia tarefas relacionadas a especialista PhD Amazon KDP em ISBN, copyright, ficha catalográfica, metadados bibliográficos e consistência editorial.. O orquestrador pode encaminhar-lhe o trabalho na etapa correspondente; a resposta retorna ao fluxo principal para integração e conferência de fontes, critérios e permissões. A ficha documenta a função, mas não comprova que o perfil tenha sido chamado ou executado a tarefa.
\end{funcao}

\begin{arquitetura}
\textbf{Modo de execução declarado:} subagente
\end{arquitetura}

\elemento{Kdp orquestrador Ph D (v1.0.0)}{\metainfo{Categoria}{Ferramentas e apoio transversal} \hfill \metainfo{Tipo}{especialista} \hfill \metainfo{Identificador}{KDP Orchestrator PhD} \hfill \metainfo{Ficha}{kdp-orchestrator-phd.md}}

\begin{descricao}
Orquestrador PhD Amazon KDP para coordenar miolo, capa, ePub, metadados, preflight e QA final de livros físicos e digitais.
\end{descricao}

\begin{funcao}
No Core, este perfil apoia tarefas relacionadas a orquestrador PhD Amazon KDP para coordenar miolo, capa, ePub, metadados, preflight e QA final de livros físicos e digitais.. O orquestrador pode encaminhar-lhe o trabalho na etapa correspondente; a resposta retorna ao fluxo principal para integração e conferência de fontes, critérios e permissões. A ficha documenta a função, mas não comprova que o perfil tenha sido chamado ou executado a tarefa.
\end{funcao}

\begin{arquitetura}
\textbf{Modo de execução declarado:} subagente
\end{arquitetura}

\elemento{Kdp Preflight auditor Ph D (v1.0.0)}{\metainfo{Categoria}{Ferramentas e apoio transversal} \hfill \metainfo{Tipo}{especialista} \hfill \metainfo{Identificador}{KDP Preflight Auditor PhD} \hfill \metainfo{Ficha}{kdp-preflight-auditor-phd.md}}

\begin{descricao}
Auditor PhD Amazon KDP de preflight PDF para MediaBox, CropBox, fontes, imagens, hyperlinks, anotações e texto fora das margens.
\end{descricao}

\begin{funcao}
No Core, este perfil apoia tarefas relacionadas a auditor PhD Amazon KDP de preflight PDF para MediaBox, CropBox, fontes, imagens, hyperlinks, anotações e texto fora das margens.. O orquestrador pode encaminhar-lhe o trabalho na etapa correspondente; a resposta retorna ao fluxo principal para integração e conferência de fontes, critérios e permissões. A ficha documenta a função, mas não comprova que o perfil tenha sido chamado ou executado a tarefa.
\end{funcao}

\begin{arquitetura}
\textbf{Modo de execução declarado:} subagente
\end{arquitetura}

\elemento{Linguistic corrector (v1.0.0)}{\metainfo{Categoria}{Ferramentas e apoio transversal} \hfill \metainfo{Tipo}{especialista} \hfill \metainfo{Identificador}{linguistic-corrector} \hfill \metainfo{Ficha}{linguistic-corrector.md}}

\begin{descricao}
Corretor ortografico, gramatical e linguistico especializado em detectar e remover contaminacao de caracteres chineses em saidas PT-BR. Ultima barreira antes da entrega ao usuario. Garantir que TODO texto entregue ao usuario esteja 100\% em portugues brasileiro formal, sem residuos de caracteres chineses, japoneses ou coreanos provenientes do contexto em chines do ecossistema. - Escanear todo texto de saida em busca de caracteres Unicode CJK - Blocos CJK monitorados: Unified Ideographs, Hiragana, Katakana, Hangul, CJK Punctuation - Identificar posicao exata (linha, coluna) de cada caractere contaminante - Classificar por tipo: chinese, japanese, korean, cjk\_punctuation - Texto Markdown: Preservar blocos de codig
\end{descricao}

\begin{funcao}
No Core, este perfil apoia tarefas relacionadas a corretor ortografico, gramatical e linguistico especializado em detectar e remover contaminacao de caracteres chineses em saidas PT-BR.. O orquestrador pode encaminhar-lhe o trabalho na etapa correspondente; a resposta retorna ao fluxo principal para integração e conferência de fontes, critérios e permissões. A ficha documenta a função, mas não comprova que o perfil tenha sido chamado ou executado a tarefa.
\end{funcao}

\begin{arquitetura}
\textbf{Modo de execução declarado:} subagente · \textbf{Permissões:} edit: deny; bash: deny
\end{arquitetura}

\elemento{Literário imagem Sepia (v1.0.0)}{\metainfo{Categoria}{Ferramentas e apoio transversal} \hfill \metainfo{Tipo}{especialista} \hfill \metainfo{Identificador}{Literary Image Sepia} \hfill \metainfo{Ficha}{literary-image-sepia.md}}

\begin{descricao}
Você é o Agente de Imagens Sépia do ecossistema OpenCode, especializado em processamento de imagens para o romance Molambudos — O Diário do Paciente 1.260. Sua função é traduzir visualmente as fotografias citadas na narrativa — descritas como "antigas" e "amareladas" — aplicando um pipeline de envelhecimento que inclui filtro sépia, vinheta e grão de papel. - Texto extraído do capítulo/fragmento do Molambudos - Metadados: personagem, ano, local, contexto narrativo - Busca: Antigravity web search por imagem similar à descrição - Geração: Antigravity image generation para cenas históricas específicas - Fallback: imagem sintética com silhueta + texto Imagem original 1. Filtro Sépia 3×3 R = 0.393R + 0.769G + 0.189B
\end{descricao}

\begin{funcao}
No Core, este perfil apoia tarefas relacionadas a você é o Agente de Imagens Sépia do ecossistema OpenCode, especializado em processamento de imagens para o romance Molambudos — O Diário do Paciente…. O orquestrador pode encaminhar-lhe o trabalho na etapa correspondente; a resposta retorna ao fluxo principal para integração e conferência de fontes, critérios e permissões. A ficha documenta a função, mas não comprova que o perfil tenha sido chamado ou executado a tarefa.
\end{funcao}

\begin{arquitetura}
\textbf{Modo de execução declarado:} subagente
\end{arquitetura}

\elemento{Literário Neurolinguistic Engineering Ph D (v1.0.0)}{\metainfo{Categoria}{Ferramentas e apoio transversal} \hfill \metainfo{Tipo}{especialista} \hfill \metainfo{Identificador}{Literary Neurolinguistic Engineering PhD} \hfill \metainfo{Ficha}{literary-neurolinguistic-engineering-phd.md}}

\begin{descricao}
Especialista em engenharia neurolinguística literária — aplica padrões de hipnose ericksoniana, sugestão indireta e ancoragem rítmica na prosa literária.
\end{descricao}

\begin{funcao}
No Core, este perfil apoia tarefas relacionadas a especialista em engenharia neurolinguística literária — aplica padrões de hipnose ericksoniana, sugestão indireta e ancoragem rítmica na prosa literária.. O orquestrador pode encaminhar-lhe o trabalho na etapa correspondente; a resposta retorna ao fluxo principal para integração e conferência de fontes, critérios e permissões. A ficha documenta a função, mas não comprova que o perfil tenha sido chamado ou executado a tarefa.
\end{funcao}

\begin{arquitetura}
\textbf{Modo de execução declarado:} subagente
\end{arquitetura}

\elemento{Especialista em redução de chamadas a modelos (v1.0.0)}{\metainfo{Categoria}{Ferramentas e apoio transversal} \hfill \metainfo{Tipo}{especialista} \hfill \metainfo{Identificador}{llm-reduction} \hfill \metainfo{Ficha}{llm-reduction.md}}

\begin{descricao}
Camada de apoio que substitui chamadas a modelos de linguagem por busca textual, roteamento baseado em regras, classificação local e estratégias de teoria dos jogos; o objetivo declarado é reduzir custo, latência e dependência de rede.
\end{descricao}

\begin{funcao}
No Core, este perfil apoia tarefas relacionadas a camada de apoio que substitui chamadas a modelos de linguagem por busca textual, roteamento baseado em regras, classificação local e estratégias de teoria dos…. O orquestrador pode encaminhar-lhe o trabalho na etapa correspondente; a resposta retorna ao fluxo principal para integração e conferência de fontes, critérios e permissões. A ficha documenta a função, mas não comprova que o perfil tenha sido chamado ou executado a tarefa.
\end{funcao}

\begin{arquitetura}
\textbf{Modo de execução declarado:} subagente · \textbf{Permissões:} edit: deny; bash: deny
\end{arquitetura}

\elemento{Marceloclaro}{\metainfo{Categoria}{Ferramentas e apoio transversal} \hfill \metainfo{Tipo}{especialista} \hfill \metainfo{Identificador}{marceloclaro} \hfill \metainfo{Ficha}{marceloclaro.md}}

\begin{descricao}
Avatar de Marcelo Claro: Controle Supremo, Criador e Orquestrador Central de todo o OpenCode e OpenCode Ecosystem.
\end{descricao}

\begin{funcao}
No Core, este perfil apoia tarefas relacionadas a avatar de Marcelo Claro: Controle Supremo, Criador e Orquestrador Central de todo o OpenCode e OpenCode Ecosystem.. O orquestrador pode encaminhar-lhe o trabalho na etapa correspondente; a resposta retorna ao fluxo principal para integração e conferência de fontes, critérios e permissões. A ficha documenta a função, mas não comprova que o perfil tenha sido chamado ou executado a tarefa.
\end{funcao}

\begin{arquitetura}
\textbf{Modo de execução declarado:} principal · \textbf{Permissões:} edit: ask; bash: ask; task: allow
\end{arquitetura}

\elemento{Orquestrador de Master}{\metainfo{Categoria}{Ferramentas e apoio transversal} \hfill \metainfo{Tipo}{especialista} \hfill \metainfo{Identificador}{MasterOrchestrator} \hfill \metainfo{Ficha}{master-orchestrator.md}}

\begin{descricao}
Orquestrador mestre e controlador de ciclo de vida (end-to-end) para o Ecossistema OpenCode e Polimata. Inicializa e finaliza as demandas do usuário, garantindo auditoria e reprodutibilidade de longo prazo.
\end{descricao}

\begin{funcao}
No Core, este perfil apoia tarefas relacionadas a orquestrador mestre e controlador de ciclo de vida (end-to-end) para o Ecossistema OpenCode e Polimata.. O orquestrador pode encaminhar-lhe o trabalho na etapa correspondente; a resposta retorna ao fluxo principal para integração e conferência de fontes, critérios e permissões. A ficha documenta a função, mas não comprova que o perfil tenha sido chamado ou executado a tarefa.
\end{funcao}

\begin{arquitetura}
\textbf{Modo de execução declarado:} subagente
\end{arquitetura}

\elemento{Minizinc mcp}{\metainfo{Categoria}{Ferramentas e apoio transversal} \hfill \metainfo{Tipo}{integration} \hfill \metainfo{Identificador}{minizinc-mcp} \hfill \metainfo{Ficha}{minizinc-mcp.md}}

\begin{descricao}
Ponte orquestrável do MiniZinc Constraint Solver MCP Server (r33drichards/minizinc-mcp, MIT): modelagem CSP/COP via tool solve\_constraint (solver gecode). Monta payloads, checa solver + deps e emite config stdio/SSE.
\end{descricao}

\begin{funcao}
No Core, este perfil apoia tarefas relacionadas a ponte orquestrável do MiniZinc Constraint Solver MCP Server (r33drichards/minizinc-mcp, MIT): modelagem CSP/COP via tool solve\_constraint (solver gecode).. O orquestrador pode encaminhar-lhe o trabalho na etapa correspondente; a resposta retorna ao fluxo principal para integração e conferência de fontes, critérios e permissões. A ficha documenta a função, mas não comprova que o perfil tenha sido chamado ou executado a tarefa.
\end{funcao}

\begin{arquitetura}
\textbf{Modo de execução declarado:} subagente · \textbf{Permissões:} edit: deny; bash: deny
\end{arquitetura}

\elemento{Médico Cardiologista (v1.0.0)}{\metainfo{Categoria}{Ferramentas e apoio transversal} \hfill \metainfo{Tipo}{especialista} \hfill \metainfo{Identificador}{Médico Cardiologista} \hfill \metainfo{Ficha}{medico-cardiologista.md}}

\begin{descricao}
--- name: Médico Cardiologista — Especialista em Cardiologia description: Agente especializado Médico Cardiologista — Especialista em Cardiologia version: '1.0.0' skills: - id: sindromes-coronarianas-
\end{descricao}

\begin{funcao}
No Core, este perfil apoia tarefas relacionadas a --- name: Médico Cardiologista — Especialista em Cardiologia description: Agente especializado Médico Cardiologista — Especialista em Cardiologia version: '1.0.0' skills: - id: sindromes-coronarianas-. O orquestrador pode encaminhar-lhe o trabalho na etapa correspondente; a resposta retorna ao fluxo principal para integração e conferência de fontes, critérios e permissões. A ficha documenta a função, mas não comprova que o perfil tenha sido chamado ou executado a tarefa.
\end{funcao}

\begin{arquitetura}
\textbf{Modo de execução declarado:} subagente
\end{arquitetura}

\elemento{Médico Clínico Geral — Orquestrador Clínico (v1.0.0)}{\metainfo{Categoria}{Ferramentas e apoio transversal} \hfill \metainfo{Tipo}{especialista} \hfill \metainfo{Identificador}{Médico Clínico Geral — Orquestrador Clínico} \hfill \metainfo{Ficha}{medico-clinico-geral.md}}

\begin{descricao}
Você é o Médico Clínico Geral, orquestrador do ecossistema Médico Virtual Supremo. Sua função é realizar a triagem inicial do paciente, consolidar as análises dos especialistas cardiologista, neurologista, radiologista e infectologista, detectar conflitos entre recomendações e produzir um plano integrado e priorizado para revisão humana. - Classificar urgência (emergência, urgência, eletivo) - Identificar síndrome predominante - Coletar dados mínimos necessários - Acionar especialistas conforme necessidade - Consolidar hipóteses de todos os especialistas consultados - Detectar conflitos e divergências - Priorizar diagnósticos por gravidade e urgência - Produzir plano coerente e factível - Ações imediatas (basea
\end{descricao}

\begin{funcao}
No Core, este perfil apoia tarefas relacionadas a você é o Médico Clínico Geral, orquestrador do ecossistema Médico Virtual Supremo.. O orquestrador pode encaminhar-lhe o trabalho na etapa correspondente; a resposta retorna ao fluxo principal para integração e conferência de fontes, critérios e permissões. A ficha documenta a função, mas não comprova que o perfil tenha sido chamado ou executado a tarefa.
\end{funcao}

\begin{arquitetura}
\textbf{Modo de execução declarado:} subagente
\end{arquitetura}

\elemento{Médico Infectologista (v1.0.0)}{\metainfo{Categoria}{Ferramentas e apoio transversal} \hfill \metainfo{Tipo}{especialista} \hfill \metainfo{Identificador}{Médico Infectologista} \hfill \metainfo{Ficha}{medico-infectologista.md}}

\begin{descricao}
--- name: Médico Infectologista — Especialista em Infectologia description: Agente especializado Médico Infectologista — Especialista em Infectologia version: '1.0.0' skills: - id: sepse-qsofa-sofa-pr
\end{descricao}

\begin{funcao}
No Core, este perfil apoia tarefas relacionadas a --- name: Médico Infectologista — Especialista em Infectologia description: Agente especializado Médico Infectologista — Especialista em Infectologia version: '1.0.0' skills: - id: sepse-qsofa-sofa-pr. O orquestrador pode encaminhar-lhe o trabalho na etapa correspondente; a resposta retorna ao fluxo principal para integração e conferência de fontes, critérios e permissões. A ficha documenta a função, mas não comprova que o perfil tenha sido chamado ou executado a tarefa.
\end{funcao}

\begin{arquitetura}
\textbf{Modo de execução declarado:} subagente
\end{arquitetura}

\elemento{Médico Neurologista (v1.0.0)}{\metainfo{Categoria}{Ferramentas e apoio transversal} \hfill \metainfo{Tipo}{especialista} \hfill \metainfo{Identificador}{Médico Neurologista} \hfill \metainfo{Ficha}{medico-neurologista.md}}

\begin{descricao}
Você é o Médico Neurologista, especialista em neurologia do ecossistema Médico Virtual Supremo. Sua função é analisar casos com manifestações neurológicas, interpretar exames de neuroimagem e eletroencefalograma, e propor hipóteses e planos para revisão humana. - AVC/AVE: Isquêmico, hemorrágico, AIT, janela terapêutica - Cefaleias: Enxaqueca, cefaleia tensional, cefaleia em salvas - Epilepsia: Tipos de crise, status epilepticus, ajuste de anticonvulsivantes - Demências: Alzheimer, demência vascular, Lewy, frontotemporal - Distúrbios do movimento: Parkinson, tremor essencial, distonia - Neuropatias: Periférica, polineuropatia, síndrome do túnel do carpo 1. ABC do AVC: Hora de início, NIHSS, janela para trombólis
\end{descricao}

\begin{funcao}
No Core, este perfil apoia tarefas relacionadas a você é o Médico Neurologista, especialista em neurologia do ecossistema Médico Virtual Supremo.. O orquestrador pode encaminhar-lhe o trabalho na etapa correspondente; a resposta retorna ao fluxo principal para integração e conferência de fontes, critérios e permissões. A ficha documenta a função, mas não comprova que o perfil tenha sido chamado ou executado a tarefa.
\end{funcao}

\begin{arquitetura}
\textbf{Modo de execução declarado:} subagente
\end{arquitetura}

\elemento{Médico Radiologista (v1.0.0)}{\metainfo{Categoria}{Ferramentas e apoio transversal} \hfill \metainfo{Tipo}{especialista} \hfill \metainfo{Identificador}{Médico Radiologista} \hfill \metainfo{Ficha}{medico-radiologista.md}}

\begin{descricao}
Você é o Médico Radiologista, especialista em radiologia e diagnóstico por imagem do ecossistema Médico Virtual Supremo. Sua função é interpretar laudos de exames de imagem, correlacionar achados radiológicos com o quadro clínico e identificar achados incidentais relevantes. - Radiografia: Tórax (pneumonia, ICC, derrame), abdome, esqueleto - Tomografia: Crânio (AVC, hemorragia), tórax (TEP, massas), abdome - Ressonância: SNC (tumores, esclerose múltipla), musculoesquelética - Ultrassom: Abdome total, tireoide, partes moles, Doppler vascular - Mamografia: BI-RADS, screening, achados suspeitos - Achados incidentais: Classificação, recomendação de seguimento 1. Descreva achados usando terminologia padronizada (BI-
\end{descricao}

\begin{funcao}
No Core, este perfil apoia tarefas relacionadas a você é o Médico Radiologista, especialista em radiologia e diagnóstico por imagem do ecossistema Médico Virtual Supremo.. O orquestrador pode encaminhar-lhe o trabalho na etapa correspondente; a resposta retorna ao fluxo principal para integração e conferência de fontes, critérios e permissões. A ficha documenta a função, mas não comprova que o perfil tenha sido chamado ou executado a tarefa.
\end{funcao}

\begin{arquitetura}
\textbf{Modo de execução declarado:} subagente
\end{arquitetura}

\elemento{Médico Virtual Supremo — Apoio Clínico Auditável (v1.0.0)}{\metainfo{Categoria}{Ferramentas e apoio transversal} \hfill \metainfo{Tipo}{especialista} \hfill \metainfo{Identificador}{Médico Virtual Supremo — Apoio Clínico Auditável} \hfill \metainfo{Ficha}{medico-virtual-supremo.md}}

\begin{descricao}
Você é o Médico Virtual Supremo, um sistema de apoio à decisão clínica auditável. Sua função é organizar dados clínicos, identificar inconsistências, sintetizar evidências, propor hipóteses diferenciais e apoiar a elaboração de planos clínicos para revisão humana obrigatória. Você NÃO substitui consulta, exame físico, julgamento profissional, protocolos institucionais, farmacêutico, laboratório, radiologista, especialista ou serviço de emergência. Não deve operar como prescritor autônomo, emitir diagnóstico definitivo sem validação profissional nem alterar tratamento diretamente. ; Modo ; Público-Alvo ; Comportamento ; ; :-----; :-------------; :--------------; ; professional\_cds ; Profissional habilitado ; Hip
\end{descricao}

\begin{funcao}
No Core, este perfil apoia tarefas relacionadas a você é o Médico Virtual Supremo, um sistema de apoio à decisão clínica auditável.. O orquestrador pode encaminhar-lhe o trabalho na etapa correspondente; a resposta retorna ao fluxo principal para integração e conferência de fontes, critérios e permissões. A ficha documenta a função, mas não comprova que o perfil tenha sido chamado ou executado a tarefa.
\end{funcao}

\begin{arquitetura}
\textbf{Modo de execução declarado:} subagente
\end{arquitetura}

\elemento{Orquestrador de Nano (v1.0.0)}{\metainfo{Categoria}{Ferramentas e apoio transversal} \hfill \metainfo{Tipo}{especialista} \hfill \metainfo{Identificador}{Nano Orchestrator} \hfill \metainfo{Ficha}{nano-orchestrator.md}}

\begin{descricao}
Orquestrador de Nano-Manuscritos para produção de documentos acadêmicos de 30–500 laudas usando modelos LiteRT-LM on-device. - Divide o manuscrito em nanoblocks de \~{}10 por página - Cria grafo de dependências entre nanoblocks - Estima tokens por nanoblock (limite: 20K ctx window) - Cada nanoblock recebe uma especificação SDD formal - Critérios de aceitação explícitos por nanoblock - Gerencia rotação de até 20K tokens de contexto - Preserva nanoblocks adjacentes para coerência local - Pool de escritores LiteRT-LM (Qwen3 0.6B, Gemma4 2B/4B) - Produção paralela de nanoblocks independentes - Valida nanoblocks contra SDD - Rejeita e solicita regravação se abaixo do limiar - 3 passadas de fusão: local → transição entr
\end{descricao}

\begin{funcao}
No Core, este perfil apoia tarefas relacionadas a orquestrador de Nano-Manuscritos para produção de documentos acadêmicos de 30–500 laudas usando modelos LiteRT-LM on-device.. O orquestrador pode encaminhar-lhe o trabalho na etapa correspondente; a resposta retorna ao fluxo principal para integração e conferência de fontes, critérios e permissões. A ficha documenta a função, mas não comprova que o perfil tenha sido chamado ou executado a tarefa.
\end{funcao}

\begin{arquitetura}
\textbf{Modo de execução declarado:} subagente
\end{arquitetura}

\elemento{Agente geral OpenAgent (v1.0.0)}{\metainfo{Categoria}{Ferramentas e apoio transversal} \hfill \metainfo{Tipo}{especialista} \hfill \metainfo{Identificador}{openagent} \hfill \metainfo{Ficha}{openagent.md}}

\begin{descricao}
Agente de uso geral para responder consultas, executar tarefas e coordenar fluxos em diferentes domínios, conforme a ficha.
\end{descricao}

\begin{funcao}
No Core, este perfil apoia tarefas relacionadas a agente de uso geral para responder consultas, executar tarefas e coordenar fluxos em diferentes domínios, conforme a ficha.. O orquestrador pode encaminhar-lhe o trabalho na etapa correspondente; a resposta retorna ao fluxo principal para integração e conferência de fontes, critérios e permissões. A ficha documenta a função, mas não comprova que o perfil tenha sido chamado ou executado a tarefa.
\end{funcao}

\begin{arquitetura}
\textbf{Modo de execução declarado:} subagente · \textbf{Permissões:} edit: deny; bash: deny
\end{arquitetura}

\elemento{Opencode agente sdk}{\metainfo{Categoria}{Ferramentas e apoio transversal} \hfill \metainfo{Tipo}{integration} \hfill \metainfo{Identificador}{opencode-agent-sdk} \hfill \metainfo{Ficha}{opencode-agent-sdk.md}}

\begin{descricao}
SDK de agente FREE local-first (SPEC-935-R649, R\$ 0,00): query() com loop agêntico, @tool in-process, hooks e autodetect LiteRT-LM/Ollama/Colibri sobre HTTP OpenAI-compatível. Sem conta e sem cobrança.
\end{descricao}

\begin{funcao}
No Core, este perfil apoia tarefas relacionadas a sDK de agente FREE local-first (SPEC-935-R649, R\$ 0,00): query() com loop agêntico, @tool in-process, hooks e autodetect LiteRT-LM/Ollama/Colibri sobre HTTP OpenAI-compatível.. O orquestrador pode encaminhar-lhe o trabalho na etapa correspondente; a resposta retorna ao fluxo principal para integração e conferência de fontes, critérios e permissões. A ficha documenta a função, mas não comprova que o perfil tenha sido chamado ou executado a tarefa.
\end{funcao}

\begin{arquitetura}
\textbf{Modo de execução declarado:} subagente · \textbf{Permissões:} edit: deny; bash: deny
\end{arquitetura}

\elemento{Opencode go agente (v1.0.0)}{\metainfo{Categoria}{Ferramentas e apoio transversal} \hfill \metainfo{Tipo}{especialista} \hfill \metainfo{Identificador}{opencode-go-agent} \hfill \metainfo{Ficha}{opencode-go-agent.md}}

\begin{descricao}
Você é o OpenCode Go Agent, especialista nos modelos open-source de alta performance disponibilizados pelo gateway OpenCode Go do ecossistema. Você acessa: - Kimi (Moonshot AI) — excelente para código com contexto longo (200k tokens) - DeepSeek — raciocínio profundo e matemática avançada - GLM (Zhipu AI) — bilíngue PT/ZH, raciocínio estruturado - Qwen (Alibaba) — acadêmico e coding, thinking habilitado - MiMo/MiniMax — agentic e multimodal TODA tarefa deve seguir o ciclo: ESPECIFICAR → RED → GREEN → REFACTOR → VERIFICAR 1. ESPECIFICAR: Criar spec com spec\_registry.create\_task\_spec() 2. RED: Confirmar que critérios existem, nenhuma entrega ainda 3. GREEN: Produzir entrega mínima que satisfaz todos os critérios 4
\end{descricao}

\begin{funcao}
No Core, este perfil apoia tarefas relacionadas a você é o OpenCode Go Agent, especialista nos modelos open-source de alta performance disponibilizados pelo gateway OpenCode Go do ecossistema.. O orquestrador pode encaminhar-lhe o trabalho na etapa correspondente; a resposta retorna ao fluxo principal para integração e conferência de fontes, critérios e permissões. A ficha documenta a função, mas não comprova que o perfil tenha sido chamado ou executado a tarefa.
\end{funcao}

\begin{arquitetura}
\textbf{Modo de execução declarado:} subagente · \textbf{Permissões:} edit: deny; bash: deny
\end{arquitetura}

\elemento{Opencode zen agente (v1.0.0)}{\metainfo{Categoria}{Ferramentas e apoio transversal} \hfill \metainfo{Tipo}{especialista} \hfill \metainfo{Identificador}{opencode-zen-agent} \hfill \metainfo{Ficha}{opencode-zen-agent.md}}

\begin{descricao}
Você é o OpenCode Zen Agent, especialista nos modelos premium curados pelo gateway OpenCode Zen — um único endpoint pay-as-you-go com os melhores modelos do mercado: - GPT-5.5 / GPT-5 (OpenAI) — frontier reasoning e multimodal - Claude Opus 4 / Sonnet 4.6 / Haiku 4 (Anthropic) — escrita, código e raciocínio - Gemini 2.5 Pro / Flash (Google) — contexto de 2M tokens, multimodal - DeepSeek R2 (via Zen) — matemática e raciocínio, gratuito - Grok 4 (xAI) — raciocínio em tempo real - Mistral Large 3 — europeu, multilíngue TODA tarefa deve seguir o ciclo: ESPECIFICAR → RED → GREEN → REFACTOR → VERIFICAR 1. ESPECIFICAR: Criar spec com spec\_registry.create\_task\_spec() 2. RED: Confirmar critérios, nenhuma entrega ainda 3
\end{descricao}

\begin{funcao}
No Core, este perfil apoia tarefas relacionadas a você é o OpenCode Zen Agent, especialista nos modelos premium curados pelo gateway OpenCode Zen — um único endpoint pay-as-you-go com os melhores modelos…. O orquestrador pode encaminhar-lhe o trabalho na etapa correspondente; a resposta retorna ao fluxo principal para integração e conferência de fontes, critérios e permissões. A ficha documenta a função, mas não comprova que o perfil tenha sido chamado ou executado a tarefa.
\end{funcao}

\begin{arquitetura}
\textbf{Modo de execução declarado:} subagente · \textbf{Permissões:} edit: deny; bash: deny
\end{arquitetura}

\elemento{Agente de programação OpenCoder (v1.0.0)}{\metainfo{Categoria}{Ferramentas e apoio transversal} \hfill \metainfo{Tipo}{especialista} \hfill \metainfo{Identificador}{opencoder} \hfill \metainfo{Ficha}{opencoder.md}}

\begin{descricao}
Orquestra tarefas complexas de programação, arquitetura e refatoração de vários arquivos.
\end{descricao}

\begin{funcao}
No Core, este perfil apoia tarefas relacionadas a orquestra tarefas complexas de programação, arquitetura e refatoração de vários arquivos.. O orquestrador pode encaminhar-lhe o trabalho na etapa correspondente; a resposta retorna ao fluxo principal para integração e conferência de fontes, critérios e permissões. A ficha documenta a função, mas não comprova que o perfil tenha sido chamado ou executado a tarefa.
\end{funcao}

\begin{arquitetura}
\textbf{Modo de execução declarado:} subagente · \textbf{Permissões:} edit: deny; bash: deny
\end{arquitetura}

\elemento{Redator publicitário (v1.0.0)}{\metainfo{Categoria}{Ferramentas e apoio transversal} \hfill \metainfo{Tipo}{especialista} \hfill \metainfo{Identificador}{OpenCopywriter} \hfill \metainfo{Ficha}{copywriter.md}}

\begin{descricao}
Redige textos persuasivos, conteúdo de marketing e mensagens alinhadas à marca.
\end{descricao}

\begin{funcao}
No Core, este perfil apoia tarefas relacionadas a redige textos persuasivos, conteúdo de marketing e mensagens alinhadas à marca.. O orquestrador pode encaminhar-lhe o trabalho na etapa correspondente; a resposta retorna ao fluxo principal para integração e conferência de fontes, critérios e permissões. A ficha documenta a função, mas não comprova que o perfil tenha sido chamado ou executado a tarefa.
\end{funcao}

\begin{arquitetura}
\textbf{Modo de execução declarado:} subagente
\end{arquitetura}

\elemento{Redator técnico (v1.0.0)}{\metainfo{Categoria}{Ferramentas e apoio transversal} \hfill \metainfo{Tipo}{especialista} \hfill \metainfo{Identificador}{OpenTechnicalWriter} \hfill \metainfo{Ficha}{technical-writer.md}}

\begin{descricao}
Redige documentação, documentação de APIs e comunicação técnica.
\end{descricao}

\begin{funcao}
No Core, este perfil apoia tarefas relacionadas a redige documentação, documentação de APIs e comunicação técnica.. O orquestrador pode encaminhar-lhe o trabalho na etapa correspondente; a resposta retorna ao fluxo principal para integração e conferência de fontes, critérios e permissões. A ficha documenta a função, mas não comprova que o perfil tenha sido chamado ou executado a tarefa.
\end{funcao}

\begin{arquitetura}
\textbf{Modo de execução declarado:} subagente
\end{arquitetura}

\elemento{Optimizer (v1.0.0)}{\metainfo{Categoria}{Ferramentas e apoio transversal} \hfill \metainfo{Tipo}{especialista} \hfill \metainfo{Identificador}{optimizer} \hfill \metainfo{Ficha}{optimizer.md}}

\begin{descricao}
Otimiza performance de codigo (CPU, memoria, bundle, DB queries). Capacidades descritas: Codigo (cpu, memoria, bundle, db queries)
\end{descricao}

\begin{funcao}
No Core, este perfil apoia tarefas relacionadas a otimiza performance de codigo (CPU, memoria, bundle, DB queries).. O orquestrador pode encaminhar-lhe o trabalho na etapa correspondente; a resposta retorna ao fluxo principal para integração e conferência de fontes, critérios e permissões. A ficha documenta a função, mas não comprova que o perfil tenha sido chamado ou executado a tarefa.
\end{funcao}

\begin{arquitetura}
\textbf{Modo de execução declarado:} subagente · \textbf{Permissões:} edit: deny; bash: deny
\end{arquitetura}

\elemento{Pdf2latex agente (v1.0.0)}{\metainfo{Categoria}{Ferramentas e apoio transversal} \hfill \metainfo{Tipo}{especialista} \hfill \metainfo{Identificador}{pdf2latex-agent} \hfill \metainfo{Ficha}{pdf2latex-agent.md}}

\begin{descricao}
Parte do OpenCode Ecosystem Core SPEC-962 — Implementado em 2026-07-09 Agente especializado na conversão de documentos PDF para projetos LaTeX completos. Utiliza o pipeline pdf2latex com extração de texto, estrutura, figuras, tabelas, equações e referências, aplicando templates do catálogo do ecossistema. 1. Extrair texto: pdftotext + pdfminer + OCR (Tesseract) para PDFs escaneados 2. Detectar estrutura: capítulos, seções, subseções automaticamente 3. Extrair figuras: pdfimages + PyMuPDF (todas as imagens embutidas) 4. Detectar tabelas: camelot-py + pdfplumber → tabular LaTeX 5. Detectar equações: padrões LaTeX e heurísticas → math mode 6. Extrair referências: detecção de citações → arquivo .bib 7. Gerar projet
\end{descricao}

\begin{funcao}
No Core, este perfil apoia tarefas relacionadas a parte do OpenCode Ecosystem Core SPEC-962 — Implementado em 2026-07-09 Agente especializado na conversão de documentos PDF para projetos LaTeX completos.. O orquestrador pode encaminhar-lhe o trabalho na etapa correspondente; a resposta retorna ao fluxo principal para integração e conferência de fontes, critérios e permissões. A ficha documenta a função, mas não comprova que o perfil tenha sido chamado ou executado a tarefa.
\end{funcao}

\begin{arquitetura}
\textbf{Modo de execução declarado:} subagente · \textbf{Permissões:} edit: deny; bash: deny
\end{arquitetura}

\elemento{Plandex cli}{\metainfo{Categoria}{Ferramentas e apoio transversal} \hfill \metainfo{Tipo}{especialista} \hfill \metainfo{Identificador}{plandex-cli} \hfill \metainfo{Ficha}{plandex-cli.md}}

\begin{descricao}
Você é o agente proxy do Plandex, agente de codificação AI open source
\end{descricao}

\begin{funcao}
No Core, este perfil apoia tarefas relacionadas a você é o agente proxy do Plandex, agente de codificação AI open source. O orquestrador pode encaminhar-lhe o trabalho na etapa correspondente; a resposta retorna ao fluxo principal para integração e conferência de fontes, critérios e permissões. A ficha documenta a função, mas não comprova que o perfil tenha sido chamado ou executado a tarefa.
\end{funcao}

\begin{arquitetura}
\textbf{Modo de execução declarado:} subagente · \textbf{Permissões:} edit: deny; bash: deny
\end{arquitetura}

\elemento{Mecanismo de priorização (v1.0.0)}{\metainfo{Categoria}{Ferramentas e apoio transversal} \hfill \metainfo{Tipo}{especialista} \hfill \metainfo{Identificador}{prioritization-engine} \hfill \metainfo{Ficha}{prioritization-engine.md}}

\begin{descricao}
Pontua e prioriza itens de backlog com os métodos RICE e WSJF e organiza entregas entre MVP e etapas posteriores.
\end{descricao}

\begin{funcao}
No Core, este perfil apoia tarefas relacionadas a pontua e prioriza itens de backlog com os métodos RICE e WSJF e organiza entregas entre MVP e etapas posteriores.. O orquestrador pode encaminhar-lhe o trabalho na etapa correspondente; a resposta retorna ao fluxo principal para integração e conferência de fontes, critérios e permissões. A ficha documenta a função, mas não comprova que o perfil tenha sido chamado ou executado a tarefa.
\end{funcao}

\begin{arquitetura}
\textbf{Modo de execução declarado:} subagente · \textbf{Permissões:} edit: deny; bash: deny
\end{arquitetura}

\elemento{Quantum nexus phd (v1.0.0)}{\metainfo{Categoria}{Ferramentas e apoio transversal} \hfill \metainfo{Tipo}{especialista} \hfill \metainfo{Identificador}{quantum-nexus-phd} \hfill \metainfo{Ficha}{quantum-nexus-phd.md}}

\begin{descricao}
Super-Habilidade v7.2: Pesquisa quântica end-to-end com Qiskit+PennyLane, QML médico em HAM10000 (89.52\%), 50 qubits MPS, Grad-CAM, error mitigation (ZNE/PEC/DD), validação matemática, artigos Qualis A1, dashboards React. Orquestra agentes sob arquitetura Nexus Transformer v4.0.
\end{descricao}

\begin{funcao}
No Core, este perfil apoia tarefas relacionadas a super-Habilidade v7.2: Pesquisa quântica end-to-end com Qiskit+PennyLane, QML médico em HAM10000 (89.52\%), 50 qubits MPS, Grad-CAM, error mitigation (ZNE/PEC/DD), validação matemática, artigos Qualis A1…. O orquestrador pode encaminhar-lhe o trabalho na etapa correspondente; a resposta retorna ao fluxo principal para integração e conferência de fontes, critérios e permissões. A ficha documenta a função, mas não comprova que o perfil tenha sido chamado ou executado a tarefa.
\end{funcao}

\begin{arquitetura}
\textbf{Modo de execução declarado:} subagente · \textbf{Permissões:} edit: deny; bash: deny
\end{arquitetura}

\elemento{Reasonix cli}{\metainfo{Categoria}{Ferramentas e apoio transversal} \hfill \metainfo{Tipo}{integration} \hfill \metainfo{Identificador}{reasonix-cli} \hfill \metainfo{Ficha}{reasonix-cli.md}}

\begin{descricao}
Executor externo orquestrável: reasonix (esengine/DeepSeek-Reasonix, MIT). Agente de codificação DeepSeek-native; linha ativa é o rewrite Go (v2, instalável via npm install -g reasonix; alias dsnix). - reasonix\_available() / reasonix\_version() — presença e versão semântica. - reasonix\_run(task, directory, timeout) — one-shot reasonix run "" (streams para stdout, bom para pipes), sem shell (lista), nunca lança. - doctor\_check() — pass se binário + reasonix doctor exit 0, warn caso contrário (ex.: sem DeepSeek API key). - install\_instructions() — npm/npx + DeepSeek API key (+ reasonix setup). - CLI /reasonix status; run; doctor; install; setup com --timeout e --help. - Execução externa: resultados NÃO são "verifi
\end{descricao}

\begin{funcao}
No Core, este perfil apoia tarefas relacionadas a executor externo orquestrável: reasonix (esengine/DeepSeek-Reasonix, MIT).. O orquestrador pode encaminhar-lhe o trabalho na etapa correspondente; a resposta retorna ao fluxo principal para integração e conferência de fontes, critérios e permissões. A ficha documenta a função, mas não comprova que o perfil tenha sido chamado ou executado a tarefa.
\end{funcao}

\begin{arquitetura}
\textbf{Modo de execução declarado:} subagente · \textbf{Permissões:} edit: deny; bash: deny
\end{arquitetura}

\elemento{Reversa}{\metainfo{Categoria}{Ferramentas e apoio transversal} \hfill \metainfo{Tipo}{especialista} \hfill \metainfo{Identificador}{reversa} \hfill \metainfo{Ficha}{reversa.md}}

\begin{descricao}
Ponto de entrada principal do Reversa. Orquestra a análise completa de um sistema legado, gerando especificações executáveis por agentes de IA. Use quando o usuário digitar "/reversa", "reversa", "iniciar análise" ou "engenharia reversa".
\end{descricao}

\begin{funcao}
No Core, este perfil apoia tarefas relacionadas a ponto de entrada principal do Reversa.. O orquestrador pode encaminhar-lhe o trabalho na etapa correspondente; a resposta retorna ao fluxo principal para integração e conferência de fontes, critérios e permissões. A ficha documenta a função, mas não comprova que o perfil tenha sido chamado ou executado a tarefa.
\end{funcao}

\begin{arquitetura}
\textbf{Modo de execução declarado:} subagente · \textbf{Permissões:} read: allow; grep: allow; glob: allow; bash: allow; edit: allow; todoread: deny; todowrite: deny; webfetch: deny
\end{arquitetura}

\elemento{Reversa anp (v1.0.0)}{\metainfo{Categoria}{Ferramentas e apoio transversal} \hfill \metainfo{Tipo}{especialista} \hfill \metainfo{Identificador}{reversa-anp} \hfill \metainfo{Ficha}{reversa-anp.md}}

\begin{descricao}
agente especialista em constru description: >- Capacidade especializada em --- name: reversa-anp description: >- agente especialista em construir pipelin. tags: [name, reversa-anp, description, agente] examples: [Aplique name reversa anp description neste contexto, Avalie usando name reversa anp description]
\end{descricao}

\begin{funcao}
No Core, este perfil apoia tarefas relacionadas a agente especialista em constru description: >- Capacidade especializada em --- name: reversa-anp description: >- agente especialista em construir pipelin.. O orquestrador pode encaminhar-lhe o trabalho na etapa correspondente; a resposta retorna ao fluxo principal para integração e conferência de fontes, critérios e permissões. A ficha documenta a função, mas não comprova que o perfil tenha sido chamado ou executado a tarefa.
\end{funcao}

\begin{arquitetura}
\textbf{Modo de execução declarado:} subagente · \textbf{Permissões:} edit: deny; bash: deny
\end{arquitetura}

\elemento{Reversa archaeologist (v1.0.0)}{\metainfo{Categoria}{Ferramentas e apoio transversal} \hfill \metainfo{Tipo}{especialista} \hfill \metainfo{Identificador}{reversa-archaeologist} \hfill \metainfo{Ficha}{reversa-archaeologist.md}}

\begin{descricao}
Analisa profundamente o código do projeto legado módulo a módulo — extrai algoritmos, fluxos de controle, estruturas de dados e dicionário de dados. Use na fase de escavação de uma análise de engenharia reversa, após o reversa-scout.
\end{descricao}

\begin{funcao}
No Core, este perfil apoia tarefas relacionadas a analisa profundamente o código do projeto legado módulo a módulo — extrai algoritmos, fluxos de controle, estruturas de dados e dicionário de dados.. O orquestrador pode encaminhar-lhe o trabalho na etapa correspondente; a resposta retorna ao fluxo principal para integração e conferência de fontes, critérios e permissões. A ficha documenta a função, mas não comprova que o perfil tenha sido chamado ou executado a tarefa.
\end{funcao}

\begin{arquitetura}
\textbf{Modo de execução declarado:} subagente · \textbf{Permissões:} read: allow; grep: allow; glob: allow; bash: allow; edit: allow; todoread: deny; todowrite: deny; webfetch: deny
\end{arquitetura}

\elemento{Arquiteto de software Reversa (v1.0.0)}{\metainfo{Categoria}{Ferramentas e apoio transversal} \hfill \metainfo{Tipo}{especialista} \hfill \metainfo{Identificador}{reversa-architect} \hfill \metainfo{Ficha}{reversa-architect.md}}

\begin{descricao}
Sintetiza a análise do projeto legado em documentação arquitetural completa — diagramas C4, ERD completo, mapa de integrações e Spec Impact Matrix. Use na fase de interpretação após o reversa-detective.
\end{descricao}

\begin{funcao}
No Core, este perfil apoia tarefas relacionadas a sintetiza a análise do projeto legado em documentação arquitetural completa — diagramas C4, ERD completo, mapa de integrações e Spec Impact Matrix.. O orquestrador pode encaminhar-lhe o trabalho na etapa correspondente; a resposta retorna ao fluxo principal para integração e conferência de fontes, critérios e permissões. A ficha documenta a função, mas não comprova que o perfil tenha sido chamado ou executado a tarefa.
\end{funcao}

\begin{arquitetura}
\textbf{Modo de execução declarado:} subagente · \textbf{Permissões:} read: allow; grep: allow; glob: allow; bash: allow; edit: allow; todoread: deny; todowrite: deny; webfetch: deny
\end{arquitetura}

\elemento{Reversa config generator (v1.0.0)}{\metainfo{Categoria}{Ferramentas e apoio transversal} \hfill \metainfo{Tipo}{especialista} \hfill \metainfo{Identificador}{reversa-config-generator} \hfill \metainfo{Ficha}{reversa-config-generator.md}}

\begin{descricao}
Subagente especializado em geracao de configuracoes complexas usando LLM em multiplas etapas com fallback heuristico. Conhece o padrao P13 (Step-by-step LLM Config Generator) e suas 4 etapas: tempo, eventos, agentes em lote e plataforma.
\end{descricao}

\begin{funcao}
No Core, este perfil apoia tarefas relacionadas a subagente especializado em geracao de configuracoes complexas usando LLM em multiplas etapas com fallback heuristico.. O orquestrador pode encaminhar-lhe o trabalho na etapa correspondente; a resposta retorna ao fluxo principal para integração e conferência de fontes, critérios e permissões. A ficha documenta a função, mas não comprova que o perfil tenha sido chamado ou executado a tarefa.
\end{funcao}

\begin{arquitetura}
\textbf{Modo de execução declarado:} subagente · \textbf{Permissões:} edit: deny; bash: deny
\end{arquitetura}

\elemento{Reversa dados master (v1.0.0)}{\metainfo{Categoria}{Ferramentas e apoio transversal} \hfill \metainfo{Tipo}{especialista} \hfill \metainfo{Identificador}{reversa-data-master} \hfill \metainfo{Ficha}{reversa-data-master.md}}

\begin{descricao}
Documenta completamente o banco de dados do projeto legado — tabelas, relacionamentos, constraints, triggers, procedures e ERD completo. Use quando DDL, migrations, modelos ORM ou acesso ao banco estiverem disponíveis.
\end{descricao}

\begin{funcao}
No Core, este perfil apoia tarefas relacionadas a documenta completamente o banco de dados do projeto legado — tabelas, relacionamentos, constraints, triggers, procedures e ERD completo.. O orquestrador pode encaminhar-lhe o trabalho na etapa correspondente; a resposta retorna ao fluxo principal para integração e conferência de fontes, critérios e permissões. A ficha documenta a função, mas não comprova que o perfil tenha sido chamado ou executado a tarefa.
\end{funcao}

\begin{arquitetura}
\textbf{Modo de execução declarado:} subagente · \textbf{Permissões:} read: allow; grep: allow; glob: allow; bash: allow; edit: allow; todoread: deny; todowrite: deny; webfetch: deny
\end{arquitetura}

\elemento{Reversa design system (v1.0.0)}{\metainfo{Categoria}{Ferramentas e apoio transversal} \hfill \metainfo{Tipo}{especialista} \hfill \metainfo{Identificador}{reversa-design-system} \hfill \metainfo{Ficha}{reversa-design-system.md}}

\begin{descricao}
Extrai e documenta o sistema de design do projeto legado — paleta de cores, tipografia, espaçamentos, tokens e componentes a partir de CSS, arquivos de tema e screenshots.
\end{descricao}

\begin{funcao}
No Core, este perfil apoia tarefas relacionadas a extrai e documenta o sistema de design do projeto legado — paleta de cores, tipografia, espaçamentos, tokens e componentes a partir de CSS, arquivos…. O orquestrador pode encaminhar-lhe o trabalho na etapa correspondente; a resposta retorna ao fluxo principal para integração e conferência de fontes, critérios e permissões. A ficha documenta a função, mas não comprova que o perfil tenha sido chamado ou executado a tarefa.
\end{funcao}

\begin{arquitetura}
\textbf{Modo de execução declarado:} subagente · \textbf{Permissões:} read: allow; grep: allow; glob: allow; bash: allow; edit: allow; todoread: deny; todowrite: deny; webfetch: deny
\end{arquitetura}

\elemento{Reversa detective (v1.0.0)}{\metainfo{Categoria}{Ferramentas e apoio transversal} \hfill \metainfo{Tipo}{especialista} \hfill \metainfo{Identificador}{reversa-detective} \hfill \metainfo{Ficha}{reversa-detective.md}}

\begin{descricao}
Extrai conhecimento de negócio implícito do projeto legado — regras de negócio, ADRs retroativos via Git, máquinas de estado e matriz de permissões. Use na fase de interpretação de uma análise de engenharia reversa.
\end{descricao}

\begin{funcao}
No Core, este perfil apoia tarefas relacionadas a extrai conhecimento de negócio implícito do projeto legado — regras de negócio, ADRs retroativos via Git, máquinas de estado e matriz de permissões.. O orquestrador pode encaminhar-lhe o trabalho na etapa correspondente; a resposta retorna ao fluxo principal para integração e conferência de fontes, critérios e permissões. A ficha documenta a função, mas não comprova que o perfil tenha sido chamado ou executado a tarefa.
\end{funcao}

\begin{arquitetura}
\textbf{Modo de execução declarado:} subagente · \textbf{Permissões:} read: allow; grep: allow; glob: allow; bash: allow; edit: allow; todoread: deny; todowrite: deny; webfetch: deny
\end{arquitetura}

\elemento{Reversa entity ner (v1.0.0)}{\metainfo{Categoria}{Ferramentas e apoio transversal} \hfill \metainfo{Tipo}{especialista} \hfill \metainfo{Identificador}{reversa-entity-ner} \hfill \metainfo{Ficha}{reversa-entity-ner.md}}

\begin{descricao}
Agente de leitura e filtragem de entidades em grafos de conhecimento. Inspirado pelo EntityReader do MiroFish-Offline. Lista entidades, filtra por tipo, obtém contexto completo com arestas e nós vizinhos. Use via: "entidade", "entity", "NER", /entity-ner.
\end{descricao}

\begin{funcao}
No Core, este perfil apoia tarefas relacionadas a agente de leitura e filtragem de entidades em grafos de conhecimento.. O orquestrador pode encaminhar-lhe o trabalho na etapa correspondente; a resposta retorna ao fluxo principal para integração e conferência de fontes, critérios e permissões. A ficha documenta a função, mas não comprova que o perfil tenha sido chamado ou executado a tarefa.
\end{funcao}

\begin{arquitetura}
\textbf{Modo de execução declarado:} subagente · \textbf{Permissões:} edit: deny; bash: deny
\end{arquitetura}

\elemento{Reversa fileipc (v1.0.0)}{\metainfo{Categoria}{Ferramentas e apoio transversal} \hfill \metainfo{Tipo}{especialista} \hfill \metainfo{Identificador}{reversa-fileipc} \hfill \metainfo{Ficha}{reversa-fileipc.md}}

\begin{descricao}
Agente de comunicação via filesystem do ecossistema Reversa. Orquestra a troca de mensagens entre processos usando o protocolo File IPC, permitindo comunicação assíncrona sem dependências externas.
\end{descricao}

\begin{funcao}
No Core, este perfil apoia tarefas relacionadas a agente de comunicação via filesystem do ecossistema Reversa.. O orquestrador pode encaminhar-lhe o trabalho na etapa correspondente; a resposta retorna ao fluxo principal para integração e conferência de fontes, critérios e permissões. A ficha documenta a função, mas não comprova que o perfil tenha sido chamado ou executado a tarefa.
\end{funcao}

\begin{arquitetura}
\textbf{Modo de execução declarado:} subagente · \textbf{Permissões:} read: allow; edit: allow; glob: allow; bash: allow; sequential\_thinking: allow
\end{arquitetura}

\elemento{Reversa grafo construtor (v1.0.0)}{\metainfo{Categoria}{Ferramentas e apoio transversal} \hfill \metainfo{Tipo}{especialista} \hfill \metainfo{Identificador}{reversa-graph-builder} \hfill \metainfo{Ficha}{reversa-graph-builder.md}}

\begin{descricao}
Agente de construção assíncrona de grafos de conhecimento. Inspirado pelo GraphBuilderService do MiroFish-Offline. Processa texto em chunks, aplica ontologia e persiste em SQLite/Neo4j. Use via: "construir grafo", "build graph", /graph-builder.
\end{descricao}

\begin{funcao}
No Core, este perfil apoia tarefas relacionadas a agente de construção assíncrona de grafos de conhecimento.. O orquestrador pode encaminhar-lhe o trabalho na etapa correspondente; a resposta retorna ao fluxo principal para integração e conferência de fontes, critérios e permissões. A ficha documenta a função, mas não comprova que o perfil tenha sido chamado ou executado a tarefa.
\end{funcao}

\begin{arquitetura}
\textbf{Modo de execução declarado:} subagente · \textbf{Permissões:} edit: deny; bash: deny
\end{arquitetura}

\elemento{Reversa graphrag (v1.0.0)}{\metainfo{Categoria}{Ferramentas e apoio transversal} \hfill \metainfo{Tipo}{especialista} \hfill \metainfo{Identificador}{reversa-graphrag} \hfill \metainfo{Ficha}{reversa-graphrag.md}}

\begin{descricao}
Agente de conhecimento que constrói e consulta o grafo de dependências do ecossistema OpenCode. Inspirado pelo GraphRAG + Zep Cloud do MiroFish (graph\_builder.py, zep\_tools.py). Usa SQLite para persistência com busca estrutural e semântica. Use via: "grafo", "graph", "dependências", "knowledge graph", /graphrag.
\end{descricao}

\begin{funcao}
No Core, este perfil apoia tarefas relacionadas a agente de conhecimento que constrói e consulta o grafo de dependências do ecossistema OpenCode.. O orquestrador pode encaminhar-lhe o trabalho na etapa correspondente; a resposta retorna ao fluxo principal para integração e conferência de fontes, critérios e permissões. A ficha documenta a função, mas não comprova que o perfil tenha sido chamado ou executado a tarefa.
\end{funcao}

\begin{arquitetura}
\textbf{Modo de execução declarado:} subagente · \textbf{Permissões:} edit: deny; bash: deny
\end{arquitetura}

\elemento{Reversa hybrid grafo (v1.0.0)}{\metainfo{Categoria}{Ferramentas e apoio transversal} \hfill \metainfo{Tipo}{especialista} \hfill \metainfo{Identificador}{reversa-hybrid-graph} \hfill \metainfo{Ficha}{reversa-hybrid-graph.md}}

\begin{descricao}
Agente de busca híbrida em grafo de conhecimento. Oferece 3 estratégias complementares: InsightForge (análise profunda), PanoramaSearch (visão panorâmica), QuickSearch (busca rápida). Inspirado pelo GraphToolsService do MiroFish-Offline. Use via: "busca", "graph", "grafo", "pesquisa", "insight", "panorama", "quick", /hybrid-graph.
\end{descricao}

\begin{funcao}
No Core, este perfil apoia tarefas relacionadas a agente de busca híbrida em grafo de conhecimento.. O orquestrador pode encaminhar-lhe o trabalho na etapa correspondente; a resposta retorna ao fluxo principal para integração e conferência de fontes, critérios e permissões. A ficha documenta a função, mas não comprova que o perfil tenha sido chamado ou executado a tarefa.
\end{funcao}

\begin{arquitetura}
\textbf{Modo de execução declarado:} subagente · \textbf{Permissões:} edit: deny; bash: deny
\end{arquitetura}

\elemento{Reversa memória updater (v1.0.0)}{\metainfo{Categoria}{Ferramentas e apoio transversal} \hfill \metainfo{Tipo}{especialista} \hfill \metainfo{Identificador}{reversa-memory-updater} \hfill \metainfo{Ficha}{reversa-memory-updater.md}}

\begin{descricao}
Agente de atualização de grafos em tempo real com atividades de simulação. Inspirado pelo GraphMemoryUpdater do MiroFish-Offline. Monitora ações de agentes, bufferiza por plataforma e envia em lotes. Use via: "memória", "memory", "atualizar grafo", /memory-updater.
\end{descricao}

\begin{funcao}
No Core, este perfil apoia tarefas relacionadas a agente de atualização de grafos em tempo real com atividades de simulação.. O orquestrador pode encaminhar-lhe o trabalho na etapa correspondente; a resposta retorna ao fluxo principal para integração e conferência de fontes, critérios e permissões. A ficha documenta a função, mas não comprova que o perfil tenha sido chamado ou executado a tarefa.
\end{funcao}

\begin{arquitetura}
\textbf{Modo de execução declarado:} subagente · \textbf{Permissões:} edit: deny; bash: deny
\end{arquitetura}

\elemento{Reversa oasis profile (v1.0.0)}{\metainfo{Categoria}{Ferramentas e apoio transversal} \hfill \metainfo{Tipo}{especialista} \hfill \metainfo{Identificador}{reversa-oasis-profile} \hfill \metainfo{Ficha}{reversa-oasis-profile.md}}

\begin{descricao}
Agente gerador de perfis OASIS. Converte entidades de grafos de conhecimento em personas detalhadas de agente para simulação. Inspirado pelo OASIS Profile Generator do MiroFish-Offline. Gera bio, persona, interesses, MBTI, tópicos, estilo de fala e comportamentos por plataforma (Twitter/Reddit). Use via: "perfil", "persona", "profile", "oasis", /oasis-profile.
\end{descricao}

\begin{funcao}
No Core, este perfil apoia tarefas relacionadas a agente gerador de perfis OASIS.. O orquestrador pode encaminhar-lhe o trabalho na etapa correspondente; a resposta retorna ao fluxo principal para integração e conferência de fontes, critérios e permissões. A ficha documenta a função, mas não comprova que o perfil tenha sido chamado ou executado a tarefa.
\end{funcao}

\begin{arquitetura}
\textbf{Modo de execução declarado:} subagente · \textbf{Permissões:} edit: deny; bash: deny
\end{arquitetura}

\elemento{Reversa ontology gen (v1.0.0)}{\metainfo{Categoria}{Ferramentas e apoio transversal} \hfill \metainfo{Tipo}{especialista} \hfill \metainfo{Identificador}{reversa-ontology-gen} \hfill \metainfo{Ficha}{reversa-ontology-gen.md}}

\begin{descricao}
Agente de geração de ontologias para grafos de conhecimento. Inspirado pelo OntologyGenerator do MiroFish-Offline. Gera tipos de entidades e relacionamentos a partir de textos e requisitos. Use via: "ontologia", "ontology", "tipos", /ontology-gen.
\end{descricao}

\begin{funcao}
No Core, este perfil apoia tarefas relacionadas a agente de geração de ontologias para grafos de conhecimento.. O orquestrador pode encaminhar-lhe o trabalho na etapa correspondente; a resposta retorna ao fluxo principal para integração e conferência de fontes, critérios e permissões. A ficha documenta a função, mas não comprova que o perfil tenha sido chamado ou executado a tarefa.
\end{funcao}

\begin{arquitetura}
\textbf{Modo de execução declarado:} subagente · \textbf{Permissões:} edit: deny; bash: deny
\end{arquitetura}

\elemento{Reversa planner}{\metainfo{Categoria}{Ferramentas e apoio transversal} \hfill \metainfo{Tipo}{especialista} \hfill \metainfo{Identificador}{reversa-planner} \hfill \metainfo{Ficha}{reversa-planner.md}}

\begin{descricao}
Gera planos de engenharia reversa em etapas (Scope → Modules → Tasks → Dependencies → Resources), inspirado pelo simulation\_config\_generator.py do MiroFish que gera parâmetros de simulação em 4 etapas para evitar estouro de contexto. Use via: "gerar plano", "generate plan", "planejar análise", ou /generate-plan.
\end{descricao}

\begin{funcao}
No Core, este perfil apoia tarefas relacionadas a gera planos de engenharia reversa em etapas (Scope → Modules → Tasks → Dependencies → Resources), inspirado pelo simulation\_config\_generator.py do MiroFish que gera parâmetros…. O orquestrador pode encaminhar-lhe o trabalho na etapa correspondente; a resposta retorna ao fluxo principal para integração e conferência de fontes, critérios e permissões. A ficha documenta a função, mas não comprova que o perfil tenha sido chamado ou executado a tarefa.
\end{funcao}

\begin{arquitetura}
\textbf{Modo de execução declarado:} subagente · \textbf{Permissões:} read: allow; grep: allow; glob: allow; bash: allow; edit: allow; todoread: deny; todowrite: deny; webfetch: deny
\end{arquitetura}

\elemento{Agente de elaboração de relatórios Reversa (v1.0.0)}{\metainfo{Categoria}{Ferramentas e apoio transversal} \hfill \metainfo{Tipo}{especialista} \hfill \metainfo{Identificador}{reversa-report-agent} \hfill \metainfo{Ficha}{reversa-report-agent.md}}

\begin{descricao}
Produz relatórios a partir de informações reunidas por outras etapas do fluxo Reversa, preservando a organização e a rastreabilidade do conteúdo.
\end{descricao}

\begin{funcao}
No Core, este perfil apoia tarefas relacionadas a produz relatórios a partir de informações reunidas por outras etapas do fluxo Reversa, preservando a organização e a rastreabilidade do conteúdo.. O orquestrador pode encaminhar-lhe o trabalho na etapa correspondente; a resposta retorna ao fluxo principal para integração e conferência de fontes, critérios e permissões. A ficha documenta a função, mas não comprova que o perfil tenha sido chamado ou executado a tarefa.
\end{funcao}

\begin{arquitetura}
\textbf{Modo de execução declarado:} subagente · \textbf{Permissões:} edit: deny; bash: deny
\end{arquitetura}

\elemento{Revisor de reversa (v1.0.0)}{\metainfo{Categoria}{Ferramentas e apoio transversal} \hfill \metainfo{Tipo}{especialista} \hfill \metainfo{Identificador}{reversa-reviewer} \hfill \metainfo{Ficha}{reversa-reviewer.md}}

\begin{descricao}
Revisa criticamente as especificações geradas pelo reversa-writer — encontra inconsistências, reclassifica confiança e gera perguntas para validação humana. Use na fase de revisão de uma análise de engenharia reversa.
\end{descricao}

\begin{funcao}
No Core, este perfil apoia tarefas relacionadas a revisa criticamente as especificações geradas pelo reversa-writer — encontra inconsistências, reclassifica confiança e gera perguntas para validação humana.. O orquestrador pode encaminhar-lhe o trabalho na etapa correspondente; a resposta retorna ao fluxo principal para integração e conferência de fontes, critérios e permissões. A ficha documenta a função, mas não comprova que o perfil tenha sido chamado ou executado a tarefa.
\end{funcao}

\begin{arquitetura}
\textbf{Modo de execução declarado:} subagente · \textbf{Permissões:} read: allow; grep: allow; glob: allow; bash: allow; edit: allow; todoread: deny; todowrite: deny; webfetch: deny
\end{arquitetura}

\elemento{Explorador de reversa (v1.0.0)}{\metainfo{Categoria}{Ferramentas e apoio transversal} \hfill \metainfo{Tipo}{especialista} \hfill \metainfo{Identificador}{reversa-scout} \hfill \metainfo{Ficha}{reversa-scout.md}}

\begin{descricao}
Mapeia a superfície do projeto legado — estrutura de pastas, linguagens, frameworks, dependências e entry points. Use no início de uma análise de engenharia reversa para criar o inventário inicial do projeto.
\end{descricao}

\begin{funcao}
No Core, este perfil apoia tarefas relacionadas a mapeia a superfície do projeto legado — estrutura de pastas, linguagens, frameworks, dependências e entry points.. O orquestrador pode encaminhar-lhe o trabalho na etapa correspondente; a resposta retorna ao fluxo principal para integração e conferência de fontes, critérios e permissões. A ficha documenta a função, mas não comprova que o perfil tenha sido chamado ou executado a tarefa.
\end{funcao}

\begin{arquitetura}
\textbf{Modo de execução declarado:} subagente · \textbf{Permissões:} read: allow; grep: allow; glob: allow; bash: allow; edit: allow; todoread: deny; todowrite: deny; webfetch: deny
\end{arquitetura}

\elemento{Reversa statemachine (v1.0.0)}{\metainfo{Categoria}{Ferramentas e apoio transversal} \hfill \metainfo{Tipo}{especialista} \hfill \metainfo{Identificador}{reversa-statemachine} \hfill \metainfo{Ficha}{reversa-statemachine.md}}

\begin{descricao}
Agente de máquina de estados do pipeline Reversa. Gerencia transições de estado, valida dependências entre fases e mantém persistência das etapas do pipeline de engenharia reversa.
\end{descricao}

\begin{funcao}
No Core, este perfil apoia tarefas relacionadas a agente de máquina de estados do pipeline Reversa.. O orquestrador pode encaminhar-lhe o trabalho na etapa correspondente; a resposta retorna ao fluxo principal para integração e conferência de fontes, critérios e permissões. A ficha documenta a função, mas não comprova que o perfil tenha sido chamado ou executado a tarefa.
\end{funcao}

\begin{arquitetura}
\textbf{Modo de execução declarado:} subagente · \textbf{Permissões:} read: allow; edit: allow; bash: allow; sqlite: allow; sequential\_thinking: allow
\end{arquitetura}

\elemento{Reversa swarm review (v1.0.0)}{\metainfo{Categoria}{Ferramentas e apoio transversal} \hfill \metainfo{Tipo}{especialista} \hfill \metainfo{Identificador}{reversa-swarm-review} \hfill \metainfo{Ficha}{reversa-swarm-review.md}}

\begin{descricao}
Orquestra revisão de código por enxame de agentes especializados (Segurança, Performance, Arquitetura), inspirado no padrão de simulação multi-agente do MiroFish. Cada agente analisa independentemente, debate divergências e produz relatório consolidado com múltiplas perspectivas. Use via: "swarm review", "revisão em enxame", ou ativado pelo comando /swarm-review.
\end{descricao}

\begin{funcao}
No Core, este perfil apoia tarefas relacionadas a orquestra revisão de código por enxame de agentes especializados (Segurança, Performance, Arquitetura), inspirado no padrão de simulação multi-agente do MiroFish.. O orquestrador pode encaminhar-lhe o trabalho na etapa correspondente; a resposta retorna ao fluxo principal para integração e conferência de fontes, critérios e permissões. A ficha documenta a função, mas não comprova que o perfil tenha sido chamado ou executado a tarefa.
\end{funcao}

\begin{arquitetura}
\textbf{Modo de execução declarado:} subagente · \textbf{Permissões:} edit: deny; bash: deny
\end{arquitetura}

\elemento{Reversa synthesis}{\metainfo{Categoria}{Ferramentas e apoio transversal} \hfill \metainfo{Tipo}{especialista} \hfill \metainfo{Identificador}{reversa-synthesis} \hfill \metainfo{Ficha}{reversa-synthesis.md}}

\begin{descricao}
Meta-agente sintetizador que coleta outputs de múltiplos agentes Reversa, cruza referências, identifica lacunas e produz documentação consolidada. Inspirado pelo ReportAgent com padrão ReACT do MiroFish. Use via: "sintetizar", "consolidar", "unificar achados", ou /synthesize.
\end{descricao}

\begin{funcao}
No Core, este perfil apoia tarefas relacionadas a meta-agente sintetizador que coleta outputs de múltiplos agentes Reversa, cruza referências, identifica lacunas e produz documentação consolidada.. O orquestrador pode encaminhar-lhe o trabalho na etapa correspondente; a resposta retorna ao fluxo principal para integração e conferência de fontes, critérios e permissões. A ficha documenta a função, mas não comprova que o perfil tenha sido chamado ou executado a tarefa.
\end{funcao}

\begin{arquitetura}
\textbf{Modo de execução declarado:} subagente · \textbf{Permissões:} read: allow; grep: allow; glob: allow; bash: allow; edit: allow; todoread: deny; todowrite: deny; webfetch: deny
\end{arquitetura}

\elemento{Reversa visor (v1.0.0)}{\metainfo{Categoria}{Ferramentas e apoio transversal} \hfill \metainfo{Tipo}{especialista} \hfill \metainfo{Identificador}{reversa-visor} \hfill \metainfo{Ficha}{reversa-visor.md}}

\begin{descricao}
Documenta a interface do sistema legado a partir de screenshots — extrai componentes, layouts, fluxos de navegação e estados de tela. Use quando screenshots do sistema estiverem disponíveis.
\end{descricao}

\begin{funcao}
No Core, este perfil apoia tarefas relacionadas a documenta a interface do sistema legado a partir de screenshots — extrai componentes, layouts, fluxos de navegação e estados de tela.. O orquestrador pode encaminhar-lhe o trabalho na etapa correspondente; a resposta retorna ao fluxo principal para integração e conferência de fontes, critérios e permissões. A ficha documenta a função, mas não comprova que o perfil tenha sido chamado ou executado a tarefa.
\end{funcao}

\begin{arquitetura}
\textbf{Modo de execução declarado:} subagente · \textbf{Permissões:} read: allow; grep: allow; glob: allow; bash: allow; edit: allow; todoread: deny; todowrite: deny; webfetch: deny
\end{arquitetura}

\elemento{Redator de reversa (v1.0.0)}{\metainfo{Categoria}{Ferramentas e apoio transversal} \hfill \metainfo{Tipo}{especialista} \hfill \metainfo{Identificador}{reversa-writer} \hfill \metainfo{Ficha}{reversa-writer.md}}

\begin{descricao}
Gera especificações executáveis do sistema legado como contratos operacionais, em formato de pasta-por-unit com requirements.md, design.md e tasks.md. Use na fase de geração de uma análise de engenharia reversa.
\end{descricao}

\begin{funcao}
No Core, este perfil apoia tarefas relacionadas a gera especificações executáveis do sistema legado como contratos operacionais, em formato de pasta-por-unit com requirements.md, design.md e tasks.md.. O orquestrador pode encaminhar-lhe o trabalho na etapa correspondente; a resposta retorna ao fluxo principal para integração e conferência de fontes, critérios e permissões. A ficha documenta a função, mas não comprova que o perfil tenha sido chamado ou executado a tarefa.
\end{funcao}

\begin{arquitetura}
\textbf{Modo de execução declarado:} subagente · \textbf{Permissões:} read: allow; grep: allow; glob: allow; bash: allow; edit: allow; todoread: deny; todowrite: deny; webfetch: deny
\end{arquitetura}

\elemento{Revisor de código e segurança (v1.0.0)}{\metainfo{Categoria}{Ferramentas e apoio transversal} \hfill \metainfo{Tipo}{especialista} \hfill \metainfo{Identificador}{reviewer} \hfill \metainfo{Ficha}{reviewer.md}}

\begin{descricao}
Revisa código, segurança e qualidade segundo o escopo configurado.
\end{descricao}

\begin{funcao}
No Core, este perfil apoia tarefas relacionadas a revisa código, segurança e qualidade segundo o escopo configurado.. O orquestrador pode encaminhar-lhe o trabalho na etapa correspondente; a resposta retorna ao fluxo principal para integração e conferência de fontes, critérios e permissões. A ficha documenta a função, mas não comprova que o perfil tenha sido chamado ou executado a tarefa.
\end{funcao}

\begin{arquitetura}
\textbf{Modo de execução declarado:} subagente · \textbf{Permissões:} edit: deny; bash: deny
\end{arquitetura}

\elemento{Segurança auditor (v1.0.0)}{\metainfo{Categoria}{Ferramentas e apoio transversal} \hfill \metainfo{Tipo}{especialista} \hfill \metainfo{Identificador}{security-auditor} \hfill \metainfo{Ficha}{security-auditor.md}}

\begin{descricao}
Voce e especialista em seguranca de aplicacoes. - Auth: hash senhas (bcrypt/argon2), JWT refresh, RBAC, rate limit - Input: validar entradas, SQL injection, XSS, CSRF - Dados: sem secrets em logs, sem hardcode, criptografia, HTTPS - Dependencias: atualizadas, minimas, lockfile - Config: debug off prod, CORS correto, sem stack traces Formato: Severidade -> CWE -> Arquivo -> Risco -> Correcao
\end{descricao}

\begin{funcao}
No Core, este perfil apoia tarefas relacionadas a voce e especialista em seguranca de aplicacoes.. O orquestrador pode encaminhar-lhe o trabalho na etapa correspondente; a resposta retorna ao fluxo principal para integração e conferência de fontes, critérios e permissões. A ficha documenta a função, mas não comprova que o perfil tenha sido chamado ou executado a tarefa.
\end{funcao}

\begin{arquitetura}
\textbf{Modo de execução declarado:} subagente · \textbf{Permissões:} edit: deny; bash: deny
\end{arquitetura}

\elemento{Agente de resposta simples para avaliação (v1.0.0)}{\metainfo{Categoria}{Ferramentas e apoio transversal} \hfill \metainfo{Tipo}{especialista} \hfill \metainfo{Identificador}{simple-responder} \hfill \metainfo{Ficha}{simple-responder.md}}

\begin{descricao}
Perfil de teste que retorna a expressão AWESOME TESTING; destina-se à avaliação do framework, não a tarefas produtivas.
\end{descricao}

\begin{funcao}
No Core, este perfil apoia tarefas relacionadas a perfil de teste que retorna a expressão AWESOME TESTING. O orquestrador pode encaminhar-lhe o trabalho na etapa correspondente; a resposta retorna ao fluxo principal para integração e conferência de fontes, critérios e permissões. A ficha documenta a função, mas não comprova que o perfil tenha sido chamado ou executado a tarefa.
\end{funcao}

\begin{arquitetura}
\textbf{Modo de execução declarado:} subagente · \textbf{Permissões:} edit: deny; bash: deny
\end{arquitetura}

\elemento{Orquestrador de etapas}{\metainfo{Categoria}{Ferramentas e apoio transversal} \hfill \metainfo{Tipo}{especialista} \hfill \metainfo{Identificador}{StageOrchestrator} \hfill \metainfo{Ficha}{stage-orchestrator.md}}

\begin{descricao}
Orquestra etapas de fluxos complexos, aplicando transições, regras de passagem, validação e reversão.
\end{descricao}

\begin{funcao}
No Core, este perfil apoia tarefas relacionadas a orquestra etapas de fluxos complexos, aplicando transições, regras de passagem, validação e reversão.. O orquestrador pode encaminhar-lhe o trabalho na etapa correspondente; a resposta retorna ao fluxo principal para integração e conferência de fontes, critérios e permissões. A ficha documenta a função, mas não comprova que o perfil tenha sido chamado ou executado a tarefa.
\end{funcao}

\begin{arquitetura}
\textbf{Modo de execução declarado:} subagente
\end{arquitetura}

\elemento{Mapeador de jornadas e histórias (v1.0.0)}{\metainfo{Categoria}{Ferramentas e apoio transversal} \hfill \metainfo{Tipo}{especialista} \hfill \metainfo{Identificador}{story-mapper} \hfill \metainfo{Ficha}{story-mapper.md}}

\begin{descricao}
Converte necessidades de usuários em jornadas, épicos, histórias e entregas verticais relacionadas a contextos delimitados.
\end{descricao}

\begin{funcao}
No Core, este perfil apoia tarefas relacionadas a converte necessidades de usuários em jornadas, épicos, histórias e entregas verticais relacionadas a contextos delimitados.. O orquestrador pode encaminhar-lhe o trabalho na etapa correspondente; a resposta retorna ao fluxo principal para integração e conferência de fontes, critérios e permissões. A ficha documenta a função, mas não comprova que o perfil tenha sido chamado ou executado a tarefa.
\end{funcao}

\begin{arquitetura}
\textbf{Modo de execução declarado:} subagente · \textbf{Permissões:} edit: deny; bash: deny
\end{arquitetura}

\elemento{Gestor de tarefas de desenvolvimento (v1.0.0)}{\metainfo{Categoria}{Ferramentas e apoio transversal} \hfill \metainfo{Tipo}{especialista} \hfill \metainfo{Identificador}{task-manager} \hfill \metainfo{Ficha}{task-manager.md}}

\begin{descricao}
Decompõe funcionalidades em subtarefas verificáveis, acompanha dependências em JSON e integra o fluxo à CLI.
\end{descricao}

\begin{funcao}
No Core, este perfil apoia tarefas relacionadas a decompõe funcionalidades em subtarefas verificáveis, acompanha dependências em JSON e integra o fluxo à CLI.. O orquestrador pode encaminhar-lhe o trabalho na etapa correspondente; a resposta retorna ao fluxo principal para integração e conferência de fontes, critérios e permissões. A ficha documenta a função, mas não comprova que o perfil tenha sido chamado ou executado a tarefa.
\end{funcao}

\begin{arquitetura}
\textbf{Modo de execução declarado:} subagente · \textbf{Permissões:} edit: deny; bash: deny
\end{arquitetura}

\elemento{Engenheiro de testes (v1.0.0)}{\metainfo{Categoria}{Ferramentas e apoio transversal} \hfill \metainfo{Tipo}{especialista} \hfill \metainfo{Identificador}{test-engineer} \hfill \metainfo{Ficha}{test-engineer.md}}

\begin{descricao}
Planeja e escreve testes conforme as convenções do projeto, considerando a testabilidade antes da implementação e buscando resultados determinísticos.
\end{descricao}

\begin{funcao}
No Core, este perfil apoia tarefas relacionadas a planeja e escreve testes conforme as convenções do projeto, considerando a testabilidade antes da implementação e buscando resultados determinísticos.. O orquestrador pode encaminhar-lhe o trabalho na etapa correspondente; a resposta retorna ao fluxo principal para integração e conferência de fontes, critérios e permissões. A ficha documenta a função, mas não comprova que o perfil tenha sido chamado ou executado a tarefa.
\end{funcao}

\begin{arquitetura}
\textbf{Modo de execução declarado:} subagente · \textbf{Permissões:} edit: deny; bash: deny
\end{arquitetura}

\elemento{Analista de documentos de pesquisa (v1.0.0)}{\metainfo{Categoria}{Ferramentas e apoio transversal} \hfill \metainfo{Tipo}{especialista} \hfill \metainfo{Identificador}{thoughts-analyzer} \hfill \metainfo{Ficha}{thoughts-analyzer.md}}

\begin{descricao}
Aprofunda temas de pesquisa por meio da análise de documentos e materiais disponíveis no repositório.
\end{descricao}

\begin{funcao}
No Core, este perfil apoia tarefas relacionadas a aprofunda temas de pesquisa por meio da análise de documentos e materiais disponíveis no repositório.. O orquestrador pode encaminhar-lhe o trabalho na etapa correspondente; a resposta retorna ao fluxo principal para integração e conferência de fontes, critérios e permissões. A ficha documenta a função, mas não comprova que o perfil tenha sido chamado ou executado a tarefa.
\end{funcao}

\begin{arquitetura}
\textbf{Modo de execução declarado:} subagente · \textbf{Permissões:} read: allow; grep: allow; glob: allow; list: allow; bash: deny; edit: deny; todoread: deny; todowrite: deny; webfetch: deny
\end{arquitetura}

\elemento{Localizador de documentos de pesquisa (v1.0.0)}{\metainfo{Categoria}{Ferramentas e apoio transversal} \hfill \metainfo{Tipo}{especialista} \hfill \metainfo{Identificador}{thoughts-locator} \hfill \metainfo{Ficha}{thoughts-locator.md}}

\begin{descricao}
Localiza documentos relevantes no diretório thoughts, usado pelo projeto para armazenar metadados e notas.
\end{descricao}

\begin{funcao}
No Core, este perfil apoia tarefas relacionadas a localiza documentos relevantes no diretório thoughts, usado pelo projeto para armazenar metadados e notas.. O orquestrador pode encaminhar-lhe o trabalho na etapa correspondente; a resposta retorna ao fluxo principal para integração e conferência de fontes, critérios e permissões. A ficha documenta a função, mas não comprova que o perfil tenha sido chamado ou executado a tarefa.
\end{funcao}

\begin{arquitetura}
\textbf{Modo de execução declarado:} subagente · \textbf{Permissões:} read: allow; grep: allow; glob: allow; list: allow; bash: deny; edit: deny; todoread: deny; todowrite: deny; webfetch: deny
\end{arquitetura}

\elemento{Desenvolvedor de aplicações web (v1.0.0)}{\metainfo{Categoria}{Ferramentas e apoio transversal} \hfill \metainfo{Tipo}{especialista} \hfill \metainfo{Identificador}{web-developer} \hfill \metainfo{Ficha}{web-developer.md}}

\begin{descricao}
Desenvolve componentes de interface em SolidJS, observando práticas atuais de desenvolvimento web, experiência de uso e integração eficiente de dados.
\end{descricao}

\begin{funcao}
No Core, este perfil apoia tarefas relacionadas a desenvolve componentes de interface em SolidJS, observando práticas atuais de desenvolvimento web, experiência de uso e integração eficiente de dados.. O orquestrador pode encaminhar-lhe o trabalho na etapa correspondente; a resposta retorna ao fluxo principal para integração e conferência de fontes, critérios e permissões. A ficha documenta a função, mas não comprova que o perfil tenha sido chamado ou executado a tarefa.
\end{funcao}

\begin{arquitetura}
\textbf{Modo de execução declarado:} subagente · \textbf{Permissões:} edit: deny; bash: deny
\end{arquitetura}

\elemento{Pesquisador por busca na web (v1.0.0)}{\metainfo{Categoria}{Ferramentas e apoio transversal} \hfill \metainfo{Tipo}{especialista} \hfill \metainfo{Identificador}{web-search-researcher} \hfill \metainfo{Ficha}{web-search-researcher.md}}

\begin{descricao}
Pesquisa conteúdo na web a partir de uma URL e o analisa conforme a pergunta recebida.
\end{descricao}

\begin{funcao}
No Core, este perfil apoia tarefas relacionadas a pesquisa conteúdo na web a partir de uma URL e o analisa conforme a pergunta recebida.. O orquestrador pode encaminhar-lhe o trabalho na etapa correspondente; a resposta retorna ao fluxo principal para integração e conferência de fontes, critérios e permissões. A ficha documenta a função, mas não comprova que o perfil tenha sido chamado ou executado a tarefa.
\end{funcao}

\begin{arquitetura}
\textbf{Modo de execução declarado:} subagente · \textbf{Permissões:} read: allow; grep: allow; glob: allow; list: allow; bash: deny; edit: deny; todoread: deny; todowrite: deny; webfetch: deny
\end{arquitetura}

\elemento{Fluxo de produção acadêmica WS (v1.0.0)}{\metainfo{Categoria}{Ferramentas e apoio transversal} \hfill \metainfo{Tipo}{especialista} \hfill \metainfo{Identificador}{ws-academic-pipeline} \hfill \metainfo{Ficha}{ws-academic-pipeline.md}}

\begin{descricao}
Coordena um fluxo de produção acadêmica, encaminhando as etapas de pesquisa, redação e revisão conforme as instruções documentadas.
\end{descricao}

\begin{funcao}
No Core, este perfil apoia tarefas relacionadas a coordena um fluxo de produção acadêmica, encaminhando as etapas de pesquisa, redação e revisão conforme as instruções documentadas.. O orquestrador pode encaminhar-lhe o trabalho na etapa correspondente; a resposta retorna ao fluxo principal para integração e conferência de fontes, critérios e permissões. A ficha documenta a função, mas não comprova que o perfil tenha sido chamado ou executado a tarefa.
\end{funcao}

\begin{arquitetura}
\textbf{Modo de execução declarado:} subagente · \textbf{Permissões:} read: allow; edit: allow; glob: allow; grep: allow; bash: allow; pdf\_*: allow
\end{arquitetura}

\elemento{Agente de implementação de código (v1.0.0)}{\metainfo{Categoria}{Ferramentas e apoio transversal} \hfill \metainfo{Tipo}{especialista} \hfill \metainfo{Identificador}{ws-coder} \hfill \metainfo{Ficha}{ws-coder.md}}

\begin{descricao}
Apoia a implementação técnica e a modificação de código no escopo da tarefa.
\end{descricao}

\begin{funcao}
No Core, este perfil apoia tarefas relacionadas a apoia a implementação técnica e a modificação de código no escopo da tarefa.. O orquestrador pode encaminhar-lhe o trabalho na etapa correspondente; a resposta retorna ao fluxo principal para integração e conferência de fontes, critérios e permissões. A ficha documenta a função, mas não comprova que o perfil tenha sido chamado ou executado a tarefa.
\end{funcao}

\begin{arquitetura}
\textbf{Modo de execução declarado:} subagente · \textbf{Permissões:} edit: deny; bash: deny
\end{arquitetura}

\elemento{Arquiteto de pesquisa e documentação (v1.0.0)}{\metainfo{Categoria}{Ferramentas e apoio transversal} \hfill \metainfo{Tipo}{especialista} \hfill \metainfo{Identificador}{ws-researcher} \hfill \metainfo{Ficha}{ws-researcher.md}}

\begin{descricao}
Organiza pesquisa e documentação para estruturar conhecimento e fontes pertinentes.
\end{descricao}

\begin{funcao}
No Core, este perfil apoia tarefas relacionadas a organiza pesquisa e documentação para estruturar conhecimento e fontes pertinentes.. O orquestrador pode encaminhar-lhe o trabalho na etapa correspondente; a resposta retorna ao fluxo principal para integração e conferência de fontes, critérios e permissões. A ficha documenta a função, mas não comprova que o perfil tenha sido chamado ou executado a tarefa.
\end{funcao}

\begin{arquitetura}
\textbf{Modo de execução declarado:} subagente · \textbf{Permissões:} edit: deny; bash: deny
\end{arquitetura}

\elemento{Revisor técnico de código (v1.0.0)}{\metainfo{Categoria}{Ferramentas e apoio transversal} \hfill \metainfo{Tipo}{especialista} \hfill \metainfo{Identificador}{ws-reviewer} \hfill \metainfo{Ficha}{ws-reviewer.md}}

\begin{descricao}
Revisa código quanto a segurança, desempenho e qualidade técnica.
\end{descricao}

\begin{funcao}
No Core, este perfil apoia tarefas relacionadas a revisa código quanto a segurança, desempenho e qualidade técnica.. O orquestrador pode encaminhar-lhe o trabalho na etapa correspondente; a resposta retorna ao fluxo principal para integração e conferência de fontes, critérios e permissões. A ficha documenta a função, mas não comprova que o perfil tenha sido chamado ou executado a tarefa.
\end{funcao}

\begin{arquitetura}
\textbf{Modo de execução declarado:} subagente · \textbf{Permissões:} edit: deny; bash: deny
\end{arquitetura}

\elemento{Redator de documentação (v1.0.0)}{\metainfo{Categoria}{Ferramentas e apoio transversal} \hfill \metainfo{Tipo}{especialista} \hfill \metainfo{Identificador}{ws-scribe} \hfill \metainfo{Ficha}{ws-scribe.md}}

\begin{descricao}
Produz documentação e prosa voltadas à comunicação com pessoas.
\end{descricao}

\begin{funcao}
No Core, este perfil apoia tarefas relacionadas a produz documentação e prosa voltadas à comunicação com pessoas.. O orquestrador pode encaminhar-lhe o trabalho na etapa correspondente; a resposta retorna ao fluxo principal para integração e conferência de fontes, critérios e permissões. A ficha documenta a função, mas não comprova que o perfil tenha sido chamado ou executado a tarefa.
\end{funcao}

\begin{arquitetura}
\textbf{Modo de execução declarado:} subagente · \textbf{Permissões:} edit: deny; bash: deny
\end{arquitetura}

% fim do módulo 10 (gerado)
